% \iffalse meta-comment
%
% docxlike-format.dtx 是照 Word 的字体与段落对话框描述一个元素怎么排
%
% \fi
%
% \section{元素格式（format）}
%
% \changes{v3.2a}{2026/08/25}{加 \option{styles}：照 Word 的样式描述每种段落怎么排，内置
%   \texttt{Normal}、\texttt{Heading 1} 到 \texttt{Heading 4}、\texttt{Header}、\texttt{Footer}，
%   别的名字写了就新建、基于 \texttt{Normal}。键名照 Word 英文界面，按“字体”“段落”
%   “制表位”“边框”几个对话框分组嵌套，每键收 \texttt{auto}（交给默认表）；
%   \option{font} 与 \option{paragraph} 下各有一个 \option{snap-to-grid-when-document-grid-is-defined}，对应 Word 的两处
%   “对齐到网格”，\texttt{auto} 沿所在位置往上找（\cs{c\_\_docxlike\_format\_parent\_prop}），
%   根上没写就不贴；另有 \option{hyphenation} 与 \option{default-tab-stops}}
% \changes{v0.1.0}{2026/09/24}{字体键收字体名（\texttt{SimSun}、\texttt{宋体}、\texttt{Times New Roman}），
%   行高与上升部按名字照 Windows Word 的规则由字体度量现算，不再有写死的系数表，首行基线不另加
%   0.082\,bp；切换字体按名字声明字族，文档类可以用 \cs{docxlike\_font\_select\_new:nn} 指定替身，
%   同一副字体的中英文名、全名跟着一起登记}
% \changes{v0.1.0}{2026/09/24}{内部存储照 docx：条目名是 \texttt{w:rPr}、\texttt{w:pPr} 下的元素与
%   属性路径，值用缇、半磅、百分之一行等 docx 的单位；同一处量的几种写法归成一个槽，写一种
%   抹掉别的；公开的长度键不带单位时按 \texttt{bp}}
% \changes{v0.1.0}{2026/09/24}{Word 里没有的键挪出本包：\option{line-white}、\option{white-before}、
%   \option{white-after}、行距区间与 \texttt{length}、\texttt{last-line-centered}、\option{heading}、
%   \option{latin-bold}、\option{latin-line}、按版心底分档；换成给文档类的接入点（单位、
%   限定词、对齐取值、两个钩子、写存储）；\option{at} 只收一个值}
% \changes{v3.2a}{2026/08/16}{样式名可带文档类登记的限定词（\cs{docxlike\_format\_qualifiers:nn}），
%   带限定词的取值单独一份，读时先看当前限定词那一份，再看不带的}
% \changes{v3.2a}{2026/09/13}{\cs{docxlike\_format\_compute:n} 按 Word 的行盒模型换算：自然行高
%   与上升部按字体名照 Windows Word 的规则由字体度量算（中易 1.296875 与 1.0078125、Times 1.1499
%   与 0.9336）；倍数贴格时取整数格、固定值基线在盒顶下 0.8 倍；“一行”单位
%   认网格启不启用。四轮标定文档、三十组、一百六十
%   多条基线，来路在 \file{docxlike-page.dtx}}
% \changes{v3.2a}{2026/08/29}{\cs{docxlike\_format\_strut:} 把每一行的行盒撑成 Word 的尺寸，
%   贴格居中、不贴格顶对齐、固定值 0.8 倍、纯西文按 Times 三支（调用方点名的西文块也算，
%   双语题注的英文行因此按 Times 的上升部）；空的段落片段不成行
%   （公式独占一段时上下不再各多一行）；\hologo{LuaLaTeX} 下段末带封条}
% \changes{v3.2a}{2026/08/16}{正文接 \cs{normalsize}：字号、行距、首行基线、首行缩进四个
%   全局长度；\cs{rightskip} 收字符网格的余量；换页判据改成末行盒底：Lua 断行档
%   \cs{maxdepth} 归零、每行的深度都算进页（每行的盒就是 Word 的行盒），其余档
%   版心底留一个行深（\cs{@colroom} 补丁、\texttt{flush} 档页底留白刚性）；
%   孤行寡行四个罚分归零}
% \changes{v3.2a}{2026/08/29}{横轴两道胶按元素发：汉字缝伸量 0.02 行距、中西文胶 0.246 格
%   按字号缩放、可压 0.13 行距，空格伸量同档；同步 \pkg{xeCJK} 的
%   \cs{l\_\_xeCJK\_ecglue\_skip} 与 \pkg{luatexja} 的 kanjiskip，设完发钩子
%   \texttt{docxlike/hglue}；西文标点与乘除号并入半角标点类}
% \changes{v3.2a}{2026/09/20}{西文的半份网格增量（\option{LetterSpace}）只在 \hologo{XeLaTeX}
%   一路加，Lua 断行接管时由 Lua 在正文字号的段里自己补，交回 \TeX{} 的段 Word 本来不加}
% \changes{v3.2a}{2026/08/30}{\hologo{XeLaTeX} 一路的标点：全角标点补网格胶、相邻标点
%   缩半格、行中可压半格、行末减半走折处、盒内换等宽 kern、引号与波浪线自成字符类、
%   波浪号破折号不可断、西文接全角标点补一整格、半角闭号接汉字留零宽胶保住断点}
% \changes{v3.2a}{2026/08/31}{\pkg{xeCJK} 2025 与 2026 的中西文边界拉平：源码里的空格折成
%   中西文胶，\cs{verb}、\cs{num}、\cs{ref}、上标、链接的组边界补给老版，\cs{CJKecglue}
%   带闸不叠双份}
% \changes{v3.2a}{2026/09/20}{按目标排段的几个入口：\cs{docxlike\_format\_par:nn}（段前用
%   \cs{addvspace}，Word 的段前与前一段的段后取大不相加）、\cs{docxlike\_format\_par:nnn} 与
%   \cs{docxlike\_format\_with:nnn} 临时覆盖、\cs{docxlike\_format\_apply\_font:nN} 一族、
%   页眉页脚的 \cs{docxlike\_format\_apply\_hf:nN}、块与块之间照“上块行深加下块上升部”贴排的
%   \cs{docxlike\_format\_stack\_line:n}}
% \changes{v3.2a}{2026/08/24}{纯西文段由 \cs{docxlike\_format\_apply\_par:n} 铺版面，与别的段共用
%   \option{styles} 词汇；首行基线按 Times 的上升部}
% \changes{v3.2a}{2026/08/24}{新增 \cs{parmark}：模拟 Word 的空回车段，可选参数收行数或
%   \option{styles} 词汇；自然行高按段落标记的字体查；行距与贴格默认跟当前
%   段落走（Word 里回车出的空段继承上一段的属性，贴格的单倍空段占一整格，实测），
%   \option{snap-to-grid-when-document-grid-is-defined} 键可改；星号版的空盒之间补零罚分，页底装得下几行装几行；
%   浮动体的后距照留，空段另算}
% \changes{v3.2a}{2026/09/21}{\cs{verb} 照 Word 里能一字一字映射的写法排：Consolas 五号，
%   空格印成 U+02FD 加零罚分（\option{style}/\option{verb-space}\texttt{!=cell} 换成
%   不断行空格），长串在 \texttt{/ . -} 前可断（Lua 层），不贴字符网格；原先空格是
%   带伸缩的 \cs{nobreakspace}，不可断，两端对齐时被拉到 28.9\,bp}
%
% 规范里说的是 Word 的话：“章标题小二黑体居中，段前段后各一行，多倍行距 1.5”。
% 手算的话，这些话要折成 \texttt{28.34646bp} 这样的数散在好几处，改一处要连带重算三四个数。这里收 Word 对话框上的原话，换算交给代码：
% \begin{latex}
% \docxlikesetup
%   {
%     styles =
%       {
%         Heading 1 = { font = { asian-text-font = SimHei , size = xiaoer , font-style = bold } ,
%                       paragraph = { alignment = centered ,
%                                     spacing = { before = 1 line , after = 1 line ,
%                                                 line-spacing = multiple , at = 1.5 } } } ,
%       }
%   }
% \end{latex}
%
% \subsection{键}
%
% 公开的键照 Word 英文界面的原词，按“修改样式 → 格式”菜单分成 \option{font}、
% \option{paragraph}、\option{tabs}、\option{border} 四组，组内再照对话框的分组嵌套；
% 全表在手册的“样式”一节。内部存储照 docx：条目名是 \texttt{w:rPr}、\texttt{w:pPr} 底下的
% 元素与属性路径，值用 docx 的单位（缇是二十分之一磅，字号是半磅，按行、按字的量是百分之一，
% 倍数行距是 240 分之一），从 docx 转过来时原样填进去。一个公开的键可能落进几条：
% \begin{longtblr}[caption={公开的键与内部存储}, label={tab:impl-format-keys}]
%   {colspec={l l}}
% \toprule
% 公开的键 & 内部条目 \\ \midrule
% \option{font}/\option{asian-text-font}、\option{latin-text-font} & \texttt{rPr/rFonts/eastAsia}、\texttt{rPr/rFonts/ascii} \\
% \option{font}/\option{font-style}、\option{all-caps} & \texttt{rPr/b}、\texttt{rPr/caps} \\
% \option{font}/\option{size} & \texttt{rPr/sz}（半磅） \\
% \option{font}/\option{character-spacing} & \texttt{rPr/spacing}（缇，带符号） \\
% \option{font}/\option{snap-to-grid-when-document-grid-is-defined} & \texttt{rPr/snapToGrid} \\
% \option{paragraph}/\option{alignment} & \texttt{pPr/jc}：\texttt{left}、\texttt{center}、\texttt{right}、\texttt{both}、\texttt{distribute} \\
% \option{paragraph}/\option{indentation} & \texttt{pPr/ind/} 下的 \texttt{left}、\texttt{leftChars}、\texttt{right}、\texttt{rightChars}、\texttt{firstLine}、\texttt{firstLineChars}、\texttt{hanging}、\texttt{hangingChars} \\
% \option{paragraph}/\option{spacing} & \texttt{pPr/spacing/} 下的 \texttt{before}、\texttt{beforeLines}、\texttt{after}、\texttt{afterLines}、\texttt{line}、\texttt{lineRule} \\
% \option{paragraph}/\option{snap-to-grid-when-document-grid-is-defined}、\option{pagination} 的 \option{keep-with-next}、\option{keep-lines-together}、\option{page-break-before} & \texttt{pPr/snapToGrid}、\texttt{pPr/keepNext}、\texttt{pPr/keepLines}、\texttt{pPr/pageBreakBefore} \\
% \option{tabs}、\option{border} & \texttt{pPr/tabs}（每个停靠点“对齐 位置 前导符”，位置是缇）、\texttt{pPr/pBdr/top}、\texttt{pPr/pBdr/bottom} \\
% \option{numbering}\texttt{ = none}（“编号”按钮里的“无”） & \texttt{pPr/numPr/numId} 为 \texttt{0}：链接了列表的样式这一段不编号 \\
% \bottomrule
% \end{longtblr}
% 同一处量在 docx 里有几种写法，比如段前按磅是 \texttt{before}、按行是 \texttt{beforeLines}，
% 写一种就把同一处的别的写法抹掉，这样的一组叫一个槽（\texttt{before}、\texttt{after}、\texttt{line}、
% \texttt{ind-left}、\texttt{ind-right}、\texttt{ind-special}）。首行缩进写 \texttt{none} 存成
% \texttt{firstLine = 0}，好挡住上级样式的缩进。
%
% Word 里没有的写法（区间、文档类登记的单位）原样记在 \texttt{ext/}\meta{槽} 那几条：
% \meta{槽}、\meta{槽}\texttt{-unit}，区间的第二个端点在 \meta{槽}\texttt{-2}、\meta{槽}\texttt{-2-unit}。
% 公开的键只写 Word 的那几种；别的由文档类经“给文档类的接入点”一节的接口加。
%
% 从长度折成缇、半磅时保留三位小数：\hologo{TeX} 读 \texttt{bp} 时截掉的不到 1\,sp 会让
% \texttt{20bp} 变成 399.9999 缇，三位小数刚好收回整数；折回长度时写成 \texttt{bp} 交给
% \hologo{TeX} 自己读，与直接写 \texttt{bp} 一个 sp 都不差。Word 本身只存整数，从 docx 来的
% 值就是整数；文档类为了仿真写进来的不是整数也照样存，不强行取整。
%
% 行距的“设置值”在内部与模式合成一个词：倍数照写，固定值加 \texttt{!}，最小值加 \texttt{+}。
% 两个 \option{snap-to-grid-when-document-grid-is-defined} 分属字体与段落，是 Word 本来的样子：段落对话框里那一个
% （\texttt{w:pPr/w:snapToGrid}）管行网格，字体高级页里那一个（\texttt{w:rPr/w:snapToGrid}）
% 管字符网格。
%
% \texttt{1 line}、\texttt{2 ch} 这几种写法里的空格是有意义的，解析器按空格切词。
% 写在类源码里要把空格改成 \pkg{expl3} 的空格记号，那里在 \cs{ExplSyntaxOn} 之内，
% 字面空格会被吃掉。
%
% \subsection{样式}
%
% 样式名就是 Word 里的名字，去掉空格作内部标识（docx 的 \texttt{w:styleId}）。本包登记
% \texttt{Normal}（正文）、\texttt{Heading1} 到 \texttt{Heading4}、\texttt{Header}、\texttt{Footer}，
% 文档类登记自己要的，写到没登记的名字就新建一个。
%
% \texttt{Normal} 管的东西比别的样式多一层：除了字号行距，\cs{parindent} 与 \cs{topskip}
% 也从它算，见下面“接到正文上”一节。
%
% ^^A ======================================================================
% ^^A Module docxlike-format: 元素格式（各个类共用）
% ^^A ======================================================================
%    \begin{macrocode}
%<*docxlike-format>
%    \end{macrocode}
%
% \subsection{给文档类的钩子与登记}
%
% 两个 \pkg{lthooks} 钩子，文档类往里挂代码：
% \begin{itemize}
% \item \texttt{docxlike/latin-font-reload}：西文字母加半格余量（\cs{g\_docxlike\_latin\_pitch\_bool}）
%   要重新声明西文字族，\pkg{docxlike} 发完 \cs{defaultfontfeatures} 就用这个钩子让文档类
%   把自己的西文字体重设一遍。钩子是空的就不做这件事。
% \item \texttt{docxlike/display}：装行间公式那一套钩子时顺带执行，文档类在这里接上自己
%   对行间公式的处理。
% \end{itemize}
%    \begin{macrocode}
\hook_new:n { docxlike / latin-font-reload }
\hook_new:n { docxlike / display }
\hook_new:n { docxlike / hglue }
%    \end{macrocode}
%
% 字号名。\option{size} 与 \cs{parmark} 的字号可以写中文字号的名字，
% 文档类用 \cs{docxlike\_size\_new:nn}\marg{名字}\marg{长度} 加自己的。
%    \begin{macrocode}
\prop_new:N \g__docxlike_size_prop
\tl_new:N \l__docxlike_size_tl
\cs_new_protected:Npn \docxlike_size_new:nn #1#2
  { \prop_gput:Nnx \g__docxlike_size_prop {#1} { \dim_eval:n {#2} } }
\cs_generate_variant:Nn \prop_gput:Nnn { Nnx }
\docxlike_size_new:nn { chuhao }  { 42bp }
\docxlike_size_new:nn { xiaochu } { 36bp }
\docxlike_size_new:nn { yihao }   { 26bp }
\docxlike_size_new:nn { xiaoyi }  { 24bp }
\docxlike_size_new:nn { erhao }   { 22bp }
\docxlike_size_new:nn { xiaoer }  { 18bp }
\docxlike_size_new:nn { sanhao }  { 16bp }
\docxlike_size_new:nn { xiaosan } { 15bp }
\docxlike_size_new:nn { sihao }   { 14bp }
\docxlike_size_new:nn { xiaosi }  { 12bp }
\docxlike_size_new:nn { wuhao }   { 10.5bp }
\docxlike_size_new:nn { xiaowu }  { 9bp }
\docxlike_size_new:nn { liuhao }  { 7.5bp }
\docxlike_size_new:nn { xiaoliu } { 6.5bp }
\docxlike_size_new:nn { qihao }   { 5.5bp }
\docxlike_size_new:nn { bahao }   { 5bp }
%    \end{macrocode}
%
% \cs{\_\_docxlike\_format\_parend\_guard:} 排在 \texttt{para/end} 钩子里、段末撑杆
% \cs{docxlike\_format\_strut:} 之前，把撑杆焊在段落最后一个字上。\hologo{XeLaTeX} 下
% 是空操作——那边本来就是焊死的。
%
% \textbf{病灶。} 段末的水平列表是「\dots 。\meta{撑杆}\cs{penalty}~10000
% \cs{parfillskip}」。\pkg{xeCJK} 在句号与撑杆之间什么都不放，撑杆离不开末行。
% \pkg{luatexja} 会在和文字符与后续节点之间插一段 JFM 边界胶——句号（JFM 里是
% 半宽字形）后面是“\cs{glue}~6.02\,pt~minus~6.02\,pt”，这段胶是合法断点。
% 于是当段落末行恰好排满时，TeX 有两个选择：把“边界胶＋撑杆＋\cs{parfillskip}”
% 压上末行，或在边界胶处断行——断出来的新行只有撑杆和 \cs{parfillskip}，
% badness 为零，总是胜出。结果是段与段之间凭空多出一行空白，高度正好一个
% \cs{baselineskip}（撑杆的高深）。末行差一点排满时是另一副面孔：那段
% minus 胶被压光，句号收成半宽、整行两端撑满，而 \hologo{XeLaTeX} 是自然宽度加
% \cs{parfillskip} 补白。
%
% \textbf{封条是 \cs{ltjfakeboxbdd}。} 它告诉 \pkg{luatexja}「此处按盒子边界
% 处理」：盒末的尾随边界胶本来就不插（\cs{hbox} 实测如此），于是句号与撑杆之间
% 不再有断点，两种病症一起消失。段落以普通汉字收尾时 JFM 本来就不在字与撑杆之间
% 插胶，没病也不受影响；只有以标点收尾的段落走到这条上。
% 没装 \pkg{luatexja} 时什么也不做。
%    \begin{macrocode}
\sys_if_engine_luatex:TF
  {
    \cs_new_protected:Npn \__docxlike_format_parend_guard:
      { \cs_if_exist_use:N \ltjfakeboxbdd }
  }
  { \cs_new_protected:Npn \__docxlike_format_parend_guard: { } }
%    \end{macrocode}
%
% 目标清单与每个目标的存储。取值存进 \cs{g\_\_docxlike\_format\_}\meta{目标}\cs{\_prop}，
% 带单位的键存两条：值一条、单位一条。
%    \begin{macrocode}
\clist_new:N \g_docxlike_format_target_clist
\prop_new:N \g__docxlike_format_parent_prop
\prop_new:N \g__docxlike_format_name_prop
\cs_new_protected:Npn \docxlike_format_target_new:nn #1#2
  {
    \prop_if_exist:cF { g__docxlike_format_#1_prop }
      {
        \clist_gput_right:Nn \g_docxlike_format_target_clist {#1}
        \__docxlike_squash:nN {#1} \l__docxlike_value_try_str
        \prop_if_in:NVF \g__docxlike_format_name_prop \l__docxlike_value_try_str
          { \prop_gput:NVn \g__docxlike_format_name_prop \l__docxlike_value_try_str {#1} }
        \prop_new:c { g__docxlike_format_#1_prop }
        \prop_new:c { g__docxlike_format_#1_default_prop }
        \clist_map_inline:Nn \g__docxlike_format_layer_clist
          { \__docxlike_format_target_layer_new:nn {#1} {##1} }
      }
    \tl_if_blank:nF {#2}
      { \prop_gput:Nnn \g__docxlike_format_parent_prop {#1} {#2} }
  }
\cs_generate_variant:Nn \docxlike_format_target_new:nn { Vn }
%    \end{macrocode}
%
%    \begin{macrocode}
\str_new:N \l__docxlike_format_target_str
\seq_new:N \l__docxlike_format_word_seq
\int_new:N \l__docxlike_format_count_int
\tl_new:N \l__docxlike_format_head_tl
\tl_new:N \l__docxlike_format_tail_tl
\tl_new:N \l__docxlike_format_rest_tl
%    \end{macrocode}
%
% \cs{\_\_docxlike\_format\_put:nn}\marg{键}\marg{值} 写进当前目标的存储。
%
% 值当场展开再存。解析器交过来的是 \cs{dim\_eval:n} 与 \cs{fp\_to\_decimal:n} 这种
% 还没算的式子，里面读着 \cs{l\_\_docxlike\_format\_head\_tl} 这些临时变量；原样存下去，
% 等到取用时才算，那几个临时变量早被下一条键覆盖了。
% 存一项。取值写 \texttt{auto} 时不存，反而把已有的条目删掉——读取那一头
% \cs{\_\_docxlike\_format\_read:nnnN} 每一项都带着内建默认，条目不在表里就用那个默认，
% 而那个默认正是包自己的规矩。所以\texttt{auto} 与“没写这一项”等价，
% 这也正是它该有的含义。
%
% 举一个：\option{snap-to-grid-when-document-grid-is-defined} 的内建默认是 \texttt{true} 且要求文档网格已经
% 立起来，所以写 \texttt{auto} 就是“有网格就贴，没有就不贴”，跟着当前选项走。
%
% 数值类的键照样收 \texttt{auto}，含义一样，只是那些键本来就不太用得上。
%
% \emph{限定词。} 同一个样式在某种条件下可以要一套完全不同的取值（比如某个文档类按版心底
% 两种排法各要一套）。文档类用 \cs{docxlike\_format\_qualifiers:nn}\marg{限定词表}\marg{当前}
% 登记，\meta{当前}是可展开的代码，给出此刻起作用的限定词，没有时给空。之后样式名写成
% \meta{样式}\texttt{/}\meta{限定词}，那几条键存进这个样式的另一份存储。
%
% 所以每个样式有几份存储：不带限定词的一份，每个限定词各一份。读的时候先看当前限定词
% 那一份，没有再看不带的，还没有就用内建默认。
%
% 带限定词的取值压过不带的，所以用户在外头写一条普通的键时，还要把各限定词里同名的
% 条目抹掉——不然默认表在那里写的值会盖住用户刚写的。默认表自己写的时候不抹，
% 它就是在往各份里填。
%
% 文档类写自己那份默认表时把 \cs{g\_docxlike\_format\_defaults\_bool} 立起来，写完放下：
% 这期间写进去的值同时记成“默认值”，\texttt{auto} 回退到的就是它们。
%    \begin{macrocode}
\cs_generate_variant:Nn \prop_gput:Nnn { cnx }
\bool_new:N \g_docxlike_format_defaults_bool
\tl_new:N \l__docxlike_format_layer_tl
\clist_new:N \g__docxlike_format_layer_clist
\cs_new:Npn \__docxlike_format_name:nnn #1#2#3
  { g__docxlike_format_ #1 \exp_args:Ne \tl_if_empty:nF {#2} { _ #2 } #3 _prop }
\cs_new_protected:Npn \__docxlike_format_new:nnn #1#2#3
  {
    \prop_if_exist:cF { \__docxlike_format_name:nnn {#1} {#2} {#3} }
      { \prop_new:c { \__docxlike_format_name:nnn {#1} {#2} {#3} } }
  }
\cs_new:Npn \__docxlike_format_layer_now: { }
\cs_new_protected:Npn \docxlike_format_qualifiers:nn #1#2
  {
    \clist_gset:Nn \g__docxlike_format_layer_clist {#1}
    \cs_gset:Npn \__docxlike_format_layer_now: {#2}
    \clist_map_inline:nn {#1}
      {
        \prop_if_exist:cF { l__docxlike_format_par_save_##1_prop }
          { \prop_new:c { l__docxlike_format_par_save_##1_prop } }
        \clist_map_inline:Nn \g_docxlike_format_target_clist
          { \__docxlike_format_target_layer_new:nn {####1} {##1} }
      }
  }
\cs_new_protected:Npn \__docxlike_format_target_layer_new:nn #1#2
  {
    \__docxlike_format_new:nnn {#1} {#2} { }
    \__docxlike_format_new:nnn {#1} {#2} { _default }
  }
\prg_new_protected_conditional:Npnn \__docxlike_format_get:nnN #1#2#3 { T , F , TF }
  {
    \prop_get:cnNTF
      { \__docxlike_format_name:nnn {#1} { \__docxlike_format_layer_now: } { } } {#2} #3
      { \prg_return_true: }
      {
        \prop_get:cnNTF { \__docxlike_format_name:nnn {#1} { } { } } {#2} #3
          { \prg_return_true: } { \prg_return_false: }
      }
  }
\prg_generate_conditional_variant:Nnn \__docxlike_format_get:nnN { VnN }
  { T , F , TF }
\prg_new_conditional:Npnn \__docxlike_format_if_blank:n #1 { T , F , TF }
  {
    \bool_lazy_and:nnTF
      {
        \prop_if_empty_p:c
          { \__docxlike_format_name:nnn {#1} { \__docxlike_format_layer_now: } { } }
      }
      { \prop_if_empty_p:c { \__docxlike_format_name:nnn {#1} { } { } } }
      { \prg_return_true: } { \prg_return_false: }
  }
\prg_new_eq_conditional:NNn \docxlike_format_if_blank:n \__docxlike_format_if_blank:n
  { T , F , TF }
\cs_new_protected:Npn \__docxlike_format_put:nn #1#2
  {
    \str_if_eq:nnTF {#2} { auto }
      { \__docxlike_format_restore:n {#1} }
      { \__docxlike_format_put_value:nn {#1} {#2} }
  }
\cs_new_protected:Npn \__docxlike_format_put_value:nn #1#2
  {
    \__docxlike_format_store:VVnn \l__docxlike_format_target_str
      \l__docxlike_format_layer_tl {#1} {#2}
    \bool_if:NTF \g_docxlike_format_defaults_bool
      {
        \__docxlike_format_store:VVnnn \l__docxlike_format_target_str
          \l__docxlike_format_layer_tl { _default } {#1} {#2}
      }
      { \__docxlike_format_shadow:n {#1} }
  }
\cs_new_protected:Npn \__docxlike_format_store:nnnnn #1#2#3#4#5
  {
    \__docxlike_format_new:nnn {#1} {#2} {#3}
    \prop_gput:cnx { \__docxlike_format_name:nnn {#1} {#2} {#3} } {#4} {#5}
  }
\cs_new_protected:Npn \__docxlike_format_store:nnnn #1#2#3#4
  { \__docxlike_format_store:nnnnn {#1} {#2} { } {#3} {#4} }
\cs_generate_variant:Nn \__docxlike_format_store:nnnnn { VVnnn }
\cs_generate_variant:Nn \__docxlike_format_store:nnnn { VVnn }
\cs_new_protected:Npn \__docxlike_format_shadow:n #1
  {
    \tl_if_empty:NT \l__docxlike_format_layer_tl
      {
        \clist_map_inline:Nn \g__docxlike_format_layer_clist
          {
            \prop_gremove:cn
              { \__docxlike_format_name:nnn { \l__docxlike_format_target_str } {##1} { } }
              {#1}
          }
      }
  }
%    \end{macrocode}
%
% 删一项。区间的第二个端点只在写了两个词时才有，写一个词就得把上一次留下的抹掉。
% 存一个空值也能起同样的作用，但表里会多出永远读不到东西的条目，
% \file{testfiles/56-format-keys.lvt} 把整张表印出来对，一眼就看得见。
%    \begin{macrocode}
\cs_new_protected:Npn \__docxlike_format_drop:n #1
  {
    \__docxlike_format_erase:VVn \l__docxlike_format_target_str
      \l__docxlike_format_layer_tl {#1}
    \bool_if:NTF \g_docxlike_format_defaults_bool
      {
        \__docxlike_format_erase:VVnn \l__docxlike_format_target_str
          \l__docxlike_format_layer_tl { _default } {#1}
      }
      { \__docxlike_format_shadow:n {#1} }
  }
\cs_new_protected:Npn \__docxlike_format_erase:nnnn #1#2#3#4
  {
    \__docxlike_format_new:nnn {#1} {#2} {#3}
    \prop_gremove:cn { \__docxlike_format_name:nnn {#1} {#2} {#3} } {#4}
  }
\cs_new_protected:Npn \__docxlike_format_erase:nnn #1#2#3
  { \__docxlike_format_erase:nnnn {#1} {#2} { } {#3} }
\cs_generate_variant:Nn \__docxlike_format_erase:nnnn { VVnn }
\cs_generate_variant:Nn \__docxlike_format_erase:nnn { VVn }
%    \end{macrocode}
%
% \texttt{auto} 退回默认表里这一项的取值。不带限定词地写 \texttt{auto} 时，各限定词
% 里的那一项也一起退回去，不然用户刚才写的那条还留在外面压着。
%    \begin{macrocode}
\cs_new_protected:Npn \__docxlike_format_restore:n #1
  {
    \__docxlike_format_restore:VVn \l__docxlike_format_target_str
      \l__docxlike_format_layer_tl {#1}
    \tl_if_empty:NT \l__docxlike_format_layer_tl
      {
        \clist_map_inline:Nn \g__docxlike_format_layer_clist
          {
            \__docxlike_format_restore:Vnn \l__docxlike_format_target_str {##1} {#1}
          }
      }
  }
\cs_new_protected:Npn \__docxlike_format_restore:nnn #1#2#3
  {
    \__docxlike_format_new:nnn {#1} {#2} { _default }
    \__docxlike_format_new:nnn {#1} {#2} { }
    \prop_get:cnNTF { \__docxlike_format_name:nnn {#1} {#2} { _default } } {#3}
      \l__docxlike_format_value_tl
      {
        \prop_gput:cnV { \__docxlike_format_name:nnn {#1} {#2} { } } {#3}
          \l__docxlike_format_value_tl
      }
      { \prop_gremove:cn { \__docxlike_format_name:nnn {#1} {#2} { } } {#3} }
  }
\cs_generate_variant:Nn \__docxlike_format_restore:nnn { VVn , Vnn }
%    \end{macrocode}
%
% 把取值按空格切开，留下四样东西：词数、头一个词、末一个词、头一个词之后的全部。
% 带单位的写法有两种词序，\texttt{2 ch} 是数在前，\texttt{exactly 25bp} 是词在前，
% 各个解析器按自己的词序挑用。\texttt{hanging 4 ch} 是三个词，模式在头上，
% 数量在“其余”里，那一段再交给收数量的解析器切一次。
%    \begin{macrocode}
\cs_new_protected:Npn \__docxlike_format_split:n #1
  {
    \seq_set_split:Nnn \l__docxlike_format_word_seq { ~ } { #1 }
    \seq_remove_all:Nn \l__docxlike_format_word_seq { }
    \int_set:Nn \l__docxlike_format_count_int
      { \seq_count:N \l__docxlike_format_word_seq }
    \seq_get_left:NN  \l__docxlike_format_word_seq \l__docxlike_format_head_tl
    \seq_get_right:NN \l__docxlike_format_word_seq \l__docxlike_format_tail_tl
    \seq_pop_left:NN  \l__docxlike_format_word_seq \l__docxlike_format_rest_tl
    \tl_set:Nx \l__docxlike_format_rest_tl
      { \seq_use:Nn \l__docxlike_format_word_seq { ~ } }
  }
\prg_new_conditional:Npnn \__docxlike_format_one_word: { TF }
  {
    \int_compare:nNnTF { \l__docxlike_format_count_int } = { 1 }
      { \prg_return_true: } { \prg_return_false: }
  }
\cs_new_eq:NN \docxlike_format_put:nn \__docxlike_format_put:nn
\cs_new:Npn \__docxlike_format_twips:n #1
  { \fp_eval:n { round ( \dim_to_fp:n { \__docxlike_bp:n {#1} } * 1440 / 72.27 , 3 ) } }
\cs_new:Npn \__docxlike_format_half_points:n #1
  { \fp_eval:n { round ( \dim_to_fp:n { \__docxlike_bp:n {#1} } * 144 / 72.27 , 3 ) } }
\cs_new:Npn \__docxlike_format_from_twips:n #1
  { \fp_eval:n { ( #1 ) / 20 } bp }
\cs_new:Npn \__docxlike_format_from_half_points:n #1
  { \fp_eval:n { ( #1 ) / 2 } bp }
\prop_const_from_keyval:Nn \c__docxlike_format_slot_prop
  {
    before =
      {
        ext/before , ext/before-unit , ext/before-2 , ext/before-2-unit ,
        pPr/spacing/beforeLines , pPr/spacing/before
      } ,
    after =
      {
        ext/after , ext/after-unit , ext/after-2 , ext/after-2-unit ,
        pPr/spacing/afterLines , pPr/spacing/after
      } ,
    line =
      {
        ext/line , ext/line-unit , ext/line-2 , ext/line-2-unit ,
        pPr/spacing/line , pPr/spacing/lineRule
      } ,
    ind-left = { pPr/ind/leftChars , pPr/ind/left } ,
    ind-right = { pPr/ind/rightChars , pPr/ind/right } ,
    ind-special =
      {
        pPr/ind/firstLineChars , pPr/ind/firstLine ,
        pPr/ind/hangingChars , pPr/ind/hanging
      } ,
  }
\clist_new:N \l__docxlike_format_slot_clist
\cs_new_protected:Npn \__docxlike_format_slot_keys:n #1
  {
    \clist_set:Ne \l__docxlike_format_slot_clist
      { \prop_item:Nn \c__docxlike_format_slot_prop {#1} }
  }
\cs_new_protected:Npn \__docxlike_format_slot_drop:n #1
  {
    \__docxlike_format_slot_keys:n {#1}
    \clist_map_function:NN \l__docxlike_format_slot_clist \__docxlike_format_drop:n
  }
\cs_new_protected:Npn \__docxlike_format_slot_restore:n #1
  {
    \__docxlike_format_slot_keys:n {#1}
    \clist_map_function:NN \l__docxlike_format_slot_clist \__docxlike_format_restore:n
  }
%    \end{macrocode}
%
% 读一个槽：先看当前限定词那一份，那一份里这个槽有任何一条就只用那一份，没有再看
% 不带限定词的。找到的几条抄进 \cs{l\_\_docxlike\_format\_slot\_prop}。
%    \begin{macrocode}
\prop_new:N \l__docxlike_format_slot_prop
\tl_new:N \l__docxlike_format_slot_now_tl
\bool_new:N \l__docxlike_format_slot_found_bool
\cs_new_protected:Npn \__docxlike_format_slot_find:nn #1#2
  {
    \prop_clear:N \l__docxlike_format_slot_prop
    \__docxlike_format_slot_keys:n {#2}
    \bool_set_false:N \l__docxlike_format_slot_found_bool
    \tl_set:Ne \l__docxlike_format_slot_now_tl { \__docxlike_format_layer_now: }
    \tl_if_empty:NF \l__docxlike_format_slot_now_tl
      {
        \__docxlike_format_slot_from:e
          { \__docxlike_format_name:nnn {#1} { \l__docxlike_format_slot_now_tl } { } }
      }
    \bool_if:NF \l__docxlike_format_slot_found_bool
      { \__docxlike_format_slot_from:e { \__docxlike_format_name:nnn {#1} { } { } } }
  }
\cs_new_protected:Npn \__docxlike_format_slot_from:n #1
  {
    \prop_if_exist:cT {#1}
      {
        \clist_map_inline:Nn \l__docxlike_format_slot_clist
          {
            \prop_get:cnNT {#1} {##1} \l__docxlike_format_value_tl
              {
                \bool_set_true:N \l__docxlike_format_slot_found_bool
                \prop_put:NnV \l__docxlike_format_slot_prop {##1} \l__docxlike_format_value_tl
              }
          }
      }
  }
\cs_generate_variant:Nn \__docxlike_format_slot_from:n { e }
\prg_new_protected_conditional:Npnn \__docxlike_format_slot_if_set:nn #1#2 { T , F , TF }
  {
    \__docxlike_format_slot_find:nn {#1} {#2}
    \bool_if:NTF \l__docxlike_format_slot_found_bool
      { \prg_return_true: } { \prg_return_false: }
  }
\tl_new:N \l_docxlike_format_style_tl
\tl_new:N \l_docxlike_format_side_tl
\cs_new_protected:Npn \docxlike_format_line_unit_new:nn #1#2
  { \cs_gset_protected:cpn { __docxlike_format_line_unit_ #1 :nN } ##1##2 {#2} }
\cs_new_protected:Npn \docxlike_format_space_unit_new:nn #1#2
  { \cs_gset_protected:cpn { __docxlike_format_space_unit_ #1 :nN } ##1##2 {#2} }
\cs_new_protected:Npn \__docxlike_format_enum:nnn #1#2#3
  {
    \clist_if_in:nnTF {#2} {#3}
      { \__docxlike_format_put:nn {#1} {#3} }
      { \__docxlike_bad_value:Vnn \l_keys_key_str {#3} {#2} }
  }
\cs_new_protected:Npn \__docxlike_format_bool:nn #1#2
  { \__docxlike_format_enum:nnn {#1} { true , false , auto } {#2} }
%    \end{macrocode}
%
% \texttt{spacing} 是 Word“字体”对话框“高级”页里的字符间距，
% 对应 \texttt{w:rPr/w:spacing}。Word 那边分成“标准／加宽／紧缩”加一个磅值，
% 这里合成一个带符号的长度：正数加宽、负数紧缩、零是标准。
% \texttt{+0.2bp} 与 \texttt{0.2bp} 一样，前缀正号写不写都行。
%
% 标准没有专门的词，写任何等于零的长度就是。它与 \texttt{auto} 不同：
% \texttt{auto} 是“这一项不设，跟着外层的字距走”，\texttt{0bp} 是
% “明确设成不加不减”，会把外层的字距挡住，见 \cs{\_\_docxlike\_format\_hglue\_cc:}。
%
% \emph{与 \option{line-spacing} 那个加号不是一回事。} \texttt{20bp+} 的加号
% 在\emph{后面}，意思是“至少”；这里的加号在\emph{前面}，是数的正负号。
%    \begin{macrocode}
\cs_new_protected:Npn \__docxlike_format_char_spacing:n #1
  {
    \str_if_eq:nnTF {#1} { auto }
      { \__docxlike_format_put:nn { rPr/spacing } { auto } }
      { \__docxlike_format_put:ne { rPr/spacing } { \__docxlike_format_twips:n {#1} } }
  }
\cs_generate_variant:Nn \__docxlike_format_char_spacing:n { e }
\cs_generate_variant:Nn \__docxlike_format_half_points:n { V }
\cs_generate_variant:Nn \__docxlike_format_put:nn { ne }
%    \end{macrocode}
%
% \option{size} 收字号名或长度。名字查 \cs{docxlike\_size\_new:nn} 登记的那张表，
% 查不到就当长度读。
%    \begin{macrocode}
\cs_new_protected:Npn \__docxlike_format_size:n #1
  {
    \str_if_eq:nnTF {#1} { auto }
      { \__docxlike_format_put:nn { rPr/sz } { auto } }
      {
        \prop_get:NnNTF \g__docxlike_size_prop {#1} \l__docxlike_size_tl
          {
            \__docxlike_format_put:ne { rPr/sz }
              { \__docxlike_format_half_points:V \l__docxlike_size_tl }
          }
          { \__docxlike_format_put:ne { rPr/sz } { \__docxlike_format_half_points:n {#1} } }
      }
  }
%    \end{macrocode}
%
% \emph{\cs{\_\_docxlike\_format\_measure:nnnn}。}
% 行距的取值，度量是 \texttt{spacing} 时按词认。写一个词是刚性；写两个词是区间，
% 两个词不分先后，谁大谁小由代码判。公开的 \option{line-spacing} 只写一个词，区间是
% 文档类经 \cs{docxlike\_format\_measure:nnnn} 写的。
%
% 一个词有五种写法：
% \begin{longtblr}[caption={行距一个词的五种写法}, label={tab:impl-line-spacing}]
%   {colspec={l l l}}
% \toprule
% 写法 & 存下来的单位 & 意思 \\ \midrule
% \texttt{single}、\texttt{double} & \texttt{auto} & Word 下拉框里的词，折成 1 与 2 \\
% \texttt{1.25} & \texttt{auto} & 倍数 \\
% \texttt{25bp!} & \texttt{exact} & Word 的固定值 \\
% \texttt{10bp+} & \texttt{atLeast} & Word 的最小值 \\
% \texttt{19bp} & \texttt{length} & 就是这么长，不设下限；Word 里没有，给文档类用 \\
% \bottomrule
% \end{longtblr}
% 末一个字符定归属：\texttt{!} 与 \texttt{+} 是模式记号，数字与小数点收尾是倍数，
% 字母收尾是长度（\hologo{TeX} 的单位一律是字母）。所以
% \texttt{3mm~2} 这种一头长度一头倍数的区间也收。
%
% \texttt{length} 与 \texttt{atLeast} 差一条下限，\texttt{length} 与
% \texttt{exact} 差一处基线落点，三者不能合并。存下来的单位用 docx 里 \texttt{w:lineRule}
% 的取值，倍数是 \texttt{auto}。写进存储的值是 \texttt{auto} 时 \cs{\_\_docxlike\_format\_put:nn}
% 当作“退回默认”，所以单位走 \cs{\_\_docxlike\_format\_put\_value:nn}，照原样存。
%
% 度量不是 \texttt{spacing} 时不认词，每个端点的单位直接记成度量的名字，换算交给文档类
% 登记的单位（\cs{docxlike\_format\_line\_unit\_new:nn}）。几个键写同一个槽时，同一个
% 样式里都写了就报错，按后写的算。
%    \begin{macrocode}
\seq_new:N \l__docxlike_format_end_seq
\tl_new:N \l__docxlike_format_last_tl
\prop_new:N \l__docxlike_format_claim_prop
\prg_generate_conditional_variant:Nnn \str_if_in:nn { nV } { TF }
\cs_new_protected:Npn \__docxlike_format_claim:nn #1#2
  {
    \prop_get:NnNT \l__docxlike_format_claim_prop {#1} \l__docxlike_format_last_tl
      {
        \str_if_eq:VnF \l__docxlike_format_last_tl {#2}
          {
            \msg_error:nnVn { docxlike } { measure-clash }
              \l__docxlike_format_last_tl {#2}
          }
      }
    \prop_put:Nnn \l__docxlike_format_claim_prop {#1} {#2}
  }
\cs_new_protected:Npn \__docxlike_format_measure:nnnn #1#2#3#4
  {
    \__docxlike_format_claim:nn {#2} {#1}
    \str_if_eq:nnTF {#4} { auto }
      { \__docxlike_format_slot_restore:n {#2} }
      {
        \seq_set_split:Nnn \l__docxlike_format_end_seq { ~ } {#4}
        \seq_remove_all:Nn \l__docxlike_format_end_seq { }
        \tl_clear:N \l__docxlike_format_end_ii_unit_tl
        \tl_clear:N \l__docxlike_format_end_ii_value_tl
        \int_case:nnTF { \seq_count:N \l__docxlike_format_end_seq }
          {
            { 1 }
              {
                \__docxlike_format_end:nnne {#2} { } {#3}
                  { \seq_item:Nn \l__docxlike_format_end_seq { 1 } }
              }
            { 2 }
              {
                \__docxlike_format_end:nnne {#2} { } {#3}
                  { \seq_item:Nn \l__docxlike_format_end_seq { 1 } }
                \__docxlike_format_end:nnne {#2} { -2 } {#3}
                  { \seq_item:Nn \l__docxlike_format_end_seq { 2 } }
              }
          }
          {
            \__docxlike_format_slot_drop:n {#2}
            \__docxlike_format_end_commit:n {#2}
          }
          { \__docxlike_bad_value:nnn {#1} {#4} { } }
      }
  }
%    \end{macrocode}
%
% 落槽。行距只写一个端点、单位是 Word 的三种时写成 docx 的 \texttt{w:spacing/@w:line} 与
% \texttt{@w:lineRule}：倍数是 240 分之一，固定值与最小值是缇。别的（区间、文档类登记的单位、
% 段前段后按别的度量给）原样记在 \texttt{ext/}\meta{槽} 那几条里。
%    \begin{macrocode}
\tl_new:N \l__docxlike_format_end_unit_tl
\tl_new:N \l__docxlike_format_end_value_tl
\tl_new:N \l__docxlike_format_end_ii_unit_tl
\tl_new:N \l__docxlike_format_end_ii_value_tl
\cs_new_protected:Npn \__docxlike_format_end_commit:n #1
  {
    \bool_lazy_all:nTF
      {
        { \str_if_eq_p:nn {#1} { line } }
        { \tl_if_empty_p:N \l__docxlike_format_end_ii_unit_tl }
        {
          \str_if_eq_p:Vn \l__docxlike_format_end_unit_tl { auto }
          || \str_if_eq_p:Vn \l__docxlike_format_end_unit_tl { exact }
          || \str_if_eq_p:Vn \l__docxlike_format_end_unit_tl { atLeast }
        }
      }
      {
        \__docxlike_format_put_value:nV { pPr/spacing/lineRule } \l__docxlike_format_end_unit_tl
        \str_if_eq:VnTF \l__docxlike_format_end_unit_tl { auto }
          {
            \__docxlike_format_put_value:ne { pPr/spacing/line }
              { \fp_eval:n { \l__docxlike_format_end_value_tl * 240 } }
          }
          {
            \__docxlike_format_put_value:ne { pPr/spacing/line }
              { \__docxlike_format_twips:V \l__docxlike_format_end_value_tl }
          }
      }
      {
        \__docxlike_format_put_value:nV { ext/#1-unit } \l__docxlike_format_end_unit_tl
        \__docxlike_format_put_value:nV { ext/#1 } \l__docxlike_format_end_value_tl
        \tl_if_empty:NF \l__docxlike_format_end_ii_unit_tl
          {
            \__docxlike_format_put_value:nV { ext/#1-2-unit } \l__docxlike_format_end_ii_unit_tl
            \__docxlike_format_put_value:nV { ext/#1-2 } \l__docxlike_format_end_ii_value_tl
          }
      }
  }
\cs_generate_variant:Nn \__docxlike_format_put_value:nn { nV , ne }
\cs_generate_variant:Nn \__docxlike_format_twips:n { V }
\cs_new_eq:NN \docxlike_format_measure:nnnn \__docxlike_format_measure:nnnn
%    \end{macrocode}
%
% 一个端点。\meta{槽}\meta{后缀}\meta{度量}\meta{词}。度量不是 \texttt{spacing} 的不认词，
% 端点原样记下，单位就是度量的名字。
%    \begin{macrocode}
\cs_new_protected:Npn \__docxlike_format_end:nnnn #1#2#3#4
  {
    \str_if_eq:nnTF {#3} { spacing }
      { \__docxlike_format_end_word:nnn {#1} {#2} {#4} }
      { \__docxlike_format_end_put:nnnn {#1} {#2} {#3} {#4} }
  }
\cs_generate_variant:Nn \__docxlike_format_end:nnnn { nnne }
\cs_new_protected:Npn \__docxlike_format_end_word:nnn #1#2#3
  {
    \str_case:nnF {#3}
      {
        { single } { \__docxlike_format_end_put:nnnn {#1} {#2} { auto } { 1 } }
        { double } { \__docxlike_format_end_put:nnnn {#1} {#2} { auto } { 2 } }
      }
      { \__docxlike_format_end_mark:nnn {#1} {#2} {#3} }
  }
\cs_new_protected:Npn \__docxlike_format_end_mark:nnn #1#2#3
  {
    \str_set:Nx \l__docxlike_format_last_tl { \str_item:nn {#3} { -1 } }
    \str_case:VnF \l__docxlike_format_last_tl
      {
        { ! }
          {
            \__docxlike_format_end_put:nnne {#1} {#2} { exact }
              { \str_range:nnn {#3} { 1 } { -2 } }
          }
        { + }
          {
            \__docxlike_format_end_put:nnne {#1} {#2} { atLeast }
              { \str_range:nnn {#3} { 1 } { -2 } }
          }
      }
      {
        \str_if_in:nVTF { 0123456789. } \l__docxlike_format_last_tl
          { \__docxlike_format_end_put:nnnn {#1} {#2} { auto } {#3} }
          { \__docxlike_format_end_put:nnnn {#1} {#2} { length } {#3} }
      }
  }
\cs_new_protected:Npn \__docxlike_format_end_put:nnnn #1#2#3#4
  {
    \tl_if_empty:nTF {#2}
      {
        \tl_set:Nn \l__docxlike_format_end_unit_tl {#3}
        \tl_set:Nn \l__docxlike_format_end_value_tl {#4}
      }
      {
        \tl_set:Nn \l__docxlike_format_end_ii_unit_tl {#3}
        \tl_set:Nn \l__docxlike_format_end_ii_value_tl {#4}
      }
  }
\cs_generate_variant:Nn \__docxlike_format_end_put:nnnn { nnne }
\cs_generate_variant:Nn \str_case:nnF { VnF }
%    \end{macrocode}
%
% 文档类加的键可以按别的度量写 \texttt{before}、\texttt{after} 这两个槽（经
% \cs{docxlike\_format\_measure:nnnn}），与公开的 \option{before}、\option{after} 说的是同一段间距，
% 同一个样式里都写了就报错，按后写的算。
%
% 走 \texttt{before} 这一路时把区间的第二个端点抹掉：\texttt{1~line} 是
% “数加单位词”两个词，不是两个端点。
%    \begin{macrocode}
\cs_new_protected:Npn \__docxlike_format_space:nn #1#2
  {
    \__docxlike_format_claim:nn {#1} {#1}
    \str_if_eq:nnTF {#2} { auto }
      { \__docxlike_format_slot_restore:n {#1} }
      {
        \__docxlike_format_slot_drop:n {#1}
        \__docxlike_format_amount_put:nnnn { pPr/spacing/#1Lines } { line }
          { pPr/spacing/#1 } {#2}
      }
  }
%    \end{macrocode}
%
% \texttt{ind-left}、\texttt{ind-right} 收 \texttt{2 ch} 或长度，
% \texttt{before} 那一路收 \texttt{1 line} 或长度。两种都是数在前、单位词
% 在后，只有单位词不同，所以走同一个解析器。单位词照英文 Word 的段落对话框：
% 字符是 \texttt{ch}，行是 \texttt{line}，都不变复数（“2 ch”“2 line”）。
%    \begin{macrocode}
\cs_new_protected:Npn \__docxlike_format_amount_put:nnnn #1#2#3#4
  {
    \__docxlike_format_split:n {#4}
    \__docxlike_format_one_word:TF
      { \__docxlike_format_put:ne {#3} { \__docxlike_format_twips:n {#4} } }
      {
        \str_if_eq:VnTF \l__docxlike_format_tail_tl {#2}
          {
            \__docxlike_format_put:ne {#1}
              { \fp_eval:n { \l__docxlike_format_head_tl * 100 } }
          }
          { \__docxlike_bad_value:Vnn \l_keys_key_str {#4} { } }
      }
  }
\cs_new_protected:Npn \__docxlike_format_indent_side:nn #1#2
  {
    \str_if_eq:nnTF {#2} { auto }
      { \__docxlike_format_slot_restore:n { ind-#1 } }
      {
        \__docxlike_format_slot_drop:n { ind-#1 }
        \__docxlike_format_amount_put:nnnn { pPr/ind/#1Chars } { ch } { pPr/ind/#1 } {#2}
      }
  }
%    \end{macrocode}
%
% \texttt{ind-special} 是“下拉框加数值框”，头一个词是模式，其余是数量。
% \texttt{none} 后面不跟数量。
%    \begin{macrocode}
\cs_new_protected:Npn \__docxlike_format_indent_special:n #1
  {
    \__docxlike_format_split:n {#1}
    \str_case:VnF \l__docxlike_format_head_tl
      {
        { auto } { \__docxlike_format_slot_restore:n { ind-special } }
        { none }
          {
            \__docxlike_format_slot_drop:n { ind-special }
            \__docxlike_format_put:nn { pPr/ind/firstLine } { 0 }
          }
        { first-line }
          { \__docxlike_format_special_put:nV { firstLine } \l__docxlike_format_rest_tl }
        { hanging }
          { \__docxlike_format_special_put:nV { hanging } \l__docxlike_format_rest_tl }
      }
      { \__docxlike_bad_value:nnn { indentation / special } {#1}
          { none , first-line , hanging } }
  }
\cs_generate_variant:Nn \__docxlike_format_indent_special:n { e }
\cs_new_protected:Npn \__docxlike_format_special_put:nn #1#2
  {
    \__docxlike_format_slot_drop:n { ind-special }
    \__docxlike_format_amount_put:nnnn { pPr/ind/#1Chars } { ch } { pPr/ind/#1 } {#2}
  }
\cs_generate_variant:Nn \__docxlike_format_special_put:nn { nV }
\keys_define:nn { docxlike / style }
  {
    font .code:n = { \__docxlike_keys_set:nn { docxlike / style / font } {#1} } ,
    paragraph .code:n =
      { \__docxlike_keys_set:nn { docxlike / style / paragraph } {#1} } ,
    tabs .code:n = { \__docxlike_style_tabs:n {#1} } ,
    border .code:n = { \__docxlike_keys_set:nn { docxlike / style / border } {#1} } ,
    numbering .choices:nn = { none }
      { \__docxlike_format_put:nn { pPr/numPr/numId } { 0 } } ,
    unknown .code:n =
      { \__docxlike_unknown_key: } ,
  }
\keys_define:nn { docxlike / style / font }
  {
    asian-text-font .code:n =
      { \__docxlike_format_put:nn { rPr/rFonts/eastAsia } {#1} } ,
    latin-text-font .code:n =
      { \__docxlike_format_put:nn { rPr/rFonts/ascii } {#1} } ,
    font-style .code:n = { \__docxlike_style_font_style:n {#1} } ,
    size .code:n = { \__docxlike_format_size:n {#1} } ,
    all-caps .code:n = { \__docxlike_format_bool:nn { rPr/caps } {#1} } ,
    all-caps .default:n = { true } ,
    kerning-for-fonts .code:n = { \__docxlike_style_kern:n {#1} } ,
    character-spacing .code:n = { \__docxlike_style_char_spacing:n {#1} } ,
    snap-to-grid-when-document-grid-is-defined .code:n = { \__docxlike_format_bool:nn { rPr/snapToGrid } {#1} } ,
    snap-to-grid-when-document-grid-is-defined .default:n = { true } ,
    unknown .code:n =
      { \__docxlike_unknown_key: } ,
  }
\keys_define:nn { docxlike / style / paragraph }
  {
    alignment .code:n = { \__docxlike_style_alignment:n {#1} } ,
    indentation .code:n = { \__docxlike_style_indentation:n {#1} } ,
    spacing .code:n = { \__docxlike_style_spacing:n {#1} } ,
    snap-to-grid-when-document-grid-is-defined .code:n = { \__docxlike_format_bool:nn { pPr/snapToGrid } {#1} } ,
    snap-to-grid-when-document-grid-is-defined .default:n = { true } ,
    pagination .code:n =
      { \__docxlike_keys_set:nn { docxlike / style / paragraph / pagination } {#1} } ,
    unknown .code:n =
      { \__docxlike_unknown_key: } ,
  }
% \option{kerning-for-fonts} 是字体对话框高级页的“为字体调整字间距”：取“多少磅以上”的阈值，
% \texttt{false} 是不勾。存成 docx 的 \texttt{rPr/kern}（半磅，0 是关）。Word 的连续标点缩半格
% （“），”“。”）是它干的，见 \file{docxlike-linebreak.dtx}。
\cs_new_protected:Npn \__docxlike_style_kern:n #1
  {
    \str_case:nnF {#1}
      {
        { false } { \__docxlike_format_put:nn { rPr/kern } { 0 } }
        { auto } { \__docxlike_format_put:nn { rPr/kern } { auto } }
      }
      {
        \__docxlike_format_put:ne { rPr/kern }
          { \fp_eval:n { round ( \dim_to_decimal_in_bp:n { \__docxlike_bp:n {#1} } * 2 ) } }
      }
  }
\cs_new_protected:Npn \__docxlike_style_alignment:n #1
  {
    \str_case:nnF {#1}
      {
        { left } { \__docxlike_format_put:nn { pPr/jc } { left } }
        { centered } { \__docxlike_format_put:nn { pPr/jc } { center } }
        { right } { \__docxlike_format_put:nn { pPr/jc } { right } }
        { justified } { \__docxlike_format_put:nn { pPr/jc } { both } }
        { distributed } { \__docxlike_format_put:nn { pPr/jc } { distribute } }
        { auto } { \__docxlike_format_put:nn { pPr/jc } { auto } }
      }
      {
        \prop_if_in:NnTF \g__docxlike_format_align_extra_prop {#1}
          { \__docxlike_format_put:nn { pPr/jc } {#1} }
          {
            \__docxlike_bad_value:nne { alignment } {#1}
              {
                left , centered , right , justified , distributed
                \prop_map_function:NN \g__docxlike_format_align_extra_prop
                  \__docxlike_style_align_name:nn
              }
          }
      }
  }
\cs_new:Npn \__docxlike_style_align_name:nn #1#2 { , #1 }
\cs_generate_variant:Nn \__docxlike_bad_value:nnn { nne }
\prop_new:N \g__docxlike_format_align_extra_prop
\cs_new_protected:Npn \docxlike_format_alignment_new:nn #1#2
  {
    \prop_gput:Nnn \g__docxlike_format_align_extra_prop {#1} {#2}
    \prop_gput:Nnn \g__docxlike_format_align_prop {#1} {#2}
  }
\keys_define:nn { docxlike / style / paragraph / pagination }
  {
    keep-with-next .code:n = { \__docxlike_format_bool:nn { pPr/keepNext } {#1} } ,
    keep-with-next .default:n = { true } ,
    keep-lines-together .code:n = { \__docxlike_format_bool:nn { pPr/keepLines } {#1} } ,
    keep-lines-together .default:n = { true } ,
    page-break-before .code:n = { \__docxlike_format_bool:nn { pPr/pageBreakBefore } {#1} } ,
    page-break-before .default:n = { true } ,
    unknown .code:n =
      { \__docxlike_unknown_key: } ,
  }
%    \end{macrocode}
%
% \option{font-style} 照“字形”下拉框：\texttt{regular}、\texttt{bold}；\texttt{italic} 与
% \texttt{bold-italic} 还没有做，写了报错，不静默放过。
%    \begin{macrocode}
\msg_new:nnnn { docxlike } { not-implemented }
  { '#1'~is~not~implemented~yet. }
  { The~value~is~ignored. }
\cs_new_protected:Npn \__docxlike_style_font_style:n #1
  {
    \str_case:nnF {#1}
      {
        { regular } { \__docxlike_format_put:nn { rPr/b } { false } }
        { bold } { \__docxlike_format_put:nn { rPr/b } { true } }
        { auto } { \__docxlike_format_put:nn { rPr/b } { auto } }
        { italic } { \msg_error:nnn { docxlike } { not-implemented } { font-style~=~italic } }
        { bold-italic }
          { \msg_error:nnn { docxlike } { not-implemented } { font-style~=~bold-italic } }
      }
      { \__docxlike_bad_value:nnn { font-style } {#1} { regular , bold , italic , bold-italic } }
  }
%    \end{macrocode}
%
% \option{character-spacing} 照“字体 → 高级”的字符间距：\option{spacing} 是
% \texttt{normal}、\texttt{expanded}、\texttt{condensed}，\option{by} 是磅值。
%    \begin{macrocode}
\tl_new:N \l__docxlike_style_cs_spacing_tl
\dim_new:N \l__docxlike_style_cs_by_dim
\keys_define:nn { docxlike / style / font / character-spacing }
  {
    spacing .choices:nn = { normal , expanded , condensed }
      { \tl_set_eq:NN \l__docxlike_style_cs_spacing_tl \l_keys_choice_tl } ,
    by .code:n =
      { \dim_set:Nn \l__docxlike_style_cs_by_dim { \__docxlike_bp:n {#1} } } ,
    unknown .code:n =
      { \__docxlike_unknown_key: } ,
  }
\cs_new_protected:Npn \__docxlike_style_char_spacing:n #1
  {
    \str_if_eq:nnTF {#1} { auto }
      { \__docxlike_format_char_spacing:n { auto } }
      {
        \group_begin:
          \tl_set:Nn \l__docxlike_style_cs_spacing_tl { normal }
          \dim_zero:N \l__docxlike_style_cs_by_dim
          \keys_set:nn { docxlike / style / font / character-spacing } {#1}
          \__docxlike_format_char_spacing:e
            {
              \str_case:Vn \l__docxlike_style_cs_spacing_tl
                {
                  { normal } { 0pt }
                  { expanded } { \dim_use:N \l__docxlike_style_cs_by_dim }
                  { condensed } { \dim_eval:n { - \l__docxlike_style_cs_by_dim } }
                }
            }
        \group_end:
      }
  }
%    \end{macrocode}
%
% \option{indentation} 照“缩进”那一组：\option{left}、\option{right}，\option{special} 是
% \texttt{none}、\texttt{first-line}、\texttt{hanging}，\option{by} 是它的磅值或字数
% （\texttt{2 ch}）。
%    \begin{macrocode}
\tl_new:N \l__docxlike_style_special_tl
\tl_new:N \l__docxlike_style_by_tl
\keys_define:nn { docxlike / style / paragraph / indentation }
  {
    left .code:n = { \__docxlike_format_indent_side:nn { left } {#1} } ,
    right .code:n = { \__docxlike_format_indent_side:nn { right } {#1} } ,
    special .tl_set:N = \l__docxlike_style_special_tl ,
    by .tl_set:N = \l__docxlike_style_by_tl ,
    unknown .code:n =
      { \__docxlike_unknown_key: } ,
  }
\cs_new_protected:Npn \__docxlike_style_indentation:n #1
  {
    \tl_clear:N \l__docxlike_style_special_tl
    \tl_clear:N \l__docxlike_style_by_tl
    \__docxlike_keys_set:nn { docxlike / style / paragraph / indentation } {#1}
    \tl_if_empty:NF \l__docxlike_style_special_tl
      {
        \str_case:VnF \l__docxlike_style_special_tl
          {
            { none } { \__docxlike_format_indent_special:n { none } }
            { auto } { \__docxlike_format_indent_special:n { auto } }
          }
          {
            \__docxlike_format_indent_special:e
              { \l__docxlike_style_special_tl \c_space_tl \l__docxlike_style_by_tl }
          }
      }
  }
%    \end{macrocode}
%
% \option{spacing} 照“间距”那一组：\option{before}、\option{after}，\option{line-spacing} 是
% 下拉框（\texttt{single}、\texttt{1.5-lines}、\texttt{double}、\texttt{at-least}、
% \texttt{exactly}、\texttt{multiple}），\option{at} 是旁边那个“设置值”。
% \option{line-spacing} 要等 \option{at} 读完才能落。文档类往这一组加的键若也写行距那一槽，
% 用 \cs{docxlike\_format\_spacing\_later:n}\marg{代码} 排进同一个队列，这一组读完后按写的
% 先后执行，“两个都写了按后写的算”照旧成立。\option{at} 只收一个值：倍数是数，固定值、
% 最小值是长度。
%    \begin{macrocode}
\tl_new:N \l__docxlike_style_rule_tl
\tl_new:N \l__docxlike_style_at_tl
\tl_new:N \l__docxlike_style_line_tl
\keys_define:nn { docxlike / style / paragraph / spacing }
  {
    before .code:n = { \__docxlike_format_space:nn { before } {#1} } ,
    after .code:n = { \__docxlike_format_space:nn { after } {#1} } ,
    line-spacing .code:n =
      {
        \tl_set:Nn \l__docxlike_style_rule_tl {#1}
        \docxlike_format_spacing_later:n { \__docxlike_style_rule_commit: }
      } ,
    at .tl_set:N = \l__docxlike_style_at_tl ,
    unknown .code:n =
      { \__docxlike_unknown_key: } ,
  }
\cs_new_protected:Npn \docxlike_format_spacing_later:n #1
  { \tl_put_right:Nn \l__docxlike_style_line_tl {#1} }
\cs_new_protected:Npn \__docxlike_style_spacing:n #1
  {
    \tl_clear:N \l__docxlike_style_rule_tl
    \tl_clear:N \l__docxlike_style_at_tl
    \tl_clear:N \l__docxlike_style_line_tl
    \__docxlike_keys_set:nn { docxlike / style / paragraph / spacing } {#1}
    \l__docxlike_style_line_tl
  }
\cs_new_protected:Npn \__docxlike_style_rule_commit:
  {
    \tl_if_empty:NF \l__docxlike_style_rule_tl
      {
        \str_case:VnF \l__docxlike_style_rule_tl
          {
            { single } { \__docxlike_style_rule_put:n { single } }
            { 1.5-lines } { \__docxlike_style_rule_put:n { 1.5 } }
            { double } { \__docxlike_style_rule_put:n { double } }
            { auto } { \__docxlike_style_rule_put:n { auto } }
            { multiple } { \__docxlike_style_at_check:nn { number } { } }
            { exactly } { \__docxlike_style_at_check:nn { length } { ! } }
            { at-least } { \__docxlike_style_at_check:nn { length } { + } }
          }
          {
            \__docxlike_bad_value:nVn { line-spacing } \l__docxlike_style_rule_tl
              { single , 1.5-lines , double , at-least , exactly , multiple }
          }
      }
  }
\cs_new_protected:Npn \__docxlike_style_rule_put:n #1
  { \__docxlike_format_measure:nnnn { line-spacing } { line } { spacing } {#1} }
\cs_generate_variant:Nn \__docxlike_style_rule_put:n { e }
\cs_generate_variant:Nn \__docxlike_bad_value:nnn { nV }
\cs_new_protected:Npn \__docxlike_style_at_check:nn #1#2
  {
    \__docxlike_format_split:V \l__docxlike_style_at_tl
    \str_set:Nx \l__docxlike_format_last_tl
      { \str_item:Vn \l__docxlike_style_at_tl { -1 } }
    \bool_set_false:N \l__docxlike_style_at_ok_bool
    \int_compare:nNnT { \l__docxlike_format_count_int } = { 1 }
      {
        \str_if_in:nVTF { 0123456789. } \l__docxlike_format_last_tl
          { \bool_set_true:N \l__docxlike_style_at_ok_bool }
          { \str_if_eq:nnT {#1} { length } { \bool_set_true:N \l__docxlike_style_at_ok_bool } }
      }
    \bool_if:NTF \l__docxlike_style_at_ok_bool
      {
        \str_if_eq:nnTF {#1} { length }
          { \__docxlike_style_rule_put:e { \__docxlike_bp:V \l__docxlike_style_at_tl #2 } }
          { \__docxlike_style_rule_put:e { \l__docxlike_style_at_tl #2 } }
      }
      {
        \__docxlike_bad_value:nVn { at } \l__docxlike_style_at_tl
          { \str_if_eq:nnTF {#1} { number } { a~number } { a~length } }
      }
  }
\bool_new:N \l__docxlike_style_at_ok_bool
\cs_generate_variant:Nn \__docxlike_format_split:n { V }
\cs_generate_variant:Nn \str_item:nn { Vn }
\cs_generate_variant:Nn \__docxlike_format_measure:nnnn { nnne , nnnV }
%    \end{macrocode}
%
% \option{tabs} 照“制表位”对话框：逗号隔开若干个停靠点，每个写“位置 对齐方式 前导符”，
% 对齐方式 \texttt{left}、\texttt{center}、\texttt{right}，前导符可省，取 \texttt{none}、
% \texttt{dot}、\texttt{hyphen}、\texttt{underscore}、\texttt{middle-dot}。内部按
% “对齐方式 位置 前导符”存，前导符用 docx 的名字。
%    \begin{macrocode}
\seq_new:N \l__docxlike_style_tab_seq
\seq_new:N \l__docxlike_style_tab_out_seq
\cs_new_protected:Npn \__docxlike_style_tabs:n #1
  {
    \str_if_eq:nnTF {#1} { auto }
      { \__docxlike_format_put:nn { pPr/tabs } { auto } }
      {
        \seq_clear:N \l__docxlike_style_tab_out_seq
        \seq_set_from_clist:Nn \l__docxlike_style_tab_seq {#1}
        \seq_map_inline:Nn \l__docxlike_style_tab_seq
          { \__docxlike_style_tab_item:w ##1 ~ ~ \q_stop }
        \__docxlike_format_put:ne { pPr/tabs } { \seq_use:Nn \l__docxlike_style_tab_out_seq { , } }
      }
  }
\tl_new:N \l__docxlike_style_tab_lead_tl
\tl_new:N \l__docxlike_style_tab_lead_in_tl
\cs_new_protected:Npn \__docxlike_style_tab_item:w #1 ~ #2 ~ #3 \q_stop
  {
    \tl_set:Ne \l__docxlike_style_tab_lead_in_tl { \tl_trim_spaces:n {#3} }
    \tl_clear:N \l__docxlike_style_tab_lead_tl
    \tl_if_empty:NF \l__docxlike_style_tab_lead_in_tl
      {
        \__docxlike_ooxml_value:nVnNTF { tabs } \l__docxlike_style_tab_lead_in_tl
          { none , dot , hyphen , underscore , middleDot }
          \l__docxlike_style_tab_lead_tl { }
          {
            \__docxlike_bad_value:nVn { tabs } \l__docxlike_style_tab_lead_in_tl
              { none , dot , hyphen , underscore , middle-dot }
          }
      }
    \seq_put_right:Ne \l__docxlike_style_tab_out_seq
      {
        #2 ~ \__docxlike_format_twips:n {#1}
        \tl_if_empty:NF \l__docxlike_style_tab_lead_tl { ~ \l__docxlike_style_tab_lead_tl }
      }
  }
%    \end{macrocode}
%
% \option{border} 照“边框和底纹”对话框：\option{top}、\option{bottom} 各一组，
% \option{style} 线型（界面上只有图样，取 docx 的线型名改成小写加连字符，
% \texttt{thin-thick-small-gap}；\texttt{none} 是无框线）、\option{color}、
% \option{width} 线宽（长度）、\option{from-text} 距正文。默认单线、自动色、0.5 磅、1 磅。
%    \begin{macrocode}
\tl_new:N \l__docxlike_style_border_style_tl
\tl_new:N \l__docxlike_style_border_color_tl
\dim_new:N \l__docxlike_style_border_width_dim
\dim_new:N \l__docxlike_style_border_text_dim
\keys_define:nn { docxlike / style / border }
  {
    top .code:n = { \__docxlike_style_border:nn { top } {#1} } ,
    bottom .code:n = { \__docxlike_style_border:nn { bottom } {#1} } ,
    unknown .code:n =
      { \__docxlike_unknown_key: } ,
  }
\keys_define:nn { docxlike / style / border / side }
  {
    style .code:n = { \__docxlike_style_border_style:n {#1} } ,
    color .tl_set:N = \l__docxlike_style_border_color_tl ,
    width .code:n =
      { \dim_set:Nn \l__docxlike_style_border_width_dim { \__docxlike_bp:n {#1} } } ,
    from-text .code:n =
      { \dim_set:Nn \l__docxlike_style_border_text_dim { \__docxlike_bp:n {#1} } } ,
    unknown .code:n =
      { \__docxlike_unknown_key: } ,
  }
\clist_const:Nn \c__docxlike_border_style_clist
  {
    single , thick , double , triple , dotted , dashed , dashSmallGap , dotDash ,
    dotDotDash , thinThickSmallGap , thickThinSmallGap , thinThickThinSmallGap ,
    thinThickMediumGap , thickThinMediumGap , thinThickThinMediumGap ,
    thinThickLargeGap , thickThinLargeGap , thinThickThinLargeGap , wave ,
    doubleWave , dashDotStroked , threeDEmboss , threeDEngrave , outset , inset
  }
\cs_new_protected:Npn \__docxlike_style_border_style:n #1
  {
    \str_if_eq:nnTF {#1} { none }
      { \tl_set:Nn \l__docxlike_style_border_style_tl { nil } }
      {
        \__docxlike_ooxml_value:nnVNTF { style } {#1} \c__docxlike_border_style_clist
          \l__docxlike_style_border_style_tl { }
          { \__docxlike_bad_value:nnn { style } {#1} { } }
      }
  }
\cs_new_protected:Npn \__docxlike_style_border:nn #1#2
  {
    \str_if_eq:nnTF {#2} { auto }
      { \__docxlike_format_put:nn { pPr/pBdr/#1 } { auto } }
      {
        \group_begin:
          \tl_set:Nn \l__docxlike_style_border_style_tl { single }
          \tl_set:Nn \l__docxlike_style_border_color_tl { auto }
          \dim_set:Nn \l__docxlike_style_border_width_dim { 0.5bp }
          \dim_set:Nn \l__docxlike_style_border_text_dim { 1bp }
          \keys_set:nn { docxlike / style / border / side } {#2}
          \__docxlike_format_put:ne { pPr/pBdr/#1 }
            {
              { \l__docxlike_style_border_style_tl }
              { \fp_eval:n { round ( \dim_to_decimal_in_bp:n { \l__docxlike_style_border_width_dim } * 8 , 3 ) } }
              { \fp_eval:n { round ( \dim_to_decimal_in_bp:n { \l__docxlike_style_border_text_dim } , 3 ) } }
              { \l__docxlike_style_border_color_tl }
            }
        \group_end:
      }
  }
%    \end{macrocode}
%
% \emph{限定词挂在样式名上，不挂在属性里。} 条件是整份文档的事，一个样式说了不算，所以
% \meta{样式}\texttt{/}\meta{限定词} 被当成另一个目标来收，就像两个不同的样式。
%
% 用 \cs{group\_begin:} 圈起来，出了组当前限定词自动复原。存储那一头全是全局
% 赋值，不受组的影响。判重的表在 \cs{\_\_docxlike\_format\_target:nn} 里清，
% 所以带限定词的与不带的各判各的，不算冲突。
%    \begin{macrocode}
\cs_new_protected:Npn \__docxlike_format_layer:nnn #1#2#3
  {
    \group_begin:
      \tl_set:Nn \l__docxlike_format_layer_tl {#1}
      \__docxlike_format_target:nn {#2} {#3}
    \group_end:
  }
%    \end{macrocode}
%
% \emph{网格分两个方向。} Word 的 \texttt{docGrid} 同时定行距与字距，
% 但一个段落要不要吃这两样是分开的：样本文档的正文吃字距（每字 0.4543\,bp），
% 行距却按自然行高排（实测 19.44\,bp，不是一格 19.55\,bp）。所以这里拆成两轴。
%
% 写成布尔或 \texttt{auto} 是两轴的简写，两轴一起设；要分开设就写子块。
% 两种写法都收，因为绝大多数场合两轴同进同退，为了少数场合逼所有人写子块不划算。
%
% \emph{两轴的 \texttt{auto} 都是“跟着所在的位置走”。} 一个元素贴不贴格，
% Word 里看的是它所在那一段的对话框，而段落是嵌在版面里的：表题嵌在正文里、
% 子图题嵌在图题里、封面的各个字段嵌在封面里。所以每个目标有一个上级
% （\cs{c\_\_docxlike\_format\_parent\_prop}），\texttt{auto} 就沿着上级往上找，找到
% 第一个明写了 \texttt{true}／\texttt{false} 的为准；各个根上也没写时按不贴算，
% 网格是要明说才贴的东西。默认表没有一处贴纵向网格时，这条只是把“全局为假”
% 换成了“沿位置继承出来的假”；哪一页开了网格，那一页上的题注、表格、空段自然跟着，不必逐个再写。
% 之前纵轴的 \texttt{auto} 一律求值成假、横轴一律取正文，是临时写法。
% 用户写的 \texttt{auto} 仍先退回默认表（见 \cs{\_\_docxlike\_format\_put:nn}），默认表
% 明写了 \texttt{true}／\texttt{false} 就以它为准；只有默认表也没写的轴才往上找。
% 所以“跟着位置走”是默认表的写法：哪一档该跟上级，默认表就别给它写死。
%
% 横轴的 \texttt{spacing} 另有一层：没写它时这一档一个字也不发
% \cs{CJKglue}，靠继承拿到上级的字距，结果一样；写了就得发，发的时候网格那一截
% 照样按上级的来，只是再加上本档的字距增减。
%
% 靠“不发”来实现“跟着走”不行：一旦 \texttt{spacing} 逼着发，网格那一截
% 按假算就等于悄悄关掉了网格，章标题（\texttt{char-spacing = -0.4bp}）的字距
% 会比样本文档窄一整格网格余量。
%
% \emph{标题吃横向网格。} 样本文档里 \texttt{heading} 各级样式都写着
% \texttt{snapToGrid=0}，但逐页量下来，同一页上节标题与正文的每字增量总是相同：
% 前半那几页都是 $+0.6720$，后半那几页都是 $+0.4536$，正好是两节
% \texttt{docGrid} 的 \texttt{charSpace} 除以 4096。章标题也跟着换档，
% 只是每档都比正文低 $0.400$，那是它自己的 \texttt{w:spacing}。
% 所以 \texttt{snapToGrid=0} 在那个渲染里只关掉了纵轴。
%    \begin{macrocode}
\cs_new_protected:Npn \__docxlike_format_target:nn #1#2
  {
    \str_set:Nn \l__docxlike_format_target_str {#1}
    \prop_clear:N \l__docxlike_format_claim_prop
    \__docxlike_keys_set:nn { docxlike / style } {#2}
    \__docxlike_format_after_target:n {#1}
  }
%    \end{macrocode}
%
% \emph{样式集合。} \option{styles} 里每一项是一个样式：键是 Word 里的样式名
% （\texttt{Normal}、\texttt{Heading 1}），值是它的格式。样式名是名字，照 Word 原样写；
% 找样式时不计大小写、空格与连字符，\texttt{Heading 1}、\texttt{Heading1}、\texttt{heading-1}
% 是同一个。新样式的内部标识（docx 的 \texttt{w:styleId}）是去掉空格的名字。
% 没登记过的名字第一次写到就新建一个样式，与在 Word 里“新建样式”一样。
%
% 样式名后面带 \texttt{/}\meta{限定词} 的，限定词要由文档类登记过，见“限定词”。
%    \begin{macrocode}
\str_new:N \l__docxlike_style_name_str
\tl_new:N \l__docxlike_style_found_tl
\tl_new:N \l__docxlike_style_layer_tl
\msg_new:nnnn { docxlike } { unknown-qualifier }
  { The~style~name~'#1/#2'~carries~a~qualifier~no~one~declared. }
  { Qualifiers~come~from~the~document~class;~the~setting~is~ignored. }
\bool_new:N \l__docxlike_style_layer_known_bool
\cs_new_protected:Npn \__docxlike_styles_unknown:nn #1#2
  {
    \str_set:Nn \l__docxlike_style_name_str {#1}
    \str_remove_all:Nn \l__docxlike_style_name_str { ~ }
    \tl_clear:N \l__docxlike_style_layer_tl
    \str_if_in:NnT \l__docxlike_style_name_str { / }
      {
        \exp_last_unbraced:NV \__docxlike_styles_split:w
          \l__docxlike_style_name_str \q_stop
      }
    \__docxlike_squash:VN \l__docxlike_style_name_str \l__docxlike_value_try_str
    \prop_get:NVNT \g__docxlike_format_name_prop \l__docxlike_value_try_str
      \l__docxlike_style_found_tl
      { \str_set:NV \l__docxlike_style_name_str \l__docxlike_style_found_tl }
    \bool_set_true:N \l__docxlike_style_layer_known_bool
    \tl_if_empty:NF \l__docxlike_style_layer_tl
      {
        \bool_set_false:N \l__docxlike_style_layer_known_bool
        \clist_map_inline:Nn \g__docxlike_format_layer_clist
          {
            \str_if_eq:VnT \l__docxlike_style_layer_tl {##1}
              { \bool_set_true:N \l__docxlike_style_layer_known_bool }
          }
      }
    \bool_if:NTF \l__docxlike_style_layer_known_bool
      { \__docxlike_styles_store:n {#2} }
      {
        \msg_error:nnVV { docxlike } { unknown-qualifier }
          \l__docxlike_style_name_str \l__docxlike_style_layer_tl
      }
  }
\cs_new_protected:Npn \__docxlike_styles_split:w #1 / #2 \q_stop
  {
    \str_set:Nn \l__docxlike_style_name_str {#1}
    \tl_set:Nn \l__docxlike_style_layer_tl {#2}
  }
\cs_generate_variant:Nn \msg_error:nnnn { nnVV }
\cs_new_protected:Npn \__docxlike_styles_store:n #1
  {
    \prop_if_exist:cF { g__docxlike_format_ \l__docxlike_style_name_str _prop }
      { \docxlike_format_target_new:Vn \l__docxlike_style_name_str { Normal } }
    \tl_if_empty:NTF \l__docxlike_style_layer_tl
      { \__docxlike_format_target:Vn \l__docxlike_style_name_str {#1} }
      {
        \__docxlike_format_layer:VVn \l__docxlike_style_layer_tl
          \l__docxlike_style_name_str {#1}
      }
  }
\cs_generate_variant:Nn \__docxlike_format_target:nn { Vn }
\cs_generate_variant:Nn \__docxlike_format_layer:nnn { VVn }
\msg_new:nnn { docxlike } { style-without-value }
  { The~style~'#1'~was~given~without~a~value;~it~is~ignored. }
\cs_new_protected:Npn \docxlike_format_set:n #1
  {
    \tl_set:Nn \l__docxlike_keys_tl {#1}
    \tl_replace_all:Nnn \l__docxlike_keys_tl { \par } { }
    \exp_args:NnnV \keyval_parse:nnn
      { \msg_warning:nnn { docxlike } { style-without-value } }
      { \__docxlike_styles_item:nn } \l__docxlike_keys_tl
  }
\cs_new_protected:Npn \__docxlike_styles_item:nn #1#2
  { \__docxlike_styles_unknown:nn {#1} {#2} }
\keys_define:nn { docxlike }
  { styles .code:n = { \docxlike_format_set:n {#1} } }
\docxlike_format_target_new:nn { Normal } { }
\clist_map_inline:nn { Heading1 , Heading2 , Heading3 , Heading4 , Header , Footer }
  { \docxlike_format_target_new:nn {#1} { Normal } }
%    \end{macrocode}
%
% \option{hyphenation} 照“布局 → 断字”：\texttt{automatic} 自动断字，\texttt{none} 不断。
% 全文一个开关，正文、题注、摘要、纯西文段一律跟随。
%    \begin{macrocode}
\bool_new:N \g_docxlike_format_hyphen_bool
\keys_define:nn { docxlike }
  {
    hyphenation .choices:nn = { automatic , none }
      {
        \str_if_eq:VnTF \l_keys_choice_tl { automatic }
          { \bool_gset_true:N \g_docxlike_format_hyphen_bool }
          { \bool_gset_false:N \g_docxlike_format_hyphen_bool }
      } ,
  }

%    \end{macrocode}
%
% 一个目标设完之后调 \cs{\_\_docxlike\_format\_after\_target:n}\marg{目标}。文档类要在这里
% 读回设置、同步自己的开关（比如旧写法的选项与新键指的是同一件事），就用
% \cs{docxlike\_format\_set\_after\_target:n}\marg{代码} 给出代码，代码里 \texttt{\#1}
% 是目标名，读设置用 \cs{docxlike\_format\_get:nnN}。默认什么也不做。整个目标设完才调，
% 因为同一目标里的键可能前后改来改去，要看的是最终值。
%    \begin{macrocode}
\cs_new_protected:Npn \__docxlike_format_after_target:n #1 { }
\cs_new_protected:Npn \docxlike_format_set_after_target:n #1
  { \cs_gset_protected:Npn \__docxlike_format_after_target:n ##1 {#1} }
%    \end{macrocode}
%
% \cs{docxlike\_format\_get:nnN}\marg{目标}\marg{键}\meta{记号列表} 取一条设置，没设过就
% 让记号列表为空。排版那边接进来时用它。
%    \begin{macrocode}
\cs_new_protected:Npn \docxlike_format_get:nnN #1#2#3
  { \__docxlike_format_get:nnNF {#1} {#2} #3 { \tl_clear:N #3 } }
%    \end{macrocode}
%
% \subsection{给文档类的接入点}
%
% 公开的键只有 Word 里有的。某个文档类的格式规范要的别的东西（按视觉留白给间距、
% 区间、按条件换一套取值……）由文档类自己加，本包留下这些口子：
% \begin{itemize}
% \item 往 \texttt{docxlike / style / font}、\texttt{docxlike / style / paragraph}、
%   \texttt{docxlike / style / paragraph / spacing} 这几族加键，与本包的键写在同一串里；
% \item \cs{docxlike\_format\_put:nn}\marg{条目}\marg{值}：往当前样式写一条存储，值写
%   \texttt{auto} 是退回默认；读用 \cs{docxlike\_format\_get:nnN}；
% \item \cs{docxlike\_format\_measure:nnnn}\marg{键名}\marg{槽}\marg{度量}\marg{取值}：写行距
%   （槽 \texttt{line}）或段前段后（\texttt{before}、\texttt{after}）。度量是 \texttt{spacing} 时
%   按行距那张表认词，别的度量把端点的单位直接记成度量的名字；取值写两个词是区间。
%   同一个样式里几个键写同一个槽要报错，按后写的算，报错里用\meta{键名}；
% \item \cs{docxlike\_format\_spacing\_later:n}\marg{代码}：写行距那一槽的键排进
%   \option{line-spacing} 的队列，这一组读完按写的先后执行；
% \item \cs{docxlike\_format\_line\_unit\_new:nn}\marg{单位}\marg{代码}、
%   \cs{docxlike\_format\_space\_unit\_new:nn}\marg{单位}\marg{代码}：端点的单位是这个名字时，
%   \meta{代码}把值 \texttt{\#1} 折成长度设进 \texttt{\#2}。这时字号与自然行高已在
%   \cs{l\_docxlike\_format\_size\_dim}、\cs{l\_docxlike\_format\_natural\_dim}；段前段后那一种还有
%   \cs{l\_docxlike\_format\_style\_tl}（样式名）与 \cs{l\_docxlike\_format\_side\_tl}
%   （\texttt{before} 或 \texttt{after}）；
% \item \cs{docxlike\_format\_qualifiers:nn}：样式名的限定词，见“限定词”；
% \item \cs{docxlike\_format\_alignment\_new:nn}\marg{取值}\marg{代码}：给 \option{alignment}
%   加一个取值，排的时候发\meta{代码}；
% \item \cs{docxlike\_format\_set\_target\_hook:n}\marg{代码}：按样式算版面、发字体之前执行，
%   \texttt{\#1} 是样式名；要这一样式的行盒按西文算，就在这里立
%   \cs{l\_docxlike\_format\_latin\_line\_bool}；
% \item \cs{docxlike\_format\_set\_font\_hook:n}\marg{代码}：发字体命令时在加粗之后执行，
%   \texttt{\#1} 是样式名，要追加的命令交给 \cs{docxlike\_format\_emit:n}；
% \item 上面说过的 \cs{docxlike\_format\_set\_after\_target:n} 与
%   \cs{docxlike\_format\_add\_font:nn}。
% \end{itemize}
%
% \subsection{换算}
%
% 把收下来的键算成排版参数。规则是从 Word 里量出来的，四轮对照文档、三十组、
% 一百六十多条基线，量法与来路记在 \file{docxlike-page.dtx}。
%
% \emph{字体常数。} 自然行高与上升部都与字号成正比，系数只跟字体有关，由字体文件里的
% 度量照 Windows Word 的规则算，见 \file{docxlike-metrics.dtx} 的“按字体名取度量”。中易宋体是
% 1.296875 与 1.0078125，Times New Roman 是 1.1499 与 0.9336，苹方的行高是 1.82，差了四成，
% 所以必须按字体查。首行基线就落在上升部处，Windows 不另加常数。
%
% \emph{纯西文的段落行距比中文小。} Word 里一行的行高取这一行\emph{实际用到}的
% 各个字体的最大值，所以同一个设置下，有汉字的行按东亚字体（中易 1.297）算，纯英文的行
% 只有西文字体参与（Times 1.15）。样本文档的两份摘要就是这样：中英文摘要的段落设置相同
% （多倍行距 1.25、不对齐到网格），中文那份量出行距 19.455\,bp，英文那份 17.25\,bp，
% 正好是 $1.25 \times 1.297 \times 12$ 与 $1.25 \times 1.15 \times 12$。
%
% \hologo{LaTeX} 静态看不出一段里有没有汉字，所以要有人说一声。逐条内容不同的地方
% （目录的条目）由排版代码临时立 \cs{l\_docxlike\_format\_latin\_pick\_bool}，
% \cs{docxlike\_format\_compute:n} 就按西文字体算这一次，判据见 \cs{docxlike\_format\_cjk:nTF}。
%
% \emph{按样式说。} 文档类在 \cs{docxlike\_format\_set\_target\_hook:n} 里按样式立
% \cs{l\_docxlike\_format\_latin\_line\_bool}，那个样式的行盒就按西文字体算，东亚字体照发。
%
% \emph{行盒高。} 也就是 \cs{baselineskip}。倍数写法在对齐到网格时是网格行距的
% 倍数，关掉时是自然行高的倍数；一行装不下就往上取整到整数个格。固定值原样生效，
% 网格与字号都不介入，比字号小也照给（字会叠）。最小值是设定值与「网格开看网格
% 行距、网格关看自然行高」两者取大。
%
% \emph{基线在盒里的位置。} 对齐到网格时行盒居中，基线在盒顶下方
% “行盒高减自然行高的一半，再加上升部”；关掉时顶对齐，倍数多出来的量全落在基线之下，
% 基线就在上升部处。固定值是另一套，基线在盒顶下方 0.8 倍盒高，与网格开关无关。
%
% \emph{段前段后。} “行”这个单位恒指网格行距，与该段自己的行距无关，也与这一段
% 对不对齐网格无关；磅数原样加，不吸到网格。相邻两段取段后与段前的较大者，
% 不相加。页首的段前与页尾的段后由 \hologo{TeX} 自己在分页处丢弃，与 Word 一致。
%
% 折成长度以后要向下取整到整缇。Word 的 \texttt{w:beforeLines} 是百分之一行，
% 存盘时它把“系数乘网格行距”截断成整缇写进 \texttt{w:before}：网格行距 391 缇
% 时，1 行是 391，0.8 行 312.8 截成 312，0.5 行 195.5 截成 195。样本文档的
% \texttt{docx} 里章标题写的正是
% \texttt{before="391" after="312"}，节条标题是 195，与这条规则一致；标定文档里
% 把同一种间距重复二十遍回归，也是这三个数。
%    \begin{macrocode}
\docxlike_font_metric:nNNF { SimSun } \l_tmpa_tl \l_tmpb_tl
  {
    \tl_set:Nn \l_tmpa_tl { 1.296875 }
    \tl_set:Nn \l_tmpb_tl { 1.0078125 }
  }
\tl_const:Ne \c_docxlike_format_natural_tl { \l_tmpa_tl }
\tl_const:Ne \c_docxlike_format_ascent_tl { \l_tmpb_tl }
\docxlike_font_metric:nNNF { Times~New~Roman } \l_tmpa_tl
  \l_tmpb_tl
  {
    \tl_set:Nn \l_tmpa_tl { 1.14990234375 }
    \tl_set:Nn \l_tmpb_tl { 0.93359375 }
  }
\tl_const:Ne \c_docxlike_format_natural_latin_tl { \l_tmpa_tl }
\tl_const:Ne \c_docxlike_format_ascent_latin_tl { \l_tmpb_tl }
\tl_new:N \l_docxlike_format_natural_latin_tl
\tl_set_eq:NN \l_docxlike_format_natural_latin_tl \c_docxlike_format_natural_latin_tl
\tl_new:N \l_docxlike_format_ascent_latin_tl
\tl_set_eq:NN \l_docxlike_format_ascent_latin_tl \c_docxlike_format_ascent_latin_tl
\tl_new:N \l__docxlike_format_font_tl
\bool_new:N \l__docxlike_format_latin_bool
\bool_new:N \l_docxlike_format_latin_pick_bool
\bool_new:N \l_docxlike_format_latin_line_bool
\tl_new:N \l__docxlike_format_latin_tl
\bool_new:N \l__docxlike_format_exact_bool
%    \end{macrocode}
%
% 东亚字体与西文字体的两对系数由样式里的字体名查，见后面“字体”那一段。
%    \begin{macrocode}
\tl_new:N \l_docxlike_format_ascent_tl
\tl_set_eq:NN \l_docxlike_format_ascent_tl \c_docxlike_format_ascent_tl
\tl_new:N \l__docxlike_format_strut_natural_tl
\tl_set_eq:NN \l__docxlike_format_strut_natural_tl \c_docxlike_format_natural_tl
\tl_const:Nn \c_docxlike_format_exact_tl { 0.8 }
%    \end{macrocode}
%
% \emph{缩进。} \texttt{ind-special} 的三种模式折成一个带正负号的量：
% \texttt{first-line} 为正、\texttt{hanging} 为负、\texttt{none} 为零，正好是
% \hologo{LaTeX} 里 \cs{parindent} 与 \cs{hangindent} 的符号约定。
% “一字”这个单位指网格字距，与该段自己的字号无关，与“一行”指网格行距同理。
%    \begin{macrocode}
\dim_new:N \l_docxlike_format_size_dim
\dim_new:N \l_docxlike_format_natural_dim
\dim_new:N \l_docxlike_format_lead_dim
\dim_new:N \l_docxlike_format_first_dim
\dim_new:N \l_docxlike_format_before_dim
\dim_new:N \l_docxlike_format_before_plus_dim
\dim_new:N \l_docxlike_format_after_dim
\dim_new:N \l_docxlike_format_after_plus_dim
\dim_new:N \l_docxlike_format_indent_dim
\dim_new:N \l_docxlike_format_stretch_dim
\dim_new:N \l_docxlike_format_char_spacing_dim
\bool_new:N \l_docxlike_format_snap_vertical_bool
\bool_new:N \l_docxlike_format_snap_horizontal_bool
\bool_new:N \l_docxlike_format_hglue_set_bool
\tl_new:N \l_docxlike_format_line_value_tl
\tl_new:N \l_docxlike_format_line_unit_tl
%    \end{macrocode}
%
%    \begin{macrocode}
\tl_new:N \l__docxlike_format_value_tl
\tl_new:N \l__docxlike_format_unit_tl
\tl_new:N \l__docxlike_format_value_ii_tl
\tl_new:N \l__docxlike_format_unit_ii_tl
\tl_new:N \l__docxlike_format_coef_tl
\dim_new:N \l__docxlike_format_floor_dim
\dim_new:N \l__docxlike_format_span_dim
\cs_generate_variant:Nn \prop_get:NnNF { NVNF }
\cs_new_protected:Npn \__docxlike_format_read:nnnN #1#2#3#4
  { \__docxlike_format_get:nnNF {#1} {#2} #4 { \tl_set:Nn #4 {#3} } }
%    \end{macrocode}
%
% 带单位的键：\meta{值}加\meta{单位}两条。\texttt{line} 折成一行的倍数，
% \texttt{chars} 折成网格字距的倍数，文档类登记过的单位按它的代码折，其余当长度。
%
% \emph{“一行”有多长，看网格启不启用。} 启用时是网格行距；关着时 Word 仍认
% “段前 0.5 行”，那个行是这一节\emph{正文的单倍行高}，也就是 \texttt{Normal}
% 那一档的字号。\file{docGrid} 里记着的 \texttt{linePitch} 此时只是行高的四舍五入，
% 不是跨度——有两份样本写着 391 却没有 \texttt{w:type}，就是这一档。
%    \begin{macrocode}
\dim_new:N \l__docxlike_format_unit_dim
\cs_new_protected:Npn \__docxlike_format_line_unit:N #1
  {
    \bool_if:NTF \g_docxlike_grid_bool
      { \dim_set_eq:NN #1 \g_docxlike_page_lead_dim }
      {
        \__docxlike_format_sz_read:nN { Normal } \l__docxlike_format_value_tl
        \tl_if_empty:NTF \l__docxlike_format_value_tl
          { \dim_set_eq:NN #1 \g_docxlike_page_lead_dim }
          { \dim_set:Nn #1 { \l__docxlike_format_value_tl } }
      }
  }
%    \end{macrocode}
%
%    \begin{macrocode}
\cs_new_protected:Npn \__docxlike_format_amount_use:nnN #1#2#3
  {
    \str_case:nnF {#1}
      {
        { line }
          {
            \__docxlike_format_line_unit:N \l__docxlike_format_unit_dim
            \dim_set:Nn #3
              {
                \fp_to_dim:n
                  {
                    floor
                      (
                        \dim_to_fp:n { #2 \l__docxlike_format_unit_dim }
                        / \dim_to_fp:n { 0.05bp } + 0.0001
                      )
                    * \dim_to_fp:n { 0.05bp }
                  }
              }
          }
        { chars } { \dim_set:Nn #3 { #2 \g_docxlike_page_char_dim } }
      }
      {
        \cs_if_exist:cTF { __docxlike_format_space_unit_ #1 :nN }
          { \use:c { __docxlike_format_space_unit_ #1 :nN } {#2} #3 }
          { \dim_set:Nn #3 {#2} }
      }
  }
\cs_generate_variant:Nn \__docxlike_format_amount_use:nnN { VVN }
\tl_new:N \l__docxlike_format_raw_tl
\cs_new_protected:Npn \__docxlike_format_amount_read:nn #1#2
  {
    \__docxlike_format_slot_find:nn {#1} {#2}
    \tl_set:Nn \l__docxlike_format_value_tl { 0pt }
    \tl_set:Nn \l__docxlike_format_unit_tl { length }
    \tl_clear:N \l__docxlike_format_value_ii_tl
    \tl_clear:N \l__docxlike_format_unit_ii_tl
    \__docxlike_format_ext_read:nTF {#2}
      { }
      {
        \clist_map_inline:Nn \l__docxlike_format_slot_clist
          {
            \prop_get:NnNT \l__docxlike_format_slot_prop {##1} \l__docxlike_format_raw_tl
              {
                \__docxlike_format_attr_amount:nV {##1} \l__docxlike_format_raw_tl
                \clist_map_break:
              }
          }
      }
  }
%    \end{macrocode}
%
% 一条 docx 属性折回“值加单位”：名字以 \texttt{Lines} 收尾的是百分之一行，以
% \texttt{Chars} 收尾的是百分之一字，其余是缇。
%    \begin{macrocode}
\cs_new_protected:Npn \__docxlike_format_attr_amount:nn #1#2
  {
    \str_case:enF { \str_range:nnn {#1} { -5 } { -1 } }
      {
        { Lines }
          {
            \tl_set:Nn \l__docxlike_format_unit_tl { line }
            \tl_set:Ne \l__docxlike_format_value_tl { \fp_eval:n { #2 / 100 } }
          }
        { Chars }
          {
            \tl_set:Nn \l__docxlike_format_unit_tl { chars }
            \tl_set:Ne \l__docxlike_format_value_tl { \fp_eval:n { #2 / 100 } }
          }
      }
      {
        \tl_set:Nn \l__docxlike_format_unit_tl { length }
        \tl_set:Ne \l__docxlike_format_value_tl { \__docxlike_format_from_twips:n {#2} }
      }
  }
\cs_generate_variant:Nn \__docxlike_format_attr_amount:nn { nV }
\prg_new_protected_conditional:Npnn \__docxlike_format_ext_read:n #1 { TF }
  {
    \prop_get:NnNTF \l__docxlike_format_slot_prop { ext/#1 } \l__docxlike_format_value_tl
      {
        \prop_get:NnNF \l__docxlike_format_slot_prop { ext/#1-unit } \l__docxlike_format_unit_tl
          { \tl_set:Nn \l__docxlike_format_unit_tl { length } }
        \prop_get:NnNF \l__docxlike_format_slot_prop { ext/#1-2 } \l__docxlike_format_value_ii_tl
          { \tl_clear:N \l__docxlike_format_value_ii_tl }
        \prop_get:NnNF \l__docxlike_format_slot_prop { ext/#1-2-unit } \l__docxlike_format_unit_ii_tl
          { \tl_clear:N \l__docxlike_format_unit_ii_tl }
        \prg_return_true:
      }
      { \prg_return_false: }
  }
\cs_new_protected:Npn \__docxlike_format_line_read:n #1
  {
    \__docxlike_format_slot_find:nn {#1} { line }
    \tl_clear:N \l__docxlike_format_value_ii_tl
    \tl_clear:N \l__docxlike_format_unit_ii_tl
    \__docxlike_format_ext_read:nTF { line }
      { }
      {
        \prop_get:NnNTF \l__docxlike_format_slot_prop { pPr/spacing/line } \l__docxlike_format_raw_tl
          {
            \prop_get:NnNF \l__docxlike_format_slot_prop { pPr/spacing/lineRule }
              \l__docxlike_format_unit_tl
              { \tl_set:Nn \l__docxlike_format_unit_tl { auto } }
            \str_if_eq:VnTF \l__docxlike_format_unit_tl { auto }
              {
                \tl_set:Ne \l__docxlike_format_value_tl
                  { \fp_eval:n { \l__docxlike_format_raw_tl / 240 } }
              }
              {
                \tl_set:Ne \l__docxlike_format_value_tl
                  { \__docxlike_format_from_twips:V \l__docxlike_format_raw_tl }
              }
          }
          {
            \tl_set:Nn \l__docxlike_format_value_tl { 1 }
            \tl_set:Nn \l__docxlike_format_unit_tl { auto }
          }
      }
  }
\cs_new_protected:Npn \__docxlike_format_sz_read:nN #1#2
  {
    \__docxlike_format_get:nnNTF {#1} { rPr/sz } #2
      { \tl_set:Ne #2 { \__docxlike_format_from_half_points:V #2 } }
      { \tl_clear:N #2 }
  }
\cs_generate_variant:Nn \__docxlike_format_from_twips:n { V }
\cs_generate_variant:Nn \__docxlike_format_from_half_points:n { V }
\cs_new_protected:Npn \__docxlike_format_amount_pick:nnN #1#2#3
  {
    \__docxlike_format_amount_read:nn {#1} {#2}
    \__docxlike_format_amount_use:VVN \l__docxlike_format_unit_tl
      \l__docxlike_format_value_tl #3
  }
%    \end{macrocode}
%
%
% 段前段后也收区间（文档类加的键写得出来）：小端做自然值，大端减小端做伸量，与行距
% 那边同一个规矩。算之前把样式名与方向放进 \cs{l\_docxlike\_format\_style\_tl}、
% \cs{l\_docxlike\_format\_side\_tl}，文档类登记的单位要用。
%    \begin{macrocode}
\cs_new_protected:Npn \__docxlike_format_space_get:nnnNN #1#2#3#4#5
  {
    \__docxlike_format_amount_read:nn {#1} {#2}
    \tl_set:Nn \l_docxlike_format_style_tl {#1}
    \tl_set:Nn \l_docxlike_format_side_tl {#3}
    \__docxlike_format_amount_use:VVN \l__docxlike_format_unit_tl
      \l__docxlike_format_value_tl #4
    \dim_zero:N #5
    \tl_if_empty:NF \l__docxlike_format_value_ii_tl
      {
        \__docxlike_format_amount_use:VVN \l__docxlike_format_unit_ii_tl
          \l__docxlike_format_value_ii_tl \l__docxlike_format_span_dim
        \dim_compare:nNnT { \l__docxlike_format_span_dim } < { #4 }
          {
            \dim_set_eq:NN \l__docxlike_format_floor_dim #4
            \dim_set_eq:NN #4 \l__docxlike_format_span_dim
            \dim_set_eq:NN \l__docxlike_format_span_dim \l__docxlike_format_floor_dim
          }
        \dim_set:Nn #5 { \l__docxlike_format_span_dim - #4 }
      }
  }
\cs_new_protected:Npn \__docxlike_format_indent:n #1
  {
    \__docxlike_format_slot_find:nn {#1} { ind-special }
    \dim_zero:N \l_docxlike_format_indent_dim
    \clist_map_inline:nn { hangingChars , hanging , firstLineChars , firstLine }
      {
        \prop_get:NnNT \l__docxlike_format_slot_prop { pPr/ind/##1 } \l__docxlike_format_raw_tl
          {
            \__docxlike_format_attr_amount:nV {##1} \l__docxlike_format_raw_tl
            \__docxlike_format_amount_use:VVN \l__docxlike_format_unit_tl
              \l__docxlike_format_value_tl \l_docxlike_format_indent_dim
            \str_if_eq:eeT { \str_range:nnn {##1} { 1 } { 7 } } { hanging }
              { \dim_set:Nn \l_docxlike_format_indent_dim { - \l_docxlike_format_indent_dim } }
            \clist_map_break:
          }
      }
  }
%    \end{macrocode}
%
% 网格行距为零说明没走过 \cs{docxlike\_page\_setup:n}，那时候不存在网格，
% \option{snap-to-grid-when-document-grid-is-defined} 勾了也没有可对齐的东西，按关掉算。
%    \begin{macrocode}
\cs_new_protected:Npn \__docxlike_format_latin_line:n #1
  {
    \bool_set_false:N \l_docxlike_format_latin_line_bool
    \__docxlike_format_target_hook:n {#1}
  }
\cs_new_protected:Npn \__docxlike_format_target_hook:n #1 { }
\cs_new_protected:Npn \docxlike_format_set_target_hook:n #1
  { \cs_gset_protected:Npn \__docxlike_format_target_hook:n ##1 {#1} }
\tl_new:N \l__docxlike_format_parent_tl
\cs_new_protected:Npn \__docxlike_format_snap_axis:nnN #1#2#3
  {
    \__docxlike_format_get:nnNF {#1} {#2} #3
      {
        \prop_get:NnNTF \g__docxlike_format_parent_prop {#1} \l__docxlike_format_parent_tl
          { \__docxlike_format_snap_axis:VnN \l__docxlike_format_parent_tl {#2} #3 }
          { \tl_set:Nn #3 { false } }
      }
  }
\cs_generate_variant:Nn \__docxlike_format_snap_axis:nnN { VnN }
\cs_new_protected:Npn \docxlike_format_compute:n #1
  {
    \__docxlike_format_sz_read:nN {#1} \l__docxlike_format_value_tl
    \tl_if_empty:NTF \l__docxlike_format_value_tl
      { \dim_set:Nn \l_docxlike_format_size_dim { \f@size pt } }
      { \dim_set:Nn \l_docxlike_format_size_dim { \l__docxlike_format_value_tl } }
    \__docxlike_format_latin_line:n {#1}
    \__docxlike_format_font_metrics:n {#1}
    \bool_set:Nn \l__docxlike_format_latin_bool
      {
        \bool_if_p:N \l_docxlike_format_latin_pick_bool
        || \bool_if_p:N \l_docxlike_format_latin_line_bool
      }
    \bool_if:NTF \l__docxlike_format_latin_bool
      { \tl_set_eq:NN \l__docxlike_format_coef_tl \l_docxlike_format_natural_latin_tl }
      { \tl_set_eq:NN \l__docxlike_format_coef_tl \l__docxlike_format_strut_natural_tl }
    \dim_set:Nn \l_docxlike_format_natural_dim
      { \l__docxlike_format_coef_tl \l_docxlike_format_size_dim }
    \__docxlike_format_snap_axis:nnN {#1} { pPr/snapToGrid }
      \l__docxlike_format_value_tl
    \bool_set:Nn \l_docxlike_format_snap_vertical_bool
      {
        \str_if_eq_p:Vn \l__docxlike_format_value_tl { true }
        && \bool_if_p:N \g_docxlike_grid_bool
        && ! \dim_compare_p:nNn { \g_docxlike_page_lead_dim } = { 0pt }
      }
    \__docxlike_format_get:nnNTF {#1} { rPr/snapToGrid } \l__docxlike_format_value_tl
      { \bool_set_true:N \l_docxlike_format_hglue_set_bool }
      { \bool_set_false:N \l_docxlike_format_hglue_set_bool }
    \__docxlike_format_snap_axis:nnN {#1} { rPr/snapToGrid }
      \l__docxlike_format_value_tl
    \bool_set:Nn \l_docxlike_format_snap_horizontal_bool
      {
        \str_if_eq_p:Vn \l__docxlike_format_value_tl { true }
        && ! \dim_compare_p:nNn { \g_docxlike_page_char_dim } = { 0pt }
      }
    \__docxlike_format_read:nnnN {#1} { rPr/spacing } { auto } \l__docxlike_format_value_tl
    \str_if_eq:VnTF \l__docxlike_format_value_tl { auto }
      { \dim_zero:N \l_docxlike_format_char_spacing_dim }
      {
        \dim_set:Nn \l_docxlike_format_char_spacing_dim
          { \__docxlike_format_from_twips:V \l__docxlike_format_value_tl }
        \bool_set_true:N \l_docxlike_format_hglue_set_bool
      }
    \__docxlike_format_line_read:n {#1}
    \tl_set_eq:NN \l_docxlike_format_line_value_tl \l__docxlike_format_value_tl
    \tl_set_eq:NN \l_docxlike_format_line_unit_tl \l__docxlike_format_unit_tl
    \bool_set:Nn \l__docxlike_format_exact_bool
      { \str_if_eq_p:Vn \l__docxlike_format_unit_tl { exact } }
    \__docxlike_format_lead:
    \__docxlike_format_space_get:nnnNN {#1} { before } { before }
      \l_docxlike_format_before_dim \l_docxlike_format_before_plus_dim
    \__docxlike_format_space_get:nnnNN {#1} { after } { after }
      \l_docxlike_format_after_dim \l_docxlike_format_after_plus_dim
    \__docxlike_format_indent:n {#1}
  }
%    \end{macrocode}
%
% 一个端点折成一段基线距。倍数没有单位，光看写下来的字不知道是多长，所以端点
% 一直存着原样，到这里才折——字号、自然行高、网格行距这时候才齐。
%
%    \begin{macrocode}
\cs_new_protected:Npn \__docxlike_format_endpoint:nnN #1#2#3
  {
    \str_case:nnF {#1}
      {
        { auto } { \__docxlike_format_multiple:nN {#2} #3 }
        { atLeast }
          {
            \dim_set:Nn #3 {#2}
            \dim_set:Nn \l__docxlike_format_floor_dim
              { \bool_if:NTF \l_docxlike_format_snap_vertical_bool
                  { \g_docxlike_page_lead_dim } { \l_docxlike_format_natural_dim } }
            \dim_compare:nNnT { #3 } < { \l__docxlike_format_floor_dim }
              { \dim_set_eq:NN #3 \l__docxlike_format_floor_dim }
          }
      }
      {
        \cs_if_exist:cTF { __docxlike_format_line_unit_ #1 :nN }
          { \use:c { __docxlike_format_line_unit_ #1 :nN } {#2} #3 }
          { \dim_set:Nn #3 {#2} }
      }
  }
\cs_generate_variant:Nn \__docxlike_format_endpoint:nnN { VVN , nVN }
\cs_new_protected:Npn \__docxlike_format_multiple:nN #1#2
  {
    \bool_if:NTF \l_docxlike_format_snap_vertical_bool
      {
        \dim_set:Nn #2 { #1 \g_docxlike_page_lead_dim }
        \dim_set:Nn \l__docxlike_format_floor_dim
          {
            \fp_to_decimal:n
              {
                ceil
                  (
                    \dim_to_decimal:n { \l_docxlike_format_natural_dim }
                    / \dim_to_decimal:n { \g_docxlike_page_lead_dim }
                  )
              }
            \g_docxlike_page_lead_dim
          }
        \dim_compare:nNnT { #2 } < { \l__docxlike_format_floor_dim }
          { \dim_set_eq:NN #2 \l__docxlike_format_floor_dim }
      }
      { \dim_set:Nn #2 { #1 \l_docxlike_format_natural_dim } }
  }
%    \end{macrocode}
%
% 只写一个端点是刚性。写两个就是区间，两端各折一次，小的做自然值，大的减小的
% 做伸量，缩量恒为零：区间来自格式规范，低于小端就出了规范。
%
% 端点不分先后，谁大谁小到这里才知道：\texttt{3mm~2} 这种一头长度一头倍数的
% 写法，解析的时候没有字号，比不了。
%
% 基线落点看第一个端点的模式：\texttt{exact} 是 Word 的固定值，基线在盒顶
% 下方 0.8 倍盒高，与网格开关无关；其余三种走 \cs{\_\_docxlike\_format\_first:}。
%    \begin{macrocode}
\cs_new_protected:Npn \__docxlike_format_lead:
  {
    \__docxlike_format_endpoint:VVN \l__docxlike_format_unit_tl \l__docxlike_format_value_tl
      \l_docxlike_format_lead_dim
    \dim_zero:N \l_docxlike_format_stretch_dim
    \tl_if_empty:NF \l__docxlike_format_value_ii_tl
      {
        \__docxlike_format_endpoint:VVN \l__docxlike_format_unit_ii_tl
          \l__docxlike_format_value_ii_tl \l__docxlike_format_span_dim
        \dim_compare:nNnT
          { \l__docxlike_format_span_dim } < { \l_docxlike_format_lead_dim }
          {
            \dim_set_eq:NN \l__docxlike_format_floor_dim \l_docxlike_format_lead_dim
            \dim_set_eq:NN \l_docxlike_format_lead_dim \l__docxlike_format_span_dim
            \dim_set_eq:NN \l__docxlike_format_span_dim \l__docxlike_format_floor_dim
          }
        \dim_set:Nn \l_docxlike_format_stretch_dim
          { \l__docxlike_format_span_dim - \l_docxlike_format_lead_dim }
      }
    \str_if_eq:VnTF \l__docxlike_format_unit_tl { exact }
      {
        \dim_set:Nn \l_docxlike_format_first_dim
          { \c_docxlike_format_exact_tl \l_docxlike_format_lead_dim }
      }
      { \__docxlike_format_first: }
  }
%    \end{macrocode}
%
% 基线落点：对齐到网格时行盒居中，关掉时顶对齐。上升部与撑杆取同一个数，按
% 样式的东亚字体查（\cs{l\_docxlike\_format\_ascent\_tl}），纯西文按西文字体查。
%    \begin{macrocode}
\cs_new_protected:Npn \__docxlike_format_first:
  {
    \dim_set:Nn \l_docxlike_format_first_dim
      {
        \bool_if:NTF \l__docxlike_format_latin_bool
          { \l_docxlike_format_ascent_latin_tl } { \l_docxlike_format_ascent_tl }
        \l_docxlike_format_size_dim
      }
    \bool_if:NT \l_docxlike_format_snap_vertical_bool
      {
        \dim_add:Nn \l_docxlike_format_first_dim
          {
            0.5 \l_docxlike_format_lead_dim - 0.5 \l_docxlike_format_natural_dim
          }
      }
  }
%    \end{macrocode}
%
% \subsection{样式折成排版命令}
%
% 算出来的样式折成一串排版命令（\cs{l\_\_docxlike\_format\_out\_tl}），标题、页眉页脚、
% 题注这些本包自己排的段落用它。
%    \begin{macrocode}
\tl_new:N \l__docxlike_format_out_tl
\prop_new:N \g__docxlike_format_align_prop
\prop_gset_from_keyval:Nn \g__docxlike_format_align_prop
  { left = \raggedright , center = \centering , right = \raggedleft }
%    \end{macrocode}
%
% \emph{字体。} 样式里的字体是名字。按名字取两对系数（东亚字体的自然行高与上升部、西文字体的），
% 没写东亚字体按 SimSun、没写西文字体按文档的默认西文字体取（\cs{g\_docxlike\_format\_default\_latin\_tl}，
% 相当于 docx 里 \texttt{docDefaults} 的 \texttt{rFonts/ascii}，初值 Times New Roman，空白预设改成等线）；
% 取不到用 SimSun 与 Times New Roman 的。
% \emph{平衡单字节与双字节字符。} docx 设置里 \texttt{w:compat} 底下的
% \texttt{w:balanceSingleByteDoubleByteWidth}，Word 界面在“高级”选项的版式那一组，
% 键名 \option{balance-sbcs-characters-and-dbcs-characters}，学校的模板都开着。Windows Word 实测
% （\file{test/headings} 的 \file{heads-balance-*}，兼容模式 11）：开着时正文里贴着汉字的半角空格、
% 连着的几个半角空格各占半格（半个字加半份当时的字距），西文词间的单个空格不变（15 磅标题
% “1.1  Introduction”两个空格各 7.45，“Hello World”仍是 3.75）；空格的字形还是西文字体的，
% 不抬行高。编号格式里的空格不受它管。本包自己排的段落（\cs{docxlike\_format\_par:nn}）排之前
% 照这条改一遍内容里的空格（半个字按当时的字号 \cs{f@size} 算，不按最近一次算过的样式）：最外层的空格记号与 \verb|\ |，前后挨着空格或汉字的换成
% \cs{docxlike\_format\_sbcs\_space:}；两边都是汉字的空格记号不动，那是源码换行留下的，\pkg{xeCJK}
% 与 \pkg{luatexja} 本来就不排它（要在两个汉字之间敲空格写 \verb|\ |）。这个空格做成不伸缩的胶而不是 kern：\pkg{luatexja} 看得穿
% kern，会在前面的西文与后面的汉字之间补一段中西文间距，隔着胶就不补。正文段落的空格由
% 断行层按自己的规矩排。
%
% 调用方立了 \cs{l\_docxlike\_format\_asian\_ascent\_bool} 时，西文那一对的上升部换成东亚的
% （取大），行深不变，自然行高跟着加：西文题名配带汉字的编号就是这样的行盒（见
% \file{docxlike-heading.dtx}）。
% \cs{docxlike\_format\_font\_metrics:nnnn} 把四个数设进排版时用的那几个量，随字体命令一起
% 发出去，每个样式的撑杆就用自己字体的数。
%    \begin{macrocode}
\cs_new_protected:Npn \__docxlike_format_font_metrics:n #1
  {
    \__docxlike_format_read:nnnN {#1} { rPr/rFonts/eastAsia } { SimSun } \l__docxlike_format_font_tl
    \docxlike_font_metric:VNNF \l__docxlike_format_font_tl
      \l__docxlike_format_strut_natural_tl \l_docxlike_format_ascent_tl
      {
        \tl_set_eq:NN \l__docxlike_format_strut_natural_tl \c_docxlike_format_natural_tl
        \tl_set_eq:NN \l_docxlike_format_ascent_tl \c_docxlike_format_ascent_tl
      }
    \__docxlike_format_read:nnVN {#1} { rPr/rFonts/ascii } \g_docxlike_format_default_latin_tl
      \l__docxlike_format_font_tl
    \docxlike_font_metric:VNNF \l__docxlike_format_font_tl
      \l_docxlike_format_natural_latin_tl \l_docxlike_format_ascent_latin_tl
      {
        \tl_set_eq:NN \l_docxlike_format_natural_latin_tl \c_docxlike_format_natural_latin_tl
        \tl_set_eq:NN \l_docxlike_format_ascent_latin_tl \c_docxlike_format_ascent_latin_tl
      }
    \bool_if:NT \l_docxlike_format_asian_ascent_bool
      {
        \fp_compare:nNnT { \l_docxlike_format_ascent_tl } > { \l_docxlike_format_ascent_latin_tl }
          {
            \tl_set:Ne \l_docxlike_format_natural_latin_tl
              {
                \fp_eval:n
                  {
                    \l_docxlike_format_natural_latin_tl - \l_docxlike_format_ascent_latin_tl
                    + \l_docxlike_format_ascent_tl
                  }
              }
            \tl_set_eq:NN \l_docxlike_format_ascent_latin_tl \l_docxlike_format_ascent_tl
          }
      }
  }
\bool_new:N \l_docxlike_format_asian_ascent_bool
\bool_new:N \g_docxlike_format_balance_bool
\keys_define:nn { docxlike }
  {
    balance-sbcs-characters-and-dbcs-characters .bool_gset:N = \g_docxlike_format_balance_bool ,
    balance-sbcs-characters-and-dbcs-characters .default:n = true ,
  }
\cs_new_protected:Npn \docxlike_format_sbcs_space:
  {
    \skip_horizontal:n
      { ( \f@size pt + \l_docxlike_format_ccglue_dim ) / 2 }
  }
\seq_new:N \l__docxlike_format_balance_seq
\seq_new:N \l__docxlike_format_balance_kind_seq
\int_new:N \l__docxlike_format_balance_depth_int
\tl_new:N \l__docxlike_format_balance_tl
\tl_new:N \l__docxlike_format_balance_item_tl
\tl_const:Ne \c__docxlike_format_cspace_tl { \exp_not:N \exp_not:n { \exp_not:c { ~ } } }
\cs_new_protected:Npn \__docxlike_format_balance:N #1
  {
    \seq_clear:N \l__docxlike_format_balance_seq
    \seq_clear:N \l__docxlike_format_balance_kind_seq
    \int_zero:N \l__docxlike_format_balance_depth_int
    \tl_analysis_map_inline:Nn #1
      {
        \int_compare:nNnT { "##3 } = { 2 } { \int_decr:N \l__docxlike_format_balance_depth_int }
        \tl_set:Nn \l__docxlike_format_balance_item_tl {##1}
        \seq_put_right:Nn \l__docxlike_format_balance_seq {##1}
        \seq_put_right:Ne \l__docxlike_format_balance_kind_seq
          {
            \int_compare:nNnTF { \l__docxlike_format_balance_depth_int } = { 0 }
              {
                \int_compare:nNnTF { "##3 } = { 10 }
                  { s }
                  {
                    \tl_if_eq:NNTF \l__docxlike_format_balance_item_tl \c__docxlike_format_cspace_tl
                      { x }
                      {
                        \bool_lazy_and:nnTF
                          { \int_compare_p:nNn { "##3 } > { 10 } }
                          { \__docxlike_format_cjk_code_p:n {##2} }
                          { c } { o }
                      }
                  }
              }
              { o }
          }
        \int_compare:nNnT { "##3 } = { 1 } { \int_incr:N \l__docxlike_format_balance_depth_int }
      }
    \tl_clear:N \l__docxlike_format_balance_tl
    \int_step_inline:nn { \seq_count:N \l__docxlike_format_balance_seq }
      {
        \bool_lazy_all:nTF
          {
            { \__docxlike_format_balance_space_p:n {##1} }
            {
              \__docxlike_format_balance_space_p:n { ##1 - 1 }
              || \str_if_eq_p:ee { \__docxlike_format_balance_kind:n { ##1 - 1 } } { c }
              || \__docxlike_format_balance_space_p:n { ##1 + 1 }
              || \str_if_eq_p:ee { \__docxlike_format_balance_kind:n { ##1 + 1 } } { c }
            }
            {
              ! (
                \str_if_eq_p:ee { \__docxlike_format_balance_kind:n {##1} } { s }
                && \str_if_eq_p:ee { \__docxlike_format_balance_kind:n { ##1 - 1 } } { c }
                && \str_if_eq_p:ee { \__docxlike_format_balance_kind:n { ##1 + 1 } } { c }
              )
            }
          }
          { \tl_put_right:Nn \l__docxlike_format_balance_tl { \exp_not:N \docxlike_format_sbcs_space: } }
          {
            \tl_put_right:Ne \l__docxlike_format_balance_tl
              { \seq_item:Nn \l__docxlike_format_balance_seq {##1} }
          }
      }
    \tl_set:Ne #1 { \l__docxlike_format_balance_tl }
  }
\prg_new_conditional:Npnn \__docxlike_format_balance_space:n #1 { p }
  {
    \bool_lazy_or:nnTF
      { \str_if_eq_p:ee { \__docxlike_format_balance_kind:n {#1} } { s } }
      { \str_if_eq_p:ee { \__docxlike_format_balance_kind:n {#1} } { x } }
      { \prg_return_true: } { \prg_return_false: }
  }
\cs_new:Npn \__docxlike_format_balance_kind:n #1
  {
    \int_compare:nNnTF {#1} < { 1 } { o }
      {
        \int_compare:nNnTF {#1} > { \seq_count:N \l__docxlike_format_balance_kind_seq } { o }
          { \seq_item:Nn \l__docxlike_format_balance_kind_seq { \int_eval:n {#1} } }
      }
  }
\prg_new_conditional:Npnn \__docxlike_format_cjk_code:n #1 { p }
  {
    \bool_lazy_or:nnTF
      { \int_compare_p:nNn {#1} > { "2E7F } && \int_compare_p:nNn {#1} < { "A000 } }
      {
        ( \int_compare_p:nNn {#1} > { "F8FF } && \int_compare_p:nNn {#1} < { "FB00 } )
        || ( \int_compare_p:nNn {#1} > { "FEFF } && \int_compare_p:nNn {#1} < { "FFF0 } )
        || \int_compare_p:nNn {#1} > { "1FFFF }
      }
      { \prg_return_true: } { \prg_return_false: }
  }
\tl_new:N \g_docxlike_format_default_latin_tl
\tl_gset:Nn \g_docxlike_format_default_latin_tl { Times~New~Roman }
\cs_generate_variant:Nn \__docxlike_format_read:nnnN { nnV }
\prg_generate_conditional_variant:Nnn \docxlike_font_metric:nNN { V } { F }
\cs_new_protected:Npn \docxlike_format_font_metrics:nnnn #1#2#3#4
  {
    \tl_set:Nn \l__docxlike_format_strut_natural_tl {#1}
    \tl_set:Nn \l_docxlike_format_ascent_tl {#2}
    \tl_set:Nn \l_docxlike_format_natural_latin_tl {#3}
    \tl_set:Nn \l_docxlike_format_ascent_latin_tl {#4}
  }
%    \end{macrocode}
%
% \emph{按名字切换字体。} 文档类可以说“这个名字在本文档里用什么命令切过去”：
% \cs{docxlike\_font\_select\_new:nn}\marg{字体名}\marg{代码}，比如字体集里没有 SimSun、要用
% 替身时。预设表认得的字体，同一副的别的名字（中文名、全名）一起登记。没说过的，东亚字体按名字声明一个中文字族（\hologo{XeLaTeX} 用 \pkg{xeCJK}
% 的字族，\hologo{LuaLaTeX} 用 \pkg{luatexja-fontspec} 的 \cs{newjfontfamily}，JFM 取 \texttt{zh\_CN/quanjiao}），
% 西文字体按名字声明一个 \pkg{fontspec} 字族，都只声明一次；本机没有这副字体时警告一次，不切。
% \cs{newjfontfamily} 定义的切换命令（连同它里面那一个 \texttt{\meta{命令}~code}）是局部的，
% 字体常在段落的组里第一次用到，出组就没了，所以声明之后改成全局的。
%    \begin{macrocode}
\prop_new:N \g__docxlike_font_select_prop
\tl_new:N \l__docxlike_font_code_tl
\tl_new:N \l__docxlike_font_group_tl
\cs_new_protected:Npn \docxlike_font_select_new:nn #1#2
  {
    \__docxlike_font_key:n {#1}
    \prop_get:NVNTF \g__docxlike_font_group_prop \l__docxlike_font_key_str
      \l__docxlike_font_group_tl
      {
        \clist_map_inline:Nn \l__docxlike_font_group_tl
          {
            \__docxlike_font_key:n {##1}
            \prop_gput:NVn \g__docxlike_font_select_prop \l__docxlike_font_key_str {#2}
          }
      }
      { \prop_gput:NVn \g__docxlike_font_select_prop \l__docxlike_font_key_str {#2} }
  }
\msg_new:nnnn { docxlike } { font-not-found }
  { Font~'#1'~is~not~installed;~the~current~font~is~kept. }
  { Install~it,~or~tell~docxlike~what~to~use~with~\token_to_str:N \docxlike_font_select_new:nn. }
\cs_new_protected:Npn \docxlike_font_select_east_asian:n #1
  {
    \__docxlike_font_key:n {#1}
    \prop_get:NVNTF \g__docxlike_font_select_prop \l__docxlike_font_key_str
      \l__docxlike_font_code_tl
      { \l__docxlike_font_code_tl }
      { \__docxlike_font_east_asian:n {#1} }
  }
\cs_new_protected:Npn \docxlike_font_select_latin:n #1
  {
    \__docxlike_font_key:n {#1}
    \prop_get:NVNTF \g__docxlike_font_select_prop \l__docxlike_font_key_str
      \l__docxlike_font_code_tl
      { \l__docxlike_font_code_tl }
      { \__docxlike_font_latin:n {#1} }
  }
\cs_new_protected:Npn \__docxlike_font_east_asian:n #1
  {
    \bool_lazy_or:nnT
      { \cs_if_exist_p:N \xeCJK_set_family:nnn }
      { \cs_if_exist_p:N \newjfontfamily }
      {
        \tl_if_exist:cF { g__docxlike_font_ea_ \l__docxlike_font_key_str _tl }
          {
            \tl_new:c { g__docxlike_font_ea_ \l__docxlike_font_key_str _tl }
            \__docxlike_font_if_exist:nTF {#1}
              {
                \cs_if_exist:NTF \xeCJK_set_family:nnn
                  {
                    \xeCJK_set_family:enn { docxlike-\l__docxlike_font_key_str } { } {#1}
                    \tl_gset:ce { g__docxlike_font_ea_ \l__docxlike_font_key_str _tl }
                      { \exp_not:N \CJKfamily { docxlike-\l__docxlike_font_key_str } }
                  }
                  {
                    \exp_args:Nc \newjfontfamily { docxlike@jf@ \l__docxlike_font_key_str }
                      [ JFM = zh_CN/quanjiao ] {#1}
                    \cs_gset_eq:cc { docxlike@jf@ \l__docxlike_font_key_str }
                      { docxlike@jf@ \l__docxlike_font_key_str }
                    \cs_if_exist:cT { docxlike@jf@ \l__docxlike_font_key_str \c_space_tl code }
                      {
                        \cs_gset_eq:cc { docxlike@jf@ \l__docxlike_font_key_str \c_space_tl code }
                          { docxlike@jf@ \l__docxlike_font_key_str \c_space_tl code }
                      }
                    \tl_gset:ce { g__docxlike_font_ea_ \l__docxlike_font_key_str _tl }
                      { \exp_not:c { docxlike@jf@ \l__docxlike_font_key_str } }
                  }
              }
              { \msg_warning:nnn { docxlike } { font-not-found } {#1} }
          }
        \tl_use:c { g__docxlike_font_ea_ \l__docxlike_font_key_str _tl }
      }
  }
\cs_new_protected:Npn \__docxlike_font_latin:n #1
  {
    \cs_if_exist:NT \fontspec_set_family:Nnn
      {
        \tl_if_exist:cF { g__docxlike_font_latin_ \l__docxlike_font_key_str _tl }
          {
            \tl_new:c { g__docxlike_font_latin_ \l__docxlike_font_key_str _tl }
            \__docxlike_font_if_exist:nTF {#1}
              {
                \fontspec_set_family:cnn
                  { g__docxlike_font_latin_ \l__docxlike_font_key_str _tl } { } {#1}
                \tl_gset_eq:cc { g__docxlike_font_latin_ \l__docxlike_font_key_str _tl }
                  { g__docxlike_font_latin_ \l__docxlike_font_key_str _tl }
              }
              { \msg_warning:nnn { docxlike } { font-not-found } {#1} }
          }
        \tl_if_empty:cF { g__docxlike_font_latin_ \l__docxlike_font_key_str _tl }
          {
            \fontfamily { \tl_use:c { g__docxlike_font_latin_ \l__docxlike_font_key_str _tl } }
            \selectfont
          }
      }
  }
\prg_new_protected_conditional:Npnn \__docxlike_font_if_exist:n #1 { TF }
  {
    \cs_if_exist:NTF \fontspec_font_if_exist:nTF
      { \fontspec_font_if_exist:nTF {#1} { \prg_return_true: } { \prg_return_false: } }
      { \prg_return_false: }
  }
\cs_generate_variant:Nn \fontspec_set_family:Nnn { cnn }
\cs_generate_variant:Nn \tl_gset:Nn { ce }
%    \end{macrocode}
%
% \cs{docxlike\_format\_strut:} 是一根零宽的撑杆，把它所在那一行的行盒撑成 Word 的尺寸。
%
% Word 的行盒高度由字体的上升部与下降部定，与这一行实际写了什么字无关；
% \hologo{TeX} 的行高是这一行各个字形外框的最大值。汉字的字形外框比字体的上升部
% 矮一截，小二黑体差 3.4\,bp，小四宋体差 2.3\,bp。段内看不出来——段内的行距由
% \cs{baselineskip} 定死——差别全落在段与段的接缝上：
% \hologo{TeX} 那里两段之间是“上段末行的深度加间距加下段首行的高度”，
% 三项里有两项是字形量出来的，比 Word 的窄。页首那一行也一样，
% \cs{topskip} 拦不住比它高的行。
%
% 撑杆按当时的字号与 \cs{baselineskip} 现算，一根走遍正文、标题、页眉：
% \begin{itemize}
% \item 自然行高 $H$ 与上升部 $A$ 用换算一节的常数，汉字与纯西文各一副
%   （\cs{docxlike\_format\_latin\_true:} 切换）；
% \item 行盒高 $L$ 取 \cs{baselineskip} 与 $H$ 的较大者。取大这一步是为页眉设的：
%   页眉那一行的 \cs{baselineskip} 是零，光看它撑出来的盒子没有深度，
%   页眉横线会整块上移 2\,bp；
% \item 基线在盒里的落点，不对齐网格时是 $A$（顶对齐），对齐时是
%   $(L-H)/2+A$（居中）。这一支由 \cs{l\_docxlike\_format\_line\_snap\_bool} 决定，
%   正文的取值由 \cs{docxlike\_format\_apply\_body:} 全局设好，标题在自己的
%   \option{styles} 里临时改，出了标题那个组就复原。
% \end{itemize}
%
% 改这个标志位的是两个不带参数的命令 \cs{docxlike\_format\_snap\_true:} 与
% \cs{docxlike\_format\_snap\_false:}，不是一个收真假值的命令。收值那种写法要在标题
% 排版的当口解析一个布尔表达式，而那时候 \pkg{expl3} 的 catcode 已经关了，
% 解析器读不出 \cs{\_\_bool\_f\_0:} 这类内部名字，报 \texttt{Missing number}。
%
% 字号要先落进一个长度再乘系数。\cs{f@size} 是个数字串不是长度寄存器，
% \texttt{1.0156\cs{f@size}pt} 在 \cs{dimexpr} 里读成两个数并排，报
% \texttt{Illegal unit of measure}。
%    \begin{macrocode}
\bool_new:N \l_docxlike_format_line_snap_bool
\dim_new:N \l__docxlike_format_strut_dim
\dim_new:N \l__docxlike_format_nat_dim
\dim_new:N \l__docxlike_format_box_dim
\dim_new:N \l__docxlike_format_asc_dim
\cs_new_protected:Npn \docxlike_format_snap_true:
  { \bool_set_true:N \l_docxlike_format_line_snap_bool }
\cs_new_protected:Npn \docxlike_format_snap_false:
  { \bool_set_false:N \l_docxlike_format_line_snap_bool }
%    \end{macrocode}
%
% 横轴的两道胶按元素发，不只在正文设一次：Word 那边三级标题写着
% \texttt{snapToGrid=0}，不吃正文的字距。发的位置是 \cs{\_\_docxlike\_format\_build:n}
% 拼的那串排版命令，只管那一段。
%
% \emph{伸量给多少。} \pkg{xeCJK} 原样给 0.08\cs{baselineskip}，一行 33 缝合计
% 约 51\,bp，两端对齐时行就靠拉汉字间距填满。Word 不这样：网格是刚的，填行靠压
% 本行的标点，汉字之间纹丝不动。样本文档里那种整行 34 字自然宽 423.01\,bp、
% 行尾离版心还差 2.2\,bp，就是没拉伸的样子。
%
% 伸量太大会少装字：\hologo{TeX} 断行取全局最优，拉伸比压缩便宜得多，于是宁可
% 少放一个字去把行拉开。伸量为零又会走到另一头：\pkg{xeCJK} 自带的标点压缩被
% 全逼出来，实测有标点被压到 5.98\,bp 以下，比 Word 的下限半格还窄，行里因此
% 多塞进本不该有的字。
%
% 取 0.02：一行能吸收一整格（12.4543 除以 33 缝，每缝约 0.38\,bp）。版心
% 425.2\,bp 除以格宽是 34.14 格，全刚性时任何整数格的行都凑不满，前两遍断行全败、
% 掉进应急遍被 \cs{emergencystretch} 硬撑；扫过 0.01／0.02／0.03，只有 0.02 与样本文档
% 逐行一致。
%
% 中西文胶与空格的伸量也归到这一档（见 \cs{docxlike\_format\_latin\_pitch:} 那边的
% 说明）：伸量在 \hologo{TeX} 里兼任摊派权重，三种缝同一档，行凑不满时缺口
% 才会像 Word 那样均摊，不在某一个空格上开大洞。
%
% \hologo{XeLaTeX} 下汉字之间的胶是 \pkg{xeCJK} 的 \cs{CJKglue}，中西文之间的是 \cs{CJKecglue}。
%
% \hologo{LuaLaTeX} 下这两段胶不叫 \cs{CJKglue}：\pkg{luatexja} 做成
% \texttt{kanjiskip} 与 \texttt{xkanjiskip} 两个参数，走 \cs{ltjsetparameter}，
% 名字不同管的是同一件事。同一组值原样喂过去，先过一趟寄存器——
% \cs{ltjsetparameter} 不吃 \hologo{LaTeX} 的胶写法。\hologo{XeLaTeX} 下
% \cs{ltjsetparameter} 不存在，这一段整个跳过。
%
% 全角空格 U+3000 在 Word 里是空格，不是汉字，旁边不加中西文间距（编号
% “1.1　标题”里数字与全角空格之间没有那一截）；\pkg{luatexja} 把它当汉字，
% 这里用 \texttt{jaxspmode} 把它两侧的 \texttt{xkanjiskip} 关掉。
%
%    \begin{macrocode}
\tl_const:Nn \c__docxlike_format_ccstretch_tl { 0.02 }
\tl_const:Nn \c__docxlike_format_ecstretch_tl { 0.02 }
\tl_const:Nn \c__docxlike_format_ecshrink_tl { 0.13 }
%    \end{macrocode}
%
% \emph{中西文胶只该发给字母和数字。} 1--5 页版实测 \texttt{即→×} 是 12.32、
% \texttt{“→,} 是 12.35、\texttt{汉→,} 约一格——Word 的中西文空隙只加在汉字与
% 字母数字之间，西文标点、乘除号这类符号旁没有。按字符豁免不走窥探（试过：
% \cs{CJKecglue} 的调用点后面时常跟着 \cs{scan\_stop:} 或 \pkg{xeCJK} 的内部宏，
% \cs{peek\_after:Nw} 读不到真字符），走字符类：这些字符声明进 \pkg{xeCJK}
% 本来就为半角标点准备的 \texttt{HalfLeft}/\texttt{HalfRight} 类，与
% \texttt{Default} 拆开之后，边界规则在 \cs{docxlike\_format\_punct\_grid:} 里按类各配。
%    \begin{macrocode}
\dim_new:N \l__docxlike_format_ecglue_dim
\skip_new:N \l__docxlike_format_kanji_skip
\cs_new_protected:Npn \__docxlike_format_ecglue_emit:
  {
    \hskip \l__docxlike_format_ecglue_dim
      \@plus \c__docxlike_format_ecstretch_tl \baselineskip
      \@minus \c__docxlike_format_ecshrink_tl \baselineskip
  }
\dim_new:N \l_docxlike_format_ccglue_dim
\cs_new_protected:Npn \docxlike_format_hglue:nn #1#2
  {
    \dim_set:Nn \l_docxlike_format_ccglue_dim { \dim_eval:n {#1} }
    \cs_set:Npx \CJKglue
      {
        \exp_not:N \hskip \dim_eval:n {#1}
        \exp_not:N \@plus \c__docxlike_format_ccstretch_tl \exp_not:N \baselineskip
      }
    \dim_set:Nn \l__docxlike_format_ecglue_dim { \dim_eval:n {#2} }
    \cs_set_protected:Npn \CJKecglue { \__docxlike_format_ecglue_guard: }
    \cs_if_exist:NT \ltjsetparameter
      {
        \skip_set:Nn \l__docxlike_format_kanji_skip
          {
            \dim_eval:n {#1}
            plus \c__docxlike_format_ccstretch_tl \baselineskip
          }
        \ltjsetparameter { kanjiskip = \l__docxlike_format_kanji_skip }
        \skip_set:Nn \l__docxlike_format_kanji_skip
          {
            \dim_eval:n {#2}
            plus \c__docxlike_format_ecstretch_tl \baselineskip
            minus \c__docxlike_format_ecshrink_tl \baselineskip
          }
        \ltjsetparameter { xkanjiskip = \l__docxlike_format_kanji_skip }
        \ltjsetparameter { jaxspmode = { "3000 , inhibit } }
      }
%    \end{macrocode}
%
% \emph{还得把 \pkg{xeCJK} 的中西胶寄存器一并设了。} \cs{CJKecglue} 我们赋成
% 一个发胶的宏，\pkg{xeCJK} 的 \cs{xeCJK\_glue\_to\_skip:nN} 只在自己改
% \cs{CJKecglue} 时才把它折成 \cs{l\_\_xeCJK\_ecglue\_skip}，我们绕过那条路，
% 寄存器就一直停在包的默认值（汉字字体的空格，12\,bp 下 3.912）。
% \TeX~Live 2025 的边界几乎都直接调 \cs{CJKecglue}，看不出来；2026 重写边界之后
% \cs{\_\_xeCJK\_replace\_space:}、\cs{\_\_xeCJK\_use\_ecglue\_skip:} 这些
% 改读寄存器，陈旧值就露了出来，源码里带空格的中西文混排两版对不上。
%    \begin{macrocode}
    \skip_if_exist:NT \l__xeCJK_ecglue_skip
      {
        \skip_set:Nn \l__xeCJK_ecglue_skip
          {
            \l__docxlike_format_ecglue_dim
            plus \c__docxlike_format_ecstretch_tl \baselineskip
          }
      }
%    \end{macrocode}
%
% 设完之后用钩子 \texttt{docxlike/hglue} 通知文档类，字距在 \cs{l\_docxlike\_format\_ccglue\_dim}：
% 文档类另有自己一套字距（比如别的中文宏包记着的那一份）要跟着同步，就挂在这里。
%    \begin{macrocode}
    \hook_use:n { docxlike / hglue }
  }
\cs_new_protected:Npn \__docxlike_format_strut:nn #1#2
  {
    \leavevmode
    \__docxlike_format_strut_dims:nn {#1} {#2}
    \vrule ~ width ~ \c_zero_dim
      ~ height ~ \l__docxlike_format_asc_dim
      ~ depth ~ \dim_eval:n
          { \l__docxlike_format_box_dim - \l__docxlike_format_asc_dim }
    \scan_stop:
  }
\cs_new_protected:Npn \__docxlike_format_strut_dims:nn #1#2
  {
    \dim_set:Nn \l__docxlike_format_strut_dim { \f@size pt }
    \dim_set:Nn \l__docxlike_format_nat_dim
      { #1 \l__docxlike_format_strut_dim }
    \bool_if:NTF \l_docxlike_format_strut_exact_bool
      {
        \dim_set_eq:NN \l__docxlike_format_box_dim \baselineskip
        \dim_set:Nn \l__docxlike_format_asc_dim
          { \c_docxlike_format_exact_tl \l__docxlike_format_box_dim }
      }
      {
        \dim_set:Nn \l__docxlike_format_box_dim
          { \dim_max:nn { \baselineskip } { \l__docxlike_format_nat_dim } }
        \dim_set:Nn \l__docxlike_format_asc_dim
          { #2 \l__docxlike_format_strut_dim }
        \bool_if:NT \l_docxlike_format_line_snap_bool
          {
            \dim_add:Nn \l__docxlike_format_asc_dim
              { 0.5 \l__docxlike_format_box_dim - 0.5 \l__docxlike_format_nat_dim }
          }
      }
  }
%    \end{macrocode}
%
% \emph{固定值行距的基线。} Word 的 \texttt{w:lineRule="exact"} 下，行盒高被钉死
% 成给定值，基线落在盒顶下 \cs{c\_docxlike\_format\_exact\_tl} 倍处，
% \emph{与字号无关、与中西文无关}。
%
% 这条是量出来的：自造探针 \file{docx}，每页只放一个测试段——这样盒顶就是版心顶，
% 基线减版心顶即上升部，一次量到绝对值，不必解方程。Word 亲自导出 PDF，用文本
% 矩阵取基线。固定值取 18 到 150\,pt 六档，字号四档，中西文各一组，
% 最小二乘得 \emph{上升部 = 0.79997 × 固定值 + 0.047}，截距在 Word 自身
% 0.25\,bp 的位置量化噪声里。每个固定值下四组字号与文种给出的上升部完全相同，
% 这一条否掉了“行距减下降部”那种随字号变的猜想。
%
% \cs{\_\_docxlike\_format\_first:} 早就按这条算了，只是它只管页首那一行的
% \cs{topskip}；撑杆管每一行，缺了这条分支，固定值行距的段落每一行都偏。
%    \begin{macrocode}
\bool_new:N \l_docxlike_format_strut_exact_bool
\cs_new_protected:Npn \docxlike_format_exact_true:
  { \bool_set_true:N \l_docxlike_format_strut_exact_bool }
\cs_new_protected:Npn \docxlike_format_exact_false:
  { \bool_set_false:N \l_docxlike_format_strut_exact_bool }
\bool_new:N \l_docxlike_format_strut_latin_bool
\cs_new_protected:Npn \docxlike_format_latin_true:
  { \bool_set_true:N \l_docxlike_format_strut_latin_bool }
\cs_new_protected:Npn \docxlike_format_latin_false:
  { \bool_set_false:N \l_docxlike_format_strut_latin_bool }
%    \end{macrocode}
%
% \textbf{公式独占一段时上下各多一行空白。} 写成
% \begin{latex}
% 上一段文字。
%
% \begin{equation}...\end{equation}
%
% 下一段文字。
% \end{latex}
% 排出来公式上下各多一整行。两头的来路不同：
%
% \emph{上面那一行是 \hologo{LaTeX} 的老行为。} \cs{begin}\marg{equation} 在竖排
% 模式下先起一个段落，首行缩进盒让这一段非空，\hologo{TeX} 于是排出一行只有缩进
% 的空行。\cls{article} 加 \pkg{amsmath} 实测同样如此（行距 14.45 的文档里，
% 公式紧接文字是 14.45，空一行是 28.89），不是本包引入的。
%
% \emph{下面那一行是撑杆造的。} \cs{end}\marg{equation} 之后段落还开着，紧跟的
% 空行让这个空续段收尾。\hologo{TeX} 对空续段本不产生任何行，而 \texttt{para/end}
% 的撑杆把它撑成一行只有零宽 rule 的幽灵行。不装撑杆时与 \cls{ctexbook} 都没有这一行。
%
% 两头一起收在这两个钩子里：
%
% \cs{\_\_docxlike\_format\_para\_close:} 只在段落片段\emph{有内容}时发撑杆。空不空看
% \cs{lastnodetype}：横排列表为空时它是 $-1$。
%
% \cs{\_\_docxlike\_format\_para\_open:} 看一个标志。标志由行间公式的环境在竖排模式下
% 立起来（见 \cs{docxlike\_format\_display\_enter:}），那一段的撑杆就不发，缩进盒也由
% \cs{noindent} 免掉，两个成因一起去掉，空行才真的没有——只 \cs{noindent} 不够，
% 撑杆照样把那一行撑出来（实测）。
%
% 还有一个标志管\emph{两头}：\cs{docxlike\_format\_blank\_para:} 立起来之后，下一段段首
% 段尾都不发撑杆，段尾发完就放倒。给 \pkg{tabularray} 的长表用：它起表前先
% \cs{leavevmode} 开一个空段去刷新 \cs{pagetotal}，排每一页又用 \cs{noindent}
% 开一段装整张表的盒子。这两段不是内容，撑杆一发，前者成了一整行空白，与
% \texttt{presep} 叠出两行，后者给表盒底下垫出一个行深。段首那个标志管不住
% 段尾：空段里有缩进盒，\cs{lastnodetype} 不是 $-1$。
%
% 这样的段还得交回 \hologo{TeX} 断行。Lua 那头给没有一个字的行按段落字体算行盒
% （行里只有公式或图时 Word 就是这样：图有多高行就多高，行深仍是文字的），
% 装表的那一段就会在表盒底下垫出一个行深来，与撑杆垫的是同一回事。所以立标志时
% 顺手把 \cs{docxlike@linebreak@attr} 清零（段首节点上为 0 的段 Lua 不碰），段一开就
% 设回去，后面装表题、装格子的段照旧走 Lua。设回去的是清零前的值，不是 1：
% 不让 Lua 碰的环境里本来就是 0。已经立着时不再存，不然存下的是 0。段尾也设一次：
% 续页那一段是 \pkg{tabularray} 自己先用 \cs{noindent} 开的，装表时段早开着，
% 段首那一句轮不到。
%
% 页眉页脚里两个钩子什么也不做：文档类排页眉页脚前立起
% \cs{l\_docxlike\_format\_in\_header\_bool}，钩子见了就整个跳过，连行内代码收尾的
% 标志也不放倒。页眉在输出例程里排，那里的段不该改动正文排到一半的状态。
%    \begin{macrocode}
\bool_new:N \g__docxlike_format_display_para_bool
\bool_new:N \g__docxlike_format_blank_para_bool
\bool_new:N \l_docxlike_format_in_header_bool
\int_new:N \g__docxlike_format_blank_attr_int
\cs_new_protected:Npn \docxlike_format_blank_para:
  {
    \bool_if:NF \g__docxlike_format_blank_para_bool
      {
        \cs_if_exist:NT \docxlike@linebreak@attr
          {
            \int_gset:Nn \g__docxlike_format_blank_attr_int { \docxlike@linebreak@attr }
            \docxlike@linebreak@attr = 0 \scan_stop:
          }
      }
    \bool_gset_true:N \g__docxlike_format_blank_para_bool
  }
\cs_new_protected:Npn \__docxlike_format_blank_para_restore:
  {
    \cs_if_exist:NT \docxlike@linebreak@attr
      { \docxlike@linebreak@attr = \g__docxlike_format_blank_attr_int \scan_stop: }
  }
\cs_new_protected:Npn \__docxlike_format_para_open:
  {
    \bool_if:NF \l_docxlike_format_in_header_bool
      {
        \bool_gset_false:N \g_docxlike_format_verb_exit_bool
        \bool_if:NTF \g__docxlike_format_blank_para_bool
          { \__docxlike_format_blank_para_restore: }
          {
            \bool_if:NTF \g__docxlike_format_display_para_bool
              { \bool_gset_false:N \g__docxlike_format_display_para_bool }
              { \docxlike_format_strut: }
          }
      }
  }
\cs_new_protected:Npn \__docxlike_format_para_close:
  {
    \bool_if:NF \l_docxlike_format_in_header_bool
      {
        \bool_if:NTF \g__docxlike_format_blank_para_bool
          {
            \bool_gset_false:N \g__docxlike_format_blank_para_bool
            \__docxlike_format_blank_para_restore:
          }
          {
            \int_compare:nNnF { \tex_lastnodetype:D } = { -1 }
              { \docxlike_format_strut: }
          }
      }
  }
%    \end{macrocode}
%
% \begin{macro}{\docxlike_format_display_enter:}
% 行间公式起头时如果还在竖排模式，这一段就是“只装一条公式”的段落：立标志、
% 免缩进。挂在 \pkg{amsmath} 那一族环境上，\cs{[} 与 \env{displaymath} 都归到
% \env{equation*}（实测 \cs{meaning} 是 \cs{begin}\marg{equation*}），所以那一条
% 就够。
%    \begin{macrocode}
\cs_new_protected:Npn \docxlike_format_display_enter:
  {
    \mode_if_vertical:T
      {
        \bool_gset_true:N \g__docxlike_format_display_para_bool
        \tex_noindent:D
      }
  }
\clist_const:Nn \c__docxlike_format_display_clist
  {
    equation , equation* , displaymath , eqnarray , eqnarray* ,
    align , align* , alignat , alignat* , flalign , flalign* ,
    gather , gather* , multline , multline*
  }
\cs_new_protected:Npn \docxlike_format_display_hooks:
  {
    \clist_map_inline:Nn \c__docxlike_format_display_clist
      { \AddToHook { env/##1/begin } { \docxlike_format_display_enter: } }
  }
%    \end{macrocode}
% \end{macro}
%
% 汉字那一支的自然行高与上升部都取最近一次 \cs{docxlike\_format\_compute:n} 按
% \texttt{eastAsia} 查到的那一对。自然行高先前写死成中易宋体的 1.296875，行盒高
% 不过网格的地方看不出来，换成等线（1.351334）就差出来了。
%    \begin{macrocode}
\cs_new_protected:Npn \docxlike_format_strut:
  {
    \__docxlike_format_parend_guard:
    \bool_if:NTF \l_docxlike_format_strut_latin_bool
      {
        \__docxlike_format_strut:nn
          { \l_docxlike_format_natural_latin_tl }
          { \l_docxlike_format_ascent_latin_tl }
      }
      {
        \__docxlike_format_strut:nn
          { \l__docxlike_format_strut_natural_tl } { \l_docxlike_format_ascent_tl }
      }
  }
%    \end{macrocode}
%
% \emph{块与块之间的基线距。} Word 的行距是每一段自己的事，段与段之间不另有
% 一个量：下一段首行的基线落在“上一段末行的行深”加“本段首行的上升部”处。
% 段内每一行等距，所以这条只在一行一段的地方看得出来，目录的每一条就是一段。
%
% 上升部与行盒高由\emph{这一行内容的字体}定，不由段落样式里写的那副字定：
% 一条纯西文的条目（\texttt{Abstract}）行盒比汉字的矮，Word 排出来的步进跟着变。
% 样本文档的目录四种组合各是一个数（小四 12\,bp、多倍 1.25）：
% \begin{longtblr}[caption={样本文档目录条目的步进}, label={tab:impl-toc-pitch}]
%   {colspec={l l l l}}
% \toprule
% 上一条 & 本条 & 行深 + 上升部 & 步进 \\ \midrule
% 汉 & 汉 & 7.275 + 12.178 & 19.453 \\
% 汉 & 西 & 7.275 + 11.273 & 18.548 \\
% 西 & 汉 & 5.977 + 12.178 & 18.155 \\
% 西 & 西 & 5.977 + 11.273 & 17.250 \\
% \bottomrule
% \end{longtblr}
% 这四个数不存在表里。上升部是 \cs{l\_docxlike\_format\_first\_dim}，行深是行盒高
% 减上升部，两个都由 \cs{docxlike\_format\_compute:n} 按目标自己的字号与行距算出来。
% 1.2 倍、固定值 23\,bp、1.5 倍这些行距因此走同一条式子，各档不必自备一组数。
%
% \cs{docxlike\_format\_stack\_seed:n}\marg{目标} 起头，只存行深不设行距：头一块的
% 位置由它上面那一整块定。\cs{docxlike\_format\_stack\_line:n}\marg{目标} 每一块调一次。
% 这一块是不是纯西文，调用方排之前写进 \cs{l\_docxlike\_format\_latin\_pick\_bool}。
%
% \emph{一块折成几行的时候。} \hologo{TeX} 的 \cs{baselineskip} 管整段，
% 而 Word 的式子里只有\emph{头一行}用上一块的行深，段内后面几行用本块自己的。
% 目录里一条长条目折两行就看得出来：先前第二行也吃“上一块行深加上升部”，
% 从目录标题那一块接下来是 23.13，Word 是 19.45。
%
% 所以 \cs{baselineskip} 设成本块自己的“行深加上升部”，两者的差额
% （上一块行深减本块行深）单独发一段竖胶。不折行时两种写法的总距离一模一样
% （胶加行距等于原来那个和），折行时后面几行才对。\cs{lineskiplimit} 压到底
% 之后行间胶恒按 \cs{baselineskip} 算，\cs{prevdepth} 不参与，这一步才成立。
%
% \emph{撑杆的中西文也跟着这一块换。} 行盒的上升部由段落钩子发的撑杆定
% （\cs{docxlike\_format\_strut:} 看 \cs{l\_docxlike\_format\_strut\_latin\_bool}），而
% \cs{docxlike\_format\_compute:n} 只算取值不发这个标志位，先前一条纯西文的条目
% 因此拿着西文的行盒高、汉字的上升部：步进（行深加上升部）照样对，头一条的
% 基线却低 0.905——正是两副上升部之差。英文目录里整页跟着低，实测 26 行。
%    \begin{macrocode}
\dim_new:N \g__docxlike_format_depth_dim
\cs_new_protected:Npn \__docxlike_format_stack_depth:
  {
    \dim_gset:Nn \g__docxlike_format_depth_dim
      { \l_docxlike_format_lead_dim - \l_docxlike_format_first_dim }
  }
\cs_new_protected:Npn \docxlike_format_stack_seed:n #1
  {
    \docxlike_format_compute:n {#1}
    \__docxlike_format_stack_depth:
  }
\dim_new:N \l__docxlike_format_stack_own_dim
\cs_new_protected:Npn \docxlike_format_stack_line:n #1
  {
    \docxlike_format_compute:n {#1}
    \bool_set_eq:NN \l_docxlike_format_strut_latin_bool \l__docxlike_format_latin_bool
    \dim_set:Nn \l__docxlike_format_stack_own_dim
      { \l_docxlike_format_lead_dim - \l_docxlike_format_first_dim }
    \dim_set:Nn \baselineskip
      { \l__docxlike_format_stack_own_dim + \l_docxlike_format_first_dim }
    \skip_vertical:n
      { \g__docxlike_format_depth_dim - \l__docxlike_format_stack_own_dim }
    \dim_gset_eq:NN \g__docxlike_format_depth_dim \l__docxlike_format_stack_own_dim
  }
%    \end{macrocode}
%
% 一块内容里有没有汉字。串成字符串逐字符看码位，过了 \texttt{U+2E7F}
% （康熙部首之前那一格）就算；宏名全是 ASCII，不会误判。
%    \begin{macrocode}
\bool_new:N \l__docxlike_format_cjk_bool
\prg_new_protected_conditional:Npnn \docxlike_format_cjk:n #1 { TF , T , F }
  {
    \bool_set_false:N \l__docxlike_format_cjk_bool
    \str_map_inline:nn {#1}
      {
        \int_compare:nNnT { `##1 } > { "2E7F }
          { \bool_set_true:N \l__docxlike_format_cjk_bool \str_map_break: }
      }
    \bool_if:NTF \l__docxlike_format_cjk_bool \prg_return_true: \prg_return_false:
  }
%    \end{macrocode}
%
%
% 样式算出来的格式是整串替换，不叠在别处原有的格式上：原有的写了什么（比如
% \cs{bfseries}）都甩掉，\option{styles} 里没写 \option{bold} 就不粗。字体、字号、行距、
% 对齐、字重、段前段后全从 \option{styles} 算，一处说了算。
%
% 有两样手写的标题格式里常有、这边有意不要的。
%
% \emph{一是行距的伸缩量。} 手写的标题格式常把行距写成
% \texttt{21bp~plus~1.677267bp~minus~1.157391bp} 这种可伸缩的值。
% Word 的行距是刚性的，本包照 Word 建模，所以这边只发 \cs{fontsize}\marg{字号}\marg{行距}，没有伸缩量，
% \option{styles} 族也不打算加表达伸缩的键。这是有意的，不是漏了。
%
% \emph{二是各级标题的 \cs{bfseries}。} 旧格式里章、节、条常带着它，
% 这边一律不加粗：黑体本身够重，规范也没要求。要加粗就在
% \option{styles} 里写 \option{bold}\texttt{=true}，一处说了算。
%
% 文档类可以给某个目标的字体设置追加一段代码：
% \cs{docxlike\_format\_add\_font:nn}\marg{目标}\marg{代码}。比如让某几级标题里的
% 西文走无衬线，就给那几级登记一段发 \cs{sffamily} 的代码。只追加到登记过的目标上，
% 页眉页脚也走 \cs{\_\_docxlike\_format\_build:n}，无条件发会串味。
%    \begin{macrocode}
\prop_new:N \g__docxlike_format_font_extra_prop
\tl_new:N \l__docxlike_format_extra_tl
\cs_new_protected:Npn \docxlike_format_add_font:nn #1#2
  { \prop_gput:Nnn \g__docxlike_format_font_extra_prop {#1} {#2} }
\cs_new_protected:Npn \__docxlike_format_font_hook:n #1 { }
\cs_new_protected:Npn \docxlike_format_set_font_hook:n #1
  { \cs_gset_protected:Npn \__docxlike_format_font_hook:n ##1 {#1} }
\cs_new_protected:Npn \docxlike_format_emit:n #1
  { \tl_put_right:Nn \l__docxlike_format_out_tl {#1} }
\cs_new_protected:Npn \__docxlike_format_build:n #1
  {
    \__docxlike_format_build_font:n {#1}
    \tl_put_right:Nn \l__docxlike_format_out_tl { \docxlike_format_strut: }
  }
%    \end{macrocode}
%
% 发出去的撑杆中西文要认调用方立的 \cs{l\_docxlike\_format\_latin\_pick\_bool}，与
% \cs{docxlike\_format\_compute:n} 一致。先前这里只看目标写的字体，双语题注的英文那一行
% 行距按西文算了、撑杆却还是宋体的上升部，中文题到英文题的步进 17.0，Word 是 16.5。
%    \begin{macrocode}
\cs_new_protected:Npn \__docxlike_format_build_font:n #1
  {
    \tl_clear:N \l__docxlike_format_out_tl
    \__docxlike_format_latin_line:n {#1}
    \__docxlike_format_font_metrics:n {#1}
    \bool_set:Nn \l__docxlike_format_latin_bool
      {
        \bool_if_p:N \l_docxlike_format_latin_line_bool
        || \bool_if_p:N \l_docxlike_format_latin_pick_bool
      }
    \__docxlike_format_read:nnnN {#1} { rPr/rFonts/eastAsia } { } \l__docxlike_format_value_tl
    \tl_if_empty:NF \l__docxlike_format_value_tl
      {
        \tl_put_right:Ne \l__docxlike_format_out_tl
          {
            \exp_not:N \docxlike_font_select_east_asian:n
              { \exp_not:V \l__docxlike_format_value_tl }
          }
      }
    \prop_get:NnNT \g__docxlike_format_font_extra_prop {#1} \l__docxlike_format_extra_tl
      { \tl_put_right:NV \l__docxlike_format_out_tl \l__docxlike_format_extra_tl }
    \__docxlike_format_read:nnnN {#1} { rPr/rFonts/ascii } { } \l__docxlike_format_value_tl
    \tl_if_empty:NF \l__docxlike_format_value_tl
      {
        \tl_put_right:Ne \l__docxlike_format_out_tl
          {
            \exp_not:N \docxlike_font_select_latin:n
              { \exp_not:V \l__docxlike_format_value_tl }
          }
      }
    \tl_put_right:Ne \l__docxlike_format_out_tl
      {
        \exp_not:N \docxlike_format_font_metrics:nnnn
          { \l__docxlike_format_strut_natural_tl } { \l_docxlike_format_ascent_tl }
          { \l_docxlike_format_natural_latin_tl } { \l_docxlike_format_ascent_latin_tl }
      }
    \__docxlike_format_sz_read:nN {#1} \l__docxlike_format_value_tl
    \tl_if_empty:NF \l__docxlike_format_value_tl
      {
        \tl_put_right:Nx \l__docxlike_format_out_tl
          {
            \exp_not:N \fontsize
              { \dim_use:N \l_docxlike_format_size_dim }
              { \dim_use:N \l_docxlike_format_lead_dim }
            \exp_not:N \selectfont
          }
      }
    \__docxlike_format_read:nnnN {#1} { rPr/kern } { } \l__docxlike_format_value_tl
    \tl_if_empty:NF \l__docxlike_format_value_tl
      {
        \bool_lazy_or:nnTF
          { \int_compare_p:nNn { \l__docxlike_format_value_tl } = { 0 } }
          {
            \dim_compare_p:nNn { \l_docxlike_format_size_dim }
              < { \fp_eval:n { \l__docxlike_format_value_tl / 2 } bp }
          }
          { \tl_put_right:Nn \l__docxlike_format_out_tl { \docxlike_linebreak_kern_off: } }
          { \tl_put_right:Nn \l__docxlike_format_out_tl { \docxlike_linebreak_kern_on: } }
      }
    \__docxlike_format_read:nnnN {#1} { rPr/b } { false } \l__docxlike_format_value_tl
    \str_if_eq:VnT \l__docxlike_format_value_tl { true }
      { \tl_put_right:Nn \l__docxlike_format_out_tl { \bfseries } }
    \__docxlike_format_font_hook:n {#1}
    \__docxlike_format_pick:nnN {#1} { pPr/jc } \g__docxlike_format_align_prop
    \tl_put_right:Nx \l__docxlike_format_out_tl
      {
        \exp_not:c
          {
            docxlike_format_snap_
            \bool_if:NTF \l_docxlike_format_snap_vertical_bool { true } { false } :
          }
        \bool_if:NT \l_docxlike_format_hglue_set_bool
          {
            \exp_not:N \docxlike_format_hglue:nn
              { \dim_eval:n { \__docxlike_format_hglue_cc: } }
              { \dim_eval:n { \__docxlike_format_hglue_ec: } }
          }
        \bool_if:NTF \l__docxlike_format_latin_bool
          { \exp_not:N \docxlike_format_latin_true: }
          { \exp_not:N \docxlike_format_latin_false: }
        \bool_if:NTF \l__docxlike_format_exact_bool
          { \exp_not:N \docxlike_format_exact_true: }
          { \exp_not:N \docxlike_format_exact_false: }
      }
  }
%    \end{macrocode}
%
% \begin{macro}{\docxlike_format_apply_state_else:nn,\docxlike_format_strut_box:N}
% \cs{docxlike\_format\_apply\_state\_else:nn}\marg{目标}\marg{兜底} 与
% \cs{docxlike\_format\_apply\_font\_else:nn} 一样按目标设字体、字号、行距与撑杆的状态，
% 但\emph{不发撑杆}。表格环境的开头是竖直模式，一根 \cs{vrule} 发出去就起了一段；
% 那里要的只是状态，撑杆由 \cs{docxlike\_format\_strut\_box:N} 另装进一只盒子给
% \cs{strutbox}。
%    \begin{macrocode}
\cs_new_protected:Npn \docxlike_format_apply_state_else:nn #1#2
  {
    \__docxlike_format_if_blank:nTF {#1}
      {#2}
      {
        \docxlike_format_compute:n {#1}
        \__docxlike_format_build_font:n {#1}
        \l__docxlike_format_out_tl
      }
  }
\cs_new_protected:Npn \docxlike_format_strut_box:N #1
  { \hbox_set:Nn #1 { \docxlike_format_strut: } }
%    \end{macrocode}
% \end{macro}
%
% 两道胶的取值。\cs{CJKglue} 是汉字之间的：吃网格就取网格那一份，不吃就取零，
% 再加上 \texttt{spacing}。两项相加是 Word 的算法，\texttt{docGrid} 的字距被
% \texttt{snapToGrid=0} 关掉，run 上的 \texttt{w:spacing} 照样生效，
% 章标题正是这个情形（\texttt{snapToGrid=0} 且 \texttt{w:spacing=-4}）。
%
% \cs{CJKecglue} 是汉字与西文之间的，它是两截相加：网格那一份，加上 Word 的
% “自动调整中文与西文的间距”那一份（\texttt{w:autoSpaceDE}/\texttt{DN}，
% \cs{g\_docxlike\_page\_ecglue\_dim} 里算好的 0.225 格）。关掉网格只该去掉前一截，
% 后一截跟网格无关，去掉了中西文之间就一点空隙都没有了。
%
% 0.246 由 1--5 页版反推：\texttt{汉→数字} 松行实测 15.52，减掉一整格 12.4543
% 得 3.07，除以格宽即 0.246。旧值 0.261 是按伪影文档定的，总步进会到 15.70，
% 宽 0.18。
%    \begin{macrocode}
\tl_const:Nn \c__docxlike_format_ecglue_tl { 0.246 }
\dim_new:N \g_docxlike_page_slack_dim
\cs_new:Npn \__docxlike_format_grid_glue:
  { \g_docxlike_page_char_dim - \g_docxlike_body_size_dim }
%    \end{macrocode}
%
% 整数格除不尽的那点余量。一行放得下 $n$ 个格（\cs{g\_docxlike\_page\_cells\_int}，
% 只跟版心宽与字距有关，由版面那层算），$n$ 个字连着排占 $n$ 个字号加 $n-1$ 段
% 字距胶，版心宽减掉它就是余量。字号是元素的量，所以算法在这一层。
%    \begin{macrocode}
\int_new:N \l__docxlike_format_gaps_int
\dim_new:N \l__docxlike_format_glue_dim
\cs_new_protected:Npn \docxlike_format_slack:
  {
    \dim_set:Nn \l__docxlike_format_glue_dim { \__docxlike_format_grid_glue: }
    \int_set:Nn \l__docxlike_format_gaps_int { \g_docxlike_page_cells_int - 1 }
    \dim_gset:Nn \g_docxlike_page_slack_dim
      {
        \textwidth
        - \g_docxlike_page_cells_int \g_docxlike_body_size_dim
        - \l__docxlike_format_gaps_int \l__docxlike_format_glue_dim
      }
  }
\cs_new:Npn \__docxlike_format_hglue_cc:
  {
    \bool_if:NT \l_docxlike_format_snap_horizontal_bool
      { \__docxlike_format_grid_glue: + }
    \l_docxlike_format_char_spacing_dim
  }
\cs_new:Npn \__docxlike_format_hglue_ec:
  {
    \bool_if:NT \l_docxlike_format_snap_horizontal_bool
      { \__docxlike_format_grid_glue: + }
    \__docxlike_format_ecglue_base:
    + \l_docxlike_format_char_spacing_dim
  }
%    \end{macrocode}
%
% 中西文胶按字号缩放，不是正文那一格的定值：样本文档里 18\,bp 的章标题汉字接数字
% 是 22.28\,bp，减掉汉字 18 得 4.28，按字号缩放 $0.225 \times 18 \times 12.6719
% / 12 = 4.2768$ 与实测同位，按正文格宽只有 2.851，“第\,1\,章”的 1 贴得太近。
%
% 正文字号为零时（没走 \cs{docxlike\_page\_setup:n}）退回固定格宽，不做除法。
%    \begin{macrocode}
\cs_new:Npn \__docxlike_format_ecglue_base:
  {
    \dim_compare:nNnTF { \g_docxlike_body_size_dim } = { 0pt }
      { \c__docxlike_format_ecglue_tl \g_docxlike_page_char_dim }
      {
        \fp_to_dim:n
          {
            \c__docxlike_format_ecglue_tl
            * \dim_to_fp:n { \l_docxlike_format_size_dim }
            * \dim_to_fp:n { \g_docxlike_page_char_dim }
            / \dim_to_fp:n { \g_docxlike_body_size_dim }
          }
      }
  }
\cs_new_protected:Npn \__docxlike_format_pick:nnN #1#2#3
  {
    \__docxlike_format_read:nnnN {#1} {#2} { } \l__docxlike_format_value_tl
    \prop_get:NVNT #3 \l__docxlike_format_value_tl \l__docxlike_format_coef_tl
      { \tl_put_right:NV \l__docxlike_format_out_tl \l__docxlike_format_coef_tl }
  }
\cs_generate_variant:Nn \prop_get:NnNT { NVNT }
%    \end{macrocode}
%
% 段前段后写了区间就带一段伸量：\cs{\_\_docxlike\_format\_plus:N} 写出 \texttt{plus} 那一截，
% 伸量为零时一个字也不发。
%    \begin{macrocode}
\cs_new:Npn \__docxlike_format_plus:N #1
  {
    \dim_compare:nNnF { #1 } = { 0pt }
      { ~ plus ~ \dim_use:N #1 }
  }
%    \end{macrocode}
%
% \cs{docxlike\_format\_apply\_font:nN}\marg{目标}\meta{宏} 把目标算出来的那串排版命令
% 存进宏里（页眉页脚在输出例程里排，就是这样拿字体的）。目标是空的就不动那个宏，
% 文档类原先的定义因此原样保留；段前段后等量留在 \cs{l\_docxlike\_format\_before\_dim} 这几个
% 变量里。文档类要问样式里某一处写没写过，用 \cs{docxlike\_format\_if\_set:nnTF}\marg{目标}\marg{槽}
% （槽是 \texttt{before}、\texttt{after} 这些，上级样式写的也算）。
%    \begin{macrocode}
\prg_new_eq_conditional:NNn \docxlike_format_if_set:nn \__docxlike_format_slot_if_set:nn
  { T , F , TF }
\cs_new_protected:Npn \docxlike_format_apply_font:nN #1#2
  {
    \__docxlike_format_if_blank:nF {#1}
      {
        \docxlike_format_compute:n {#1}
        \__docxlike_format_build:n {#1}
        \cs_set_eq:NN #2 \l__docxlike_format_out_tl
      }
  }
%    \end{macrocode}
%
% \emph{页眉页脚另走一条。} 三件事只属于页眉页脚，题注、封面字段这些同样用
% \cs{docxlike\_format\_apply\_font:nN} 的地方不该吃：
% \begin{itemize}
% \item 页眉盒在输出例程里装，\pkg{luatexja} 的参数栈那时候失灵，
%   字距会变成正文的（每字多一格增量）。汉字缝与汉西隙另记在属性上，Lua 装盒时
%   照属性重设，见 \file{docxlike-linebreak.dtx}；西文的 LetterSpace 也一并去掉，
%   页眉页脚不在字符网格上。
% \item 词隙不吃正文那份按网格算死的 \texttt{spaceskip}，回到字体自己的空格
%   （Times 0.25\,em；一份样本的英文页眉量得 2.28，正文是 3.0）。
% \end{itemize}
%    \begin{macrocode}
\cs_new_protected:Npn \docxlike_format_apply_hf:nN #1#2
  {
    \docxlike_format_apply_font:nN {#1} #2
    \__docxlike_format_if_blank:nF {#1}
      {
        \tl_put_right:Nn #2 { \skip_zero:N \spaceskip }
        \bool_lazy_and:nnT
          { \l_docxlike_format_hglue_set_bool } { \cs_if_exist_p:N \docxlike@kanjiskip@attr }
          {
            \tl_put_right:Nx #2
              {
                \exp_not:N \docxlike@kanjiskip@attr =
                  \exp_not:N \tex_dimexpr:D \dim_eval:n { \__docxlike_format_hglue_cc: } \exp_not:N \scan_stop:
                \exp_not:N \docxlike@xkanjiskip@attr =
                  \exp_not:N \tex_dimexpr:D \dim_eval:n { \__docxlike_format_hglue_ec: } \exp_not:N \scan_stop:
              }
          }
      }
  }
%    \end{macrocode}
%
% \begin{macro}{\docxlike_format_apply_font_else:nn}
% \cs{docxlike\_format\_apply\_font\_else:nn}\marg{目标}\marg{兜底} 目标有取值就按它排，
% 没有就执行\meta{兜底}。同一处排版代码既要服务设了 \option{styles} 的文档、
% 也要服务没设的文档时用这个：表里有取值就按它，表是空的就退回原先写死的字号。
%    \begin{macrocode}
\tl_new:N \l__docxlike_format_else_tl
\cs_new_protected:Npn \docxlike_format_apply_font_else:nn #1#2
  {
    \__docxlike_format_if_blank:nTF {#1}
      {#2}
      {
        \docxlike_format_apply_font:nN {#1} \l__docxlike_format_else_tl
        \l__docxlike_format_else_tl
      }
  }
\cs_generate_variant:Nn \docxlike_format_apply_font_else:nn { e }
%    \end{macrocode}
%    \end{macro}
%
% \begin{macro}{\docxlike_format_par:nn}
% \cs{docxlike\_format\_par:nn}\marg{目标}\marg{内容} 按目标的版面排一段。字体、
% 字号、行距、粗体、对齐、贴网格、字距、撑杆都由 \cs{docxlike\_format\_apply\_font:nN}
% 备好，这里只补 \option{all-caps} 和段落自身的缩进与段前段后。封面这类不走
% 章节命令的地方用它排每一段，取值就和正文标题共用一套 \option{styles} 词汇。
%
% \emph{缩进。}目标的排版命令（末尾的撑杆起段）执行完，再把左缩进加进 \cs{leftskip}、右缩进
% 加进 \cs{rightskip}，这两样到段末才读；首行缩进（悬挂时为负）是段首的一个 \cs{kern}。样式的
% 对齐方式（\cs{centering} 这类）会清掉 \cs{parindent}，所以不用它。居中时 \hologo{TeX} 在首行
% 起点到右边距之间居中，Word 也是这样（\file{test/numbering} 里悬挂缩进的居中段）。段落链接了
% 多级列表的某一级时，那一级的缩进压过样式的，首行开头排编号，见 \file{docxlike-numbering.dtx}。
%
% \emph{页顶的段前。}Word 在页顶去掉段前：自然换页、分页符之后都是这样，\hologo{TeX} 在页顶
% 丢掉胶，结果一样。保留的有两种（\file{test/headings} 的 \file{heads-top*}、\file{heads-sect*}）：
% 新一节的第一段（全文第一段也算，两种兼容模式一样），与勾了段前分页的段（兼容模式低于 15；
% Word 2013 起的兼容模式分页后去掉）。所以要让章标题在新页页顶保住段前，照 Word 的做法是
% 分页后插一个连续分节符，本包与 docx 同构，也是这样。
%
% 保留下来的不是整段段前，而是“段前减前一段的段后”，不够减就没有：Word 照常按两段取大排，
% 前一段的段后落在上一页（分节符那一段段后 0、8、20、40\,bp 时，段前 24\,bp 的标题在新页上
% 留 24、16、4、0）。每段结束时记下它的段后：按样式排的段自己记，正文段落开头就记成正文
% 样式的段后。新一节的标记由 \cs{docxsectionbreak} 立起，文档开头也立着，下一段开头时放倒。
%
% 要保留时（外层的竖直列表上、当前页是空的；在盒子里排的不算，比如先装盒再放的封面），
% 先放一只零高的空盒，页顶的 \cs{topskip} 算在它上面，再补“留下的段前减 \cs{topskip}”，
% 首行基线就落在上边距加留下的段前加上升部处。
%
% \emph{段前与前一段的段后取大，不相加。}Word 两段之间的距离是前一段的段后与后一段
% 的段前里大的那个（样本文档里章标题段后 0.8 行、小标题段前 0.5 行，量出来
% 只有 0.8 行；章标题到节标题也是这样）。这里段前也用 \cs{addvspace}，前面是段后那道胶就只补
% 差额；前面是 \cs{parmark*} 的空段（盒子，不是胶）就整个加。段后那道胶之后若还跟着
% 看不见的东西（目录条目的 \cs{write}、超链接锚点、\cs{mark}、罚分），\cs{addvspace}
% 就看不到它了，Word 却照样取大；所以段后放下时记下它的量，下一段开段前中间没有别的段落、
% \cs{lastskip} 又是零，就直接补差额。主竖直列表上 \cs{lastnodetype} 总是 $-1$，
% 分不出最后是什么，只能这样认；中间夹了盒子的只有 \cs{parmark*}，它把记下的量清掉。
% 记与清都只认主竖直列表上的段落：浮动体、脚注这些盒子里的段落（包括盒子里的题注）
% 不记也不清，标题接浮动体再接标题，Word 里仍是两个标题相邻。盒子里的段落只在 \hologo{LuaTeX} 下分得出
% （\texttt{tex.nest.ptr}），别的引擎一律算数。
%    \begin{macrocode}
\tl_new:N \l__docxlike_format_par_tl
\tl_new:N \l__docxlike_format_content_tl
\dim_new:N \l_docxlike_format_left_dim
\dim_new:N \l_docxlike_format_right_dim
\cs_new_protected:Npn \__docxlike_format_side_indent:nnN #1#2#3
  {
    \__docxlike_format_amount_read:nn {#1} {#2}
    \__docxlike_format_amount_use:VVN \l__docxlike_format_unit_tl
      \l__docxlike_format_value_tl #3
  }
\cs_new_protected:Npn \__docxlike_format_number_indent:n #1 { }
\cs_new_protected:Npn \__docxlike_format_number_put:n #1 { }
\cs_new_protected:Npn \docxlike_format_par:nn #1#2
  {
    \group_begin:
      \docxlike_format_apply_font:nN {#1} \l__docxlike_format_par_tl
      \__docxlike_format_side_indent:nnN {#1} { ind-left } \l_docxlike_format_left_dim
      \__docxlike_format_side_indent:nnN {#1} { ind-right } \l_docxlike_format_right_dim
      \__docxlike_format_number_indent:n {#1}
      \dim_zero:N \parskip
      \__docxlike_format_read:nnnN {#1} { pPr/pageBreakBefore } { false } \l__docxlike_format_value_tl
      \bool_set_false:N \l__docxlike_format_top_keep_bool
      \str_if_eq:VnT \l__docxlike_format_value_tl { true }
        {
          \par \newpage
          \int_compare:nNnT { \g__docxlike_linebreak_compat_int } < { 15 }
            { \bool_set_true:N \l__docxlike_format_top_keep_bool }
        }
      \dim_compare:nNnT { \l_docxlike_format_before_dim } > { \c_zero_dim }
        {
          \par
          \bool_lazy_all:nTF
            {
              { ! \mode_if_inner_p: }
              {
                \l__docxlike_format_top_keep_bool
                || \g_docxlike_section_start_bool
              }
              { \dim_compare_p:nNn { \pagegoal } = { \c_max_dim } }
            }
            {
              \hbox:n { }
              \nointerlineskip
              \skip_vertical:n
                {
                  \dim_max:nn { 0pt }
                    { \l_docxlike_format_before_dim - \g__docxlike_format_prev_after_dim }
                  - \topskip
                }
            }
            {
              \bool_lazy_and:nnTF
                { \dim_compare_p:nNn { \g__docxlike_format_pending_after_dim } > { 0pt } }
                { \dim_compare_p:nNn { \tex_lastskip:D } = { 0pt } }
                {
                  \dim_compare:nNnT
                    { \l_docxlike_format_before_dim } > { \g__docxlike_format_pending_after_dim }
                    {
                      \skip_vertical:n
                        {
                          \l_docxlike_format_before_dim
                          - \g__docxlike_format_pending_after_dim
                          \__docxlike_format_plus:N \l_docxlike_format_before_plus_dim
                        }
                    }
                }
                {
                  \addvspace
                    {
                      \l_docxlike_format_before_dim
                      \__docxlike_format_plus:N \l_docxlike_format_before_plus_dim
                    }
                }
            }
        }
      \nointerlineskip
      \__docxlike_format_border_put:nn {#1} { top }
      \__docxlike_format_read:nnnN {#1} { pPr/keepLines } { false } \l__docxlike_format_value_tl
      \str_if_eq:VnT \l__docxlike_format_value_tl { true } { \int_set:Nn \interlinepenalty { 10000 } }
      \dim_zero:N \parindent
      \l__docxlike_format_par_tl
      \skip_set:Nn \leftskip { \leftskip + \l_docxlike_format_left_dim }
      \skip_set:Nn \rightskip { \rightskip + \l_docxlike_format_right_dim }
      \dim_compare:nNnF { \l_docxlike_format_indent_dim } = { \c_zero_dim }
        { \kern \l_docxlike_format_indent_dim }
      \__docxlike_format_number_put:n {#1}
      \tl_set:Nn \l__docxlike_format_content_tl {#2}
      \bool_if:NT \g_docxlike_format_balance_bool
        { \__docxlike_format_balance:N \l__docxlike_format_content_tl }
      \__docxlike_format_read:nnnN {#1} { rPr/caps } { false }
        \l__docxlike_format_value_tl
      \str_if_eq:VnTF \l__docxlike_format_value_tl { true }
        {
          \use:e
            {
              \exp_not:N \__docxlike_format_tabbed:nn {#1}
                { \exp_not:N \MakeUppercase { \exp_not:V \l__docxlike_format_content_tl } }
            }
        }
        { \exp_args:NnV \__docxlike_format_tabbed:nn {#1} \l__docxlike_format_content_tl }
      \par
      \__docxlike_format_border_put:nn {#1} { bottom }
      \__docxlike_format_read:nnnN {#1} { pPr/keepNext } { false } \l__docxlike_format_value_tl
      \tl_gset_eq:NN \g__docxlike_format_keep_tl \l__docxlike_format_value_tl
      \dim_gset_eq:NN \g__docxlike_format_after_dim \l_docxlike_format_after_dim
      \dim_gset_eq:NN \g__docxlike_format_after_plus_dim \l_docxlike_format_after_plus_dim
    \group_end:
    \dim_gset_eq:NN \g__docxlike_format_prev_after_dim \g__docxlike_format_after_dim
    \str_if_eq:VnT \g__docxlike_format_keep_tl { true } { \nobreak }
    \dim_compare:nNnT { \g__docxlike_format_after_dim } > { \c_zero_dim }
      {
        \skip_vertical:n
          {
            \g__docxlike_format_after_dim
            \__docxlike_format_plus:N \g__docxlike_format_after_plus_dim
          }
      }
    \mode_if_inner:F
      { \dim_gset_eq:NN \g__docxlike_format_pending_after_dim \g__docxlike_format_after_dim }
    \str_if_eq:VnT \g__docxlike_format_keep_tl { true } { \__docxlike_format_keep_next: }
  }
\dim_new:N \g__docxlike_format_pending_after_dim

\tl_new:N \g__docxlike_format_keep_tl
\dim_new:N \g__docxlike_format_after_dim
\dim_new:N \g__docxlike_format_after_plus_dim
\bool_new:N \l__docxlike_format_top_keep_bool
\dim_new:N \g__docxlike_format_prev_after_dim
\dim_new:N \g__docxlike_format_body_after_dim
\bool_new:N \g_docxlike_section_start_bool
\bool_gset_true:N \g_docxlike_section_start_bool
\hook_gput_code:nnn { para / begin } { docxlike }
  {
    \bool_gset_false:N \g_docxlike_section_start_bool
    \dim_gset_eq:NN \g__docxlike_format_prev_after_dim \g__docxlike_format_body_after_dim
    \__docxlike_format_if_outer_par:T
      { \dim_gzero:N \g__docxlike_format_pending_after_dim }
  }
\sys_if_engine_luatex:TF
  {
    \cs_new:Npn \__docxlike_format_if_outer_par:T #1
      { \int_compare:nNnT { \lua_now:n { tex.sprint(-2, tostring(tex.nest.ptr)) } } < { 2 } {#1} }
  }
  { \cs_new:Npn \__docxlike_format_if_outer_par:T #1 {#1} }
%    \end{macrocode}
%
% 段后距加在组外：组里设过的颜色出组时复原，复原的节点落在竖直列表上，段后距放在它前面，
% 下一段的 \cs{addvspace} 就看不到这道胶，两段之间成了相加（Word 是取大）。
%
% \emph{与下段同页。}段后距之前放 \cs{nobreak}，下一段开头照 \hologo{LaTeX} 标题后的办法
% （\cs{@nobreak}、\cs{everypar}）不在第一行之前断页。Word 的与下段同页也只保证下一段的
% 第一行跟上（孤行控制开着时是两行）。首行缩进照留，Word 标题后的段落也缩进。
%    \begin{macrocode}
\cs_new_protected:Npn \__docxlike_format_keep_next:
  {
    \legacy_if_set_true:n { @afterindent }
    \use:c { @afterheading }
  }
%    \end{macrocode}
%    \end{macro}
%
% \begin{macro}{\docxlike_format_par:nnn}
% \cs{docxlike\_format\_par:nnn}\marg{目标}\marg{覆盖}\marg{内容} 同上，但先把
% \meta{覆盖} 那几条键值临时按到目标上。存储是全局的，所以进出各存一次、
% 排完还回去，不留痕迹。\meta{覆盖} 相当于 Word 里直接设在这一段上的格式：写了缩进时，
% 链接的多级列表那一级的缩进不再压过它（Word 的次序是直接设的、列表级别的、样式的）。
%    \begin{macrocode}
\prop_new:N \l__docxlike_format_par_save_prop
%    \end{macrocode}
%
% \emph{几张表一起存。} 样式的取值分几张表：不带限定词的一张，每个限定词各一张。临时
% 覆盖写的是不带限定词那张，但写的时候要把各限定词里同名的条目抹掉（不然那里的值会盖住
% 刚写的），所以进出要把几张都存、都还。
%    \begin{macrocode}
\cs_new_protected:Npn \__docxlike_format_save:n #1
  {
    \__docxlike_format_new:nnn {#1} { } { }
    \prop_set_eq:Nc \l__docxlike_format_par_save_prop { g__docxlike_format_#1_prop }
    \clist_map_inline:Nn \g__docxlike_format_layer_clist
      {
        \__docxlike_format_new:nnn {#1} {##1} { }
        \prop_set_eq:cc { l__docxlike_format_par_save_##1_prop }
          { g__docxlike_format_#1_##1_prop }
      }
  }
\cs_new_protected:Npn \__docxlike_format_restore_saved:n #1
  {
    \prop_gset_eq:cN { g__docxlike_format_#1_prop } \l__docxlike_format_par_save_prop
    \clist_map_inline:Nn \g__docxlike_format_layer_clist
      {
        \prop_gset_eq:cc { g__docxlike_format_#1_##1_prop }
          { l__docxlike_format_par_save_##1_prop }
      }
  }
%    \end{macrocode}
%
% \begin{macro}{\docxlike_format_with:nnn}
% \cs{docxlike\_format\_with:nnn}\marg{目标}\marg{覆盖}\marg{代码}：把\meta{覆盖}那几条键值
% 临时按到目标上（Word 里直接设在这一段上的格式），执行\meta{代码}，再还回去。
%    \begin{macrocode}
\cs_new_protected:Npn \docxlike_format_with:nnn #1#2#3
  {
    \__docxlike_format_save:n {#1}
    \__docxlike_format_target:nn {#1} {#2}
    #3
    \__docxlike_format_restore_saved:n {#1}
  }
%    \end{macrocode}
%    \end{macro}
%
%    \begin{macrocode}
\bool_new:N \l_docxlike_format_direct_indent_bool
\cs_new_protected:Npn \docxlike_format_par:nnn #1#2#3
  {
    \tl_if_blank:nTF {#2}
      { \docxlike_format_par:nn {#1} {#3} }
      {
        \__docxlike_format_save:n {#1}
        \__docxlike_format_target:nn {#1} {#2}
        \tl_if_in:nnT {#2} { indentation }
          { \bool_set_true:N \l_docxlike_format_direct_indent_bool }
        \docxlike_format_par:nn {#1} {#3}
        \bool_set_false:N \l_docxlike_format_direct_indent_bool
        \__docxlike_format_restore_saved:n {#1}
      }
  }
%    \end{macrocode}
%    \end{macro}
%
% \subsection{段落边框}
%
% 框线写在样式的 \option{border} 里（见前面“键”一节），存成 docx 的 \texttt{w:pBdr} 那几项：
% 线型、线宽（八分之一磅）、距正文（磅）、颜色。
%
% 量出来的规则（\file{test/headfoot/README.md}）：下边框画在最后一行行盒底之下
% \option{space} 处，线与间隙都算进段高；上边框在第一行行盒顶之上。线比段落左右各伸出
% 1.5 磅。复合线型按线、隙、线……从贴文字的一侧往外排，各段的粗细见下表的写法，
% $w$ 是 \option{sz}/8。点线、虚线类是宽 $w$ 的线按虚线样式画；波浪线没有照画，
% 按同样的高度画一条细线。
%
% 相邻两段边框相同时 Word 合成一组、只在组末画，这里不合并，要合并的请只给最后一段写。
%    \begin{macrocode}
\cs_generate_variant:Nn \__docxlike_format_put:nn { ne }
%    \end{macrocode}
%
% 各线型的构成写成一串 \texttt{R}\marg{粗}（线）与 \texttt{G}\marg{宽}（隙），
% 点线、虚线类写成 \texttt{D}\marg{粗}\marg{虚线样式}。
%    \begin{macrocode}
\cs_new:Npn \__docxlike_format_border_parts:nn #1#2
  {
    \str_case:nnF {#1}
      {
        { nil } { } { none } { }
        { double } { R {#2} G {#2} R {#2} }
        { triple } { R {#2} G {#2} R {#2} G {#2} R {#2} }
        { thinThickSmallGap } { R {#2} G { 0.75bp } R { 0.75bp } }
        { thickThinSmallGap } { R { 0.75bp } G { 0.75bp } R {#2} }
        { thinThickThinSmallGap }
          { R { 0.75bp } G { 0.75bp } R {#2} G { 0.75bp } R { 0.75bp } }
        { thinThickMediumGap } { R {#2} G { #2 / 2 } R { #2 / 2 } }
        { thickThinMediumGap } { R { #2 / 2 } G { #2 / 2 } R {#2} }
        { thinThickThinMediumGap }
          { R { #2 / 2 } G { #2 / 2 } R {#2} G { #2 / 2 } R { #2 / 2 } }
        { thinThickLargeGap } { R { 1.5bp } G {#2} R { 0.75bp } }
        { thickThinLargeGap } { R { 0.75bp } G {#2} R { 1.5bp } }
        { thinThickThinLargeGap }
          { R { 0.75bp } G {#2} R { 1.5bp } G {#2} R { 0.75bp } }
        { threeDEmboss } { \__docxlike_format_border_3d:n {#2} }
        { threeDEngrave } { \__docxlike_format_border_3d:n {#2} }
        { dotted } { D {#2} { {#2} {#2} } }
        { dashed } { D {#2} { { #2 * 3 } {#2} } }
        { dashSmallGap } { D {#2} { { #2 * 3 } {#2} } }
        { dotDash } { D {#2} { { #2 * 3 } {#2} {#2} {#2} } }
        { dotDotDash } { D {#2} { { #2 * 3 } {#2} {#2} {#2} {#2} {#2} } }
        { wave } { G { 1.125bp } R { 0.75bp } G { 1.125bp } }
        { doubleWave } { G { 0.75bp } R { 0.75bp } G { 2.25bp } R { 0.75bp } G { 0.75bp } }
        { dashDotStroked } { G { 1.125bp } R { 0.75bp } G { 1.125bp } }
      }
      { R {#2} }
  }
\cs_new:Npn \__docxlike_format_border_3d:n #1
  {
    R { \dim_min:nn { 1.5bp } { \dim_max:nn { 0.75bp } { #1 / 2 } } }
    R {#1}
    R { \dim_min:nn { 1.5bp } { \dim_max:nn { 0.75bp } { #1 / 2 } } }
  }
%    \end{macrocode}
%
% 画。下边框：先空 \option{space}，再从上往下排各段；上边框反过来，从下往上排，
% 最后空 \option{space}。前后都 \cs{nointerlineskip}：线紧贴行盒，行盒的高深由撑杆定。
%    \begin{macrocode}
\tl_new:N \l__docxlike_format_border_tl
\cs_new_protected:Npn \__docxlike_format_border_put:nn #1#2
  {
    \__docxlike_format_read:nnnN {#1} { pPr/pBdr/#2 } { } \l__docxlike_format_border_tl
    \tl_if_empty:NF \l__docxlike_format_border_tl
      {
        \use:e
          {
            \exp_not:N \__docxlike_format_border_draw:nnnnn {#2}
            \exp_not:V \l__docxlike_format_border_tl
          }
      }
  }
\seq_new:N \l__docxlike_format_border_seq
\cs_new_protected:Npn \__docxlike_format_border_draw:nnnnn #1#2#3#4#5
  {
    \seq_clear:N \l__docxlike_format_border_seq
    \use:e
      {
        \exp_not:N \__docxlike_format_border_collect:w
        \__docxlike_format_border_parts:nn {#2} { \fp_eval:n { #3 / 8 } bp }
        \exp_not:N \q_recursion_tail \exp_not:N \q_recursion_stop
      }
    \str_if_eq:nnT {#1} { top } { \seq_reverse:N \l__docxlike_format_border_seq }
    \nointerlineskip
    \str_if_eq:nnT {#1} { bottom } { \skip_vertical:n { #4 bp } }
    \group_begin:
      \str_if_eq:nnF {#5} { auto }
        { \cs_if_exist:NT \color_select:nn { \color_select:nn { HTML } {#5} } }
      \seq_map_function:NN \l__docxlike_format_border_seq \use:n
    \group_end:
    \str_if_eq:nnT {#1} { top } { \skip_vertical:n { #4 bp } }
    \nointerlineskip
  }
\cs_new_protected:Npn \__docxlike_format_border_collect:w #1
  {
    \quark_if_recursion_tail_stop:N #1
    \str_case:nn {#1}
      {
        { R } { \__docxlike_format_border_rule:w }
        { G } { \__docxlike_format_border_gap:w }
        { D } { \__docxlike_format_border_dash:w }
      }
  }
\cs_new_protected:Npn \__docxlike_format_border_rule:w #1
  {
    \seq_put_right:Ne \l__docxlike_format_border_seq
      { \__docxlike_format_border_line:nn { \dim_eval:n {#1} } { } }
    \__docxlike_format_border_collect:w
  }
\cs_new_protected:Npn \__docxlike_format_border_gap:w #1
  {
    \seq_put_right:Ne \l__docxlike_format_border_seq
      { \skip_vertical:n { \dim_eval:n {#1} } }
    \__docxlike_format_border_collect:w
  }
\cs_new_protected:Npn \__docxlike_format_border_dash:w #1#2
  {
    \seq_put_right:Ne \l__docxlike_format_border_seq
      { \__docxlike_format_border_line:nn { \dim_eval:n {#1} } { #2 } }
    \__docxlike_format_border_collect:w
  }
\cs_generate_variant:Nn \seq_put_right:Nn { Ne }
%    \end{macrocode}
%
% 一段线：宽是行宽加两头各 1.5 磅，装进行宽的盒子里左右各一个 \cs{hss} 吸掉溢出，
% 盒子对外仍报行宽。虚线样式是一串长度，交替为实与空，用 \cs{leaders} 铺满。
%    \begin{macrocode}
\dim_const:Nn \c__docxlike_format_border_over_dim { 1.5bp }
\cs_new_protected:Npn \__docxlike_format_border_line:nn #1#2
  {
    \hbox_to_wd:nn { \linewidth }
      {
        \tex_hss:D
        \tl_if_blank:nTF {#2}
          {
            \tex_vrule:D height ~ #1 ~ depth ~ \c_zero_dim ~
              width ~ \dim_eval:n { \linewidth + 2 \c__docxlike_format_border_over_dim }
          }
          {
            \hbox_to_wd:nn { \linewidth + 2 \c__docxlike_format_border_over_dim }
              {
                \tex_leaders:D \__docxlike_format_border_dashbox:nn {#1} {#2}
                  \tex_hfill:D
              }
          }
        \tex_hss:D
      }
    \nointerlineskip
  }
\bool_new:N \l__docxlike_format_border_on_bool
\cs_new_protected:Npn \__docxlike_format_border_dashbox:nn #1#2
  {
    \hbox:n
      {
        \bool_set_true:N \l__docxlike_format_border_on_bool
        \tl_map_inline:nn {#2}
          {
            \bool_if:NTF \l__docxlike_format_border_on_bool
              {
                \tex_vrule:D height ~ #1 ~ depth ~ \c_zero_dim ~
                  width ~ \dim_eval:n {##1} \scan_stop:
              }
              { \skip_horizontal:n {##1} }
            \bool_set_inverse:N \l__docxlike_format_border_on_bool
          }
      }
  }
%    \end{macrocode}
%
% \subsection{制表位}
%
% 样式的 \option{tabs} 写成“位置 对齐方式 前导符”，内部存成 docx \texttt{w:tabs} 的
% “对齐方式 位置 前导符”；位置从段落左缩进量起，前导符可省，存的是 docx 的名字 \texttt{dot}、
% \texttt{hyphen}、\texttt{underscore}、\texttt{middleDot}（界面上只有图样）。
% 段落里用 \cs{docxtab} 表示一个制表符。制表位用完了按 \option{default-tab-stops}
% （Word 的“默认制表位”，中文 Word 是 21 磅）往后排。
%
% 只排一行：带制表符的段不断行，Word 的页眉页脚几乎都是这样用的。规则是量的：先按
% 左对齐排开，居中制表位让那一截以停靠点为中心、右制表位让那一截右端贴停靠点，都不压到
% 前一截上；整行排好再按段落的对齐方式分剩下的空白。
%    \begin{macrocode}
\dim_new:N \g_docxlike_default_tab_stop_dim
\dim_gset:Nn \g_docxlike_default_tab_stop_dim { 21bp }
\keys_define:nn { docxlike }
  { default-tab-stops .dim_gset:N = \g_docxlike_default_tab_stop_dim }
\msg_new:nnn { docxlike } { tab-outside-paragraph }
  { \token_to_str:N \docxtab\ is~only~allowed~in~a~paragraph~with~a~format~target. }
\NewDocumentCommand \docxtab { }
  { \msg_error:nn { docxlike } { tab-outside-paragraph } }
\seq_new:N \l__docxlike_format_tab_seg_seq
\seq_new:N \l__docxlike_format_tab_stop_seq
\tl_new:N \l__docxlike_format_tabs_tl
\dim_new:N \l__docxlike_format_tab_x_dim
\dim_new:N \l__docxlike_format_tab_start_dim
\box_new:N \l__docxlike_format_tab_seg_box
\box_new:N \l__docxlike_format_tab_line_box
\tl_new:N \l__docxlike_format_tab_kind_tl
\tl_new:N \l__docxlike_format_tab_lead_tl
\dim_new:N \l__docxlike_format_tab_pos_dim
\cs_new_protected:Npn \__docxlike_format_tabbed:nn #1#2
  {
    \tl_if_in:nnTF {#2} { \docxtab }
      { \__docxlike_format_tab_line:nn {#1} {#2} }
      {#2}
  }
\cs_new_protected:Npn \__docxlike_format_tab_line:nn #1#2
  {
    \__docxlike_format_read:nnnN {#1} { pPr/tabs } { } \l__docxlike_format_tabs_tl
    \seq_set_from_clist:NV \l__docxlike_format_tab_stop_seq \l__docxlike_format_tabs_tl
    \seq_set_split:Nnn \l__docxlike_format_tab_seg_seq { \docxtab } {#2}
    \seq_pop_left:NN \l__docxlike_format_tab_seg_seq \l__docxlike_format_border_tl
    \hbox_set:Nn \l__docxlike_format_tab_line_box { \l__docxlike_format_border_tl }
    \dim_set:Nn \l__docxlike_format_tab_x_dim { \box_wd:N \l__docxlike_format_tab_line_box }
    \seq_map_inline:Nn \l__docxlike_format_tab_seg_seq
      {
        \hbox_set:Nn \l__docxlike_format_tab_seg_box {##1}
        \__docxlike_format_tab_next:
        \str_case:VnF \l__docxlike_format_tab_kind_tl
          {
            { center }
              {
                \dim_set:Nn \l__docxlike_format_tab_start_dim
                  { \l__docxlike_format_tab_pos_dim - 0.5 \box_wd:N \l__docxlike_format_tab_seg_box }
              }
            { right }
              {
                \dim_set:Nn \l__docxlike_format_tab_start_dim
                  { \l__docxlike_format_tab_pos_dim - \box_wd:N \l__docxlike_format_tab_seg_box }
              }
          }
          { \dim_set_eq:NN \l__docxlike_format_tab_start_dim \l__docxlike_format_tab_pos_dim }
        \dim_compare:nNnT { \l__docxlike_format_tab_start_dim } < { \l__docxlike_format_tab_x_dim }
          { \dim_set_eq:NN \l__docxlike_format_tab_start_dim \l__docxlike_format_tab_x_dim }
        \hbox_set:Nn \l__docxlike_format_tab_line_box
          {
            \box_use_drop:N \l__docxlike_format_tab_line_box
            \__docxlike_format_tab_leader:n
              { \l__docxlike_format_tab_start_dim - \l__docxlike_format_tab_x_dim }
            \box_use:N \l__docxlike_format_tab_seg_box
          }
        \dim_set:Nn \l__docxlike_format_tab_x_dim
          { \l__docxlike_format_tab_start_dim + \box_wd:N \l__docxlike_format_tab_seg_box }
      }
    \__docxlike_format_read:nnnN {#1} { pPr/jc } { both } \l__docxlike_format_value_tl
    \leavevmode
    \hbox_to_wd:nn { \linewidth }
      {
        \str_case:VnF \l__docxlike_format_value_tl
          {
            { center } { \tex_hss:D \box_use_drop:N \l__docxlike_format_tab_line_box \tex_hss:D }
            { right } { \tex_hss:D \box_use_drop:N \l__docxlike_format_tab_line_box }
          }
          { \box_use_drop:N \l__docxlike_format_tab_line_box \tex_hss:D }
      }
  }
\cs_generate_variant:Nn \seq_set_from_clist:Nn { NV }
%    \end{macrocode}
%
% 找下一个停靠点：自定的制表位里第一个比当前位置靠右的；没有了就取默认制表位的下一个倍数，
% 按左对齐。
%    \begin{macrocode}
\cs_new_protected:Npn \__docxlike_format_tab_next:
  {
    \tl_set:Nn \l__docxlike_format_tab_kind_tl { left }
    \tl_clear:N \l__docxlike_format_tab_lead_tl
    \dim_set:Nn \l__docxlike_format_tab_pos_dim
      {
        \g_docxlike_default_tab_stop_dim
        * ( \int_div_truncate:nn
              { \dim_to_decimal_in_sp:n { \l__docxlike_format_tab_x_dim } }
              { \dim_to_decimal_in_sp:n { \g_docxlike_default_tab_stop_dim } }
            + 1 )
      }
    \seq_map_inline:Nn \l__docxlike_format_tab_stop_seq
      {
        \__docxlike_format_tab_stop:w ##1 ~ \q_mark ~ \q_stop
        \dim_compare:nNnT { \l__docxlike_format_tab_stop_pos_dim } > { \l__docxlike_format_tab_x_dim }
          {
            \tl_set_eq:NN \l__docxlike_format_tab_kind_tl \l__docxlike_format_tab_stop_kind_tl
            \tl_set_eq:NN \l__docxlike_format_tab_lead_tl \l__docxlike_format_tab_stop_lead_tl
            \dim_set_eq:NN \l__docxlike_format_tab_pos_dim \l__docxlike_format_tab_stop_pos_dim
            \seq_map_break:
          }
      }
  }
\tl_new:N \l__docxlike_format_tab_stop_kind_tl
\tl_new:N \l__docxlike_format_tab_stop_lead_tl
\dim_new:N \l__docxlike_format_tab_stop_pos_dim
\cs_new_protected:Npn \__docxlike_format_tab_stop:w #1 ~ #2 ~ #3 \q_stop
  {
    \tl_set:Nn \l__docxlike_format_tab_stop_kind_tl {#1}
    \dim_set:Nn \l__docxlike_format_tab_stop_pos_dim { \__docxlike_format_from_twips:n {#2} }
    \tl_set:Nn \l__docxlike_format_tab_stop_lead_tl {#3}
    \tl_remove_all:Nn \l__docxlike_format_tab_stop_lead_tl { \q_mark }
    \tl_trim_spaces:N \l__docxlike_format_tab_stop_lead_tl
  }
\cs_new_protected:Npn \__docxlike_format_tab_leader:n #1
  {
    \str_case:VnF \l__docxlike_format_tab_lead_tl
      {
        { dot } { \tex_leaders:D \hbox:n { . } \skip_horizontal:n {#1} }
        { hyphen } { \tex_leaders:D \hbox:n { - } \skip_horizontal:n {#1} }
        { underscore } { \tex_leaders:D \hbox:n { \_ } \skip_horizontal:n {#1} }
        { middleDot } { \tex_leaders:D \hbox:n { · } \skip_horizontal:n {#1} }
      }
      { \skip_horizontal:n {#1} }
  }
%    \end{macrocode}
%
% \subsection{接到正文上}
%
% 正文四个量各有各的去处：
% \begin{longtblr}[caption={正文四个量的去处}, label={tab:impl-body-dims}]
%   {colspec={l l l}}
% \toprule
% 量 & 存在哪 & 谁读它 \\ \midrule
% 字号 & \cs{g\_docxlike\_body\_size\_dim} & \cs{normalsize} \\
% 行距 & \cs{g\_docxlike\_body\_lead\_dim} & \cs{normalsize} \\
% 首行基线落点 & \cs{g\_docxlike\_body\_first\_dim} & \cs{topskip} \\
% 首行缩进 & \cs{g\_docxlike\_body\_indent\_dim} & \cs{parindent} \\
% \bottomrule
% \end{longtblr}
%
% 四个都是全局长度而不是直接写死的数：\cs{normalsize} 每次执行都重读，所以
% 设置写在导言区哪里都来得及，写进正文也能改。默认值由文档类给，
% 没写 \texttt{Normal} 就一直是那一份。
%
% 改完长度要再调一次 \cs{normalsize}。\hologo{LaTeX} 选定正文字体是在类加载完那一刻，
% 早于 \cs{AfterPreamble}，只改长度不重选的话这一页排出来的还是旧行距。
%
% \cs{parindent} 要排在 \cs{normalsize} 之后设。Word 那边“首行缩进 2 字符”在存盘时就折算成了
% 绝对长度（样本文档里是 507 缇），换字号不跟着变。
%
% \cs{topskip} 是版心顶到首行基线的距离，跟 Word 的算法对得上：Word 那边这个距离
% 就等于首行基线在自己行盒里的落点，见前面的换算一节。\hologo{TeX} 取的是
% \cs{topskip} 与首行盒高度的较大者，所以行里有比正文高的东西时首行会往下走，
% 这一条 Word 也一样。
%
% \cs{CJKglue} 是汉字之间的空隙，取值是网格字距减字号，由
% \file{docxlike-page.dtx} 算好，伸量见 \cs{c\_\_docxlike\_format\_ccstretch\_tl}。
%
% \cs{CJKecglue} 是汉字与西文之间的空隙，取值两侧合计 0.45 格。9、12、18、24\,bp
% 四个字号量下来两侧合计分别是 4.281、5.632、8.335、11.037，除以格都是 0.451
% 上下，除以字号则从 0.476 漂到 0.460——是格的比例，不是 em 的比例。单侧实测
% 汉字到西文 0.25 格、西文到汉字 0.20 格，\pkg{xeCJK} 只给一个值，取一半 0.225
% 格，合计对得上，单侧位置差 0.3\,bp 以内。
%
% \cs{rightskip} 扣掉行末那一段不该有的空隙，扣多少要看对不对齐网格，两种情形的
% 值不一样。
%
% \emph{对齐网格时扣整数格的余量。} 行宽是整数个格，不是整个版心：样本文档版心宽
% 425.18\,bp、字距 12.4544\,bp，一行 34 个格占 423.45\,bp，右边空着 1.73\,bp
% 不管。这一段不扣掉的话，\hologo{TeX} 会把 34 个字拉满整个版心，每个字缝多
% 0.066\,bp，末字右移 2.25\,bp。扣掉之后连拉开的行也对得上——样本文档里首行缩进
% 两格、只有 31 个字的那一行，被拉开填满剩下的 32 格，末字落点与满行分毫不差。
%
% \emph{不对齐网格时扣一个字距增量。} 这一档的行宽是“版心宽减一个增量”：
% 行末那个字不带字距增量，Word 把行拉到 $425.20 - 0.4543 = 424.75$。标定文档里
% 三段纯汉字满行，两端对齐的两段末字原点都落在 424.751，左对齐那段是自然的
% 423.00，样本文档正文页量出来也是 424.751。不扣的话每行右端多出 0.45\,bp。
%
% 两个数看着像，来路完全不同：一个是整数格除不尽的零头，一个是末字省掉的增量。
% 早先的合成文档里两者混过一次——那份 \file{settings.xml} 没写
% \texttt{compatibilityMode}，Word 退到旧版行为，对齐目标正好是整数个格，
% 与网格开的那一档撞上了。样本文档写的是 15，也就是 Word 2013 以后的算法。
%
% 撑杆挂在段落的首尾两个钩子上，段间的接缝就都按 Word 的行盒算。首尾各挂一根：
% 首行那根定这一段的高，末行那根定这一段的深，只挂一根管不到另一头。标题那边
% 另有一根写在 \option{styles} 里，两根重了也没事，行盒取的是最大值。
%
% 页眉还要单接一次。页眉页脚在输出例程里排，排之前立着“在页眉里”的标志
% （\cs{l\_docxlike\_format\_in\_header\_bool}），段落钩子见了就什么也不做，所以
% 文档类可以在页眉的字体命令里留一个空位 \cs{docxlike\_header\_strut:}，
% 这里把它换成真的撑杆。没留空位的文档类不受影响，所以先看在不在。
%    \begin{macrocode}
\dim_new:N \g_docxlike_body_size_dim
\dim_new:N \g_docxlike_body_lead_dim
\dim_new:N \g_docxlike_body_stretch_dim
\dim_new:N \g_docxlike_body_first_dim
\dim_new:N \g_docxlike_body_indent_dim
\dim_new:N \docxlike@body@depth
%    \end{macrocode}
%
% 西文那半。Word 把西文也排进字符网格，但半角字只占半格，所以每个西文字符后面
% 多的是网格增量的\emph{一半}，不是整份。
%
% 依据是样本文档逐字量的：\texttt{charSpace=2752} 那一节格宽 12.6720\,bp，数字步进
% 6.3360\,bp，正是半格；\texttt{charSpace=1861} 那一节格宽 12.4536\,bp，数字
% 步进 6.2500\,bp，字母比字体自然宽多 0.24\,bp 上下，都落在半份 0.227\,bp 附近。
% 两节的增量不同而结论一致，所以是半份，不是某个定值。
%
% \pkg{fontspec} 的字体特性只能在声明字体时给，而这个键要到导言区才知道，
% 所以在这儿把选中的那一档西文字体原样再调一次。重调的代码由文档类挂在
% \texttt{docxlike/latin-font-reload} 钩子上，这里不另写 \cs{setmainfont}，
% 免得两处漂移。\cs{defaultfontfeatures} 的加号形式是追加，
% 原有的特性都留着。
%
% \option{LetterSpace} 的单位是百分之一个 em，所以要把增量换算成字号的百分比。
%
% \emph{自然宽套着 \cs{dim\_eval:n} 是有意的。} 裸写「\cs{fontdimen}2\cs{font} $+$
% 格宽 $-$ 字号 \texttt{plus} \ldots」时 \cs{glueexpr} 把「$-$ 字号 \texttt{plus}
% 伸量」当成减去一段带伸量的胶，伸缩量连符号一起被减成负的，英文段落就两端对不齐。
%
% 空格也跟着改。\hologo{XeTeX} 的字间距只落在词内，词尾那一份没有，缺的补进
% 空格，于是空格 = 字体空格 + 两份增量。可压量取自然宽的四分之一——实测 Word
% 的上限落在 0.78 与 0.97 之间（12\,bp），四分之一在窗口内。
%
% \emph{伸量是摊派的权重，不是上限。} Word 对空格的\emph{总}拉伸不设上限
% （实测拉到自然宽的 3.17 倍照拉），但它不是把缺口全给空格：字符网格那节的
% 第七条实测过，31 个字的行被拉开填满 32 格，\emph{每条缝均摊}。\hologo{TeX}
% 按各段胶的伸量比例摊派，原先空格的伸量（两倍字体空格，约 6\,bp）是汉字缝
% （0.02 倍行距，约 0.39\,bp）的十五倍，中西文胶（0.08 倍行距）是四倍，
% 于是“近五年以第一作者在 Nature Medicine(2023)、Nature”这种行里 74\,bp 的
% 缺口有 70 落在同一个空格上，纸面上一个大洞。三种缝的伸量设成同一档，缺口
% 就均摊了——同一行实测各缝 12.6 上下，空格自己也在 3.4 倍自然宽，与 Word
% 的两条实测都对得上。空格的压缩量不动，那是按 Word 的窗口另标定的
% （0.25 倍字体空格，压到自然宽的 0.797，Word 实测下沿 0.78）。
%
% \emph{压缩量只给中西文胶，这是道闸门。} Word 打包一行时压的是空格、全角
% 标点与中西文边界的空隙；纯中文行每个字缝都是断点，Word 从不硬挤。\hologo{TeX}
% 的压缩量落在哪种胶上，哪种行就有打包的本钱——放进 \cs{rightskip} 等于全体
% 放行，实测中文行会各拽一个字上来（行末棘轮 79 冲到 118）；放进中西文胶，
% 只有含西文的行拿得到，恰好是“西文断不开才许挤”。量给 0.13 倍行距
% （约 2.5\,bp），把 0.225 格的边界空隙压到贴零为止；全角标点的压缩是另一套
% 模型（贪心分支），主线 \texttt{PunctStyle=plain} 不掺和。
%
% \emph{Lua 断行接管时这一段整个不做。} \pkg{luaotfload} 的 \option{LetterSpace}
% 是全篇全字号一律加，Lua 排的段落把它清零再按“正文字号才加半份”自己补，交回
% \hologo{TeX} 的段（居中的题目、标题）Word 本来不贴格也不加——样本文档英文封面的题目
% “PARTIAL POROUS EXTERNALLY PRESSURIZED”Word 一行放下，这边每字多 0.34\,bp
% 就折到下一行。\option{LetterSpace} 还会让 \pkg{luaotfload} 不做字距对，Word
% 的 \texttt{kern=2} 是开着的。
%    \begin{macrocode}
\cs_new_protected:Npn \docxlike_format_latin_pitch:
  {
    \bool_lazy_all:nT
      {
        { \bool_if_p:N \g_docxlike_latin_pitch_bool }
        { ! \hook_if_empty_p:n { docxlike / latin-font-reload } }
        {
          ! \bool_lazy_and_p:nn
            { \sys_if_engine_luatex_p: }
            {
              \cs_if_exist_p:N \g__docxlike_linebreak_word_bool
              && \bool_if_p:N \g__docxlike_linebreak_word_bool
            }
        }
      }
      {
        \use:x
          {
            \exp_not:N \defaultfontfeatures + [ \exp_not:N \rmfamily ,
              \exp_not:N \sffamily , \exp_not:N \ttfamily ]
              {
                LetterSpace =
                  \fp_eval:n
                    {
                      50 * \dim_to_fp:n
                        { \g_docxlike_page_char_dim - \g_docxlike_body_size_dim }
                      / \dim_to_fp:n { \g_docxlike_body_size_dim }
                    }
              }
          }
        \hook_use:n { docxlike / latin-font-reload }
        \normalsize
        \skip_set:Nn \spaceskip
          {
            \dim_eval:n
              {
                \fontdimen 2 \font
                + \g_docxlike_page_char_dim - \g_docxlike_body_size_dim
              }
            plus \c__docxlike_format_ccstretch_tl \baselineskip
            minus 0.25 \fontdimen 2 \font
          }
        \frenchspacing
      }
  }
%    \end{macrocode}
%
% \begin{macro}{\docxlike_format_reserve_depth:}
% 换页的判据两边不一样。Word 要末行的\emph{行盒底}不越版心底；\hologo{TeX} 的
% 页面构建器比的是 \cs{pagetotal} 与 \cs{pagegoal}，末行的深度挂在
% \cs{pagedepth} 里不算进去，等于只要末行的\emph{基线}不越版心底就行。差的正是
% 一个行深，样本文档里 7.30\,bp，所以底下常常比 Word 多排一行。
%
% 动的是“一栏还能装多少”，也就是 \cs{@colroom} 与 \cs{vsize}，两个都减掉一个
% 行深。\cs{maxdepth} 同时设成行深：深度不超过行深的东西正好落在留出的那一段里，
% 一分不占；更深的东西超出多少算多少。图表公式这类比正文深的块也就一并算对了。
%
% \cs{textheight} 与 \cs{@colht} 一个不动。这两个说的是“版心多高”与
% “页面盒子多高”，Word 那边这两个数就是整个版心：样本文档的 \texttt{sectPr} 里
% 页高 841.90、上边距 107.75、下边距 85.05，版心 649.10\,bp，合 22.899\,cm。
% 留出来的那一段是换页时不许末行盒底越过的边界，不是版心矮了一截。
%
% 分开这两件事还有个好处：页面盒子仍按整个版心打包，封面上的 \cs{vfill}、
% 整页浮动体两头的 \texttt{1fil} 分到的量都不变，位置一点不动，页脚也不用挪。
% 留出来的那一段在 \texttt{ragged} 一档由 \cs{@textbottom} 那道 fil 吸掉；
% \texttt{flush} 一档 \hologo{LaTeX} 原本把 \cs{@textbottom} 设成 \cs{relax}，
% 页面会被拉开去补这一段，所以改发一个刚性的行深，见 \cs{docxlike\_format\_flush\_bottom:}。
%
% \emph{Lua 断行那一档改成 \cs{maxdepth} 为零、一栏照满版心算。} 留一段行深的做法
% 对末行是文字行时正好，末项没有行深时（空段的盒、表格的盒）就多留了：Word 只看
% 盒底，那一段本可以装东西。逐页对拍时表后一串空段 Word 页底比我们多装一个，
% 有一页只差 1.2\,bp，就是这么来的。\cs{maxdepth} 为零时 \hologo{TeX} 把每一项的深度都
% 算进 \cs{pagetotal}，判据就成了“盒底不越版心底”，与 Word 一字不差；一栏的高度
% 照满版心，\cs{@colroom} 与 \cs{vsize} 都不减。前提是每一行的深度都是行盒的
% 深度：Lua 断行给每一行的盒装的正是 Word 的行盒（\file{docxlike-linebreak.dtx} 的
% \texttt{linebox}），段中间的行也一样。版心底齐平那一档（\texttt{flush}）
% \hologo{LaTeX} 把页面拉到末行基线齐 \cs{vsize}，末行行深会挂到版心外，那一档
% 仍走留行深的老路。
%
% \hologo{XeLaTeX} 那一档 \cs{maxdepth} 单独设成零是不够的。行盒的深度是墨迹深度，
% 汉字只有 1.32\,pt；撑杆挂在 \texttt{para/begin} 与 \texttt{para/end} 两个钩子上，
% 段落中间那些行两头都不挨，撑不到。
%
% 撑杆自带 \cs{\_\_docxlike\_format\_parend\_guard:} 封条，把自己焊在前面的字上。
% \hologo{LuaLaTeX} 下没有这道封条时，段末标点与撑杆之间有一段可断的边界胶，
% 以标点收尾、末行又恰好排满的段落会在段后断出一行只有撑杆的空白，原委见
% 前面 \cs{\_\_docxlike\_format\_parend\_guard:} 那一段。封条写在撑杆里而不是挂钩处：断口补撑杆、
% 页眉撑杆这些后来的调用点都自动带上。\hologo{XeLaTeX} 下封条是空的。
% 样本文档正文第 17 页的 \cs{tracingpages}：
% \texttt{maxdepth=6.0} 时 \texttt{t=648.68}，\texttt{maxdepth=0.0} 时
% \texttt{t=650.00}，只多算了 1.32\,pt，离 \texttt{g=651.52} 还差着，那一行照样
% 留在这一页。
%
% 正常满页的容量没有变：版心 649.09\,bp、首行基线在版心顶下 12.18\,bp、行距
% 19.45\,bp，留之前末行基线的上限是 756.84，留之后 749.57，两种算法都是 33 行。
% 只有末行基线落在底下那 7.3\,bp 带子里的页面才会少排一行，而那正是 Word 不许
% 的位置。
%
% \cs{@colroom} 每页由 \hologo{LaTeX} 从 \cs{@colht} 重设，所以要补两句，见
% \cs{docxlike\_format\_patch\_colroom:}。\cs{vsize} 不用补，它在 \cs{output} 末尾从
% \cs{@colroom} 取。
%    \begin{macrocode}
\prg_new_conditional:Npnn \__docxlike_format_if_lua_lines: { T , F , TF }
  {
    \bool_lazy_all:nTF
      {
        { \sys_if_engine_luatex_p: }
        { \bool_if_exist_p:N \g__docxlike_linebreak_word_bool }
        { \bool_if_p:N \g__docxlike_linebreak_word_bool }
        { \bool_if_p:N \g_docxlike_ragged_bottom_bool }
      }
      { \prg_return_true: } { \prg_return_false: }
  }
\cs_new_protected:Npn \docxlike_format_reserve_depth:
  {
    \dim_gset:Nn \docxlike@body@depth
      { \g_docxlike_body_lead_dim - \g_docxlike_body_first_dim }
    \__docxlike_format_if_lua_lines:TF
      {
        \dim_gzero:N \maxdepth
        \dim_gzero:N \@maxdepth
      }
      {
        \dim_compare:nNnT { \docxlike@body@depth } > { \c_zero_dim }
          {
            \dim_gset_eq:NN \maxdepth \docxlike@body@depth
            \dim_gset_eq:NN \@maxdepth \docxlike@body@depth
            \dim_gsub:Nn \@colroom { \docxlike@body@depth }
            \dim_gset_eq:NN \vsize \@colroom
            \AddToHook { begindocument/end }
              {
                \docxlike_format_patch_colroom:
                \bool_if:NF \g_docxlike_ragged_bottom_bool
                  { \docxlike_format_flush_bottom: }
              }
          }
      }
  }
%    \end{macrocode}
% \end{macro}
%
% \begin{macro}{\docxlike_format_patch_colroom:}
% \cs{@colroom} 是“这一栏还剩多少地方”，每页开头由 \cs{@startcolumn} 从
% \cs{@colht} 重设，\cs{clearpage} 那条路上 \cs{@doclearpage} 也重设一次。
% 留出来的那一段每页都会被这两句填回去，所以两处都要补。
% 单栏用不到 \cs{twocolumn} 里那一句，不补。
%
% 补丁里\emph{不能出现 \pkg{expl3} 的变量名}。\cs{patchcmd} 是把宏体反切分再重扫
% 一遍，重扫时的 catcode 是打补丁那一刻的，\pkg{etoolbox} 只替你把 \texttt{@}
% 设成字母，下划线与冒号不管。写 \cs{g\_docxlike\_body\_lead\_dim} 那种名字进去，
% 出版时会断成 \cs{g} 加一串下标，报“Undefined control sequence”。所以行深
% 存在 \cs{docxlike@body@depth} 里，名字只带 \texttt{@}，跟这一族别的量不同款是有意的。
%
% 三个记号之间没有空格记号：控制词后面的空格在切分时就被吞掉了，
% \verb|\global \@colroom \@colht| 与 \verb|\global\@colroom\@colht| 是同一串，
% 所以这一段不必退出 \pkg{expl3} 写。
%
% 打不上就报警告：那时留白只在第一页有效，第二页起换页的位置会跟 Word 差回一行，
% 不该无声无息。
%    \begin{macrocode}
\cs_new_protected:Npn \__docxlike_format_patch_room:N #1
  {
    \patchcmd #1
      { \global\@colroom\@colht }
      { \global\@colroom\dimexpr\@colht-\docxlike@body@depth\relax }
      { }
      { \msg_warning:nnn { docxlike } { page-patch-failed } {#1} }
  }
\cs_new_protected:Npn \docxlike_format_patch_colroom:
  {
    \__docxlike_format_patch_room:N \@startcolumn
    \__docxlike_format_patch_room:N \@doclearpage
  }
%    \end{macrocode}
% \end{macro}
%
% \begin{macro}{\docxlike_format_flush_bottom:}
% \cs{flushbottom} 把 \cs{@textbottom} 设成 \cs{relax}，页面盒子按 \cs{@colht}
% 打包时缺的那一段就摊到行间去了。留出来的行深本来就该空着，摊掉等于白留，
% 所以换成一个刚性的行深：盒子照样填满 \cs{@colht}，末行基线仍然停在
% 版心底减一个行深的地方。
%
% 可压量也给一个行深。整页装的是压不动的东西时，页面盒子会比 \cs{@colht} 高出
% 几个点，报 Overfull；给了可压量，\hologo{TeX} 先把这一段压回去，最坏退回原先
% 那种“末行基线压在版心底”的样子，不会更差。
%
% 参考文献那两处环境里也各发一次 \cs{flushbottom}，所以要改的是 \cs{flushbottom}
% 本身，不是发过之后再补一句。改法是打补丁不是重定义：\cs{flushbottom} 是内核的
% 名字，重定义等于把它算进本包的公开面；打补丁只换掉 \cs{@textbottom} 那一句，
% 别人往里加过什么都留着，加不上还会报出来。
%
% 它是 \cs{DeclareRobustCommand} 建的，宏体在带一个空格的那个名字里，所以用
% \cs{exp\_args:Nc} 把 \texttt{flushbottom\textvisiblespace} 拼出来再打。
%    \begin{macrocode}
\cs_new_protected:Npn \__docxlike_format_patch_flush:N #1
  {
    \patchcmd #1
      { \let\@textbottom\relax }
      { \def\@textbottom{\vskip\docxlike@body@depth minus \docxlike@body@depth} }
      { }
      { \msg_warning:nnn { docxlike } { page-patch-failed } { \flushbottom } }
  }
\cs_new_protected:Npn \docxlike_format_flush_bottom:
  {
    \exp_args:Nc \__docxlike_format_patch_flush:N { flushbottom ~ }
    \flushbottom
  }
%    \end{macrocode}
% \end{macro}
%
% \begin{macro}{\docxlike_format_penalties:}
% 样本文档的 Normal 样式写着 \texttt{<w:widowControl w:val="0"/>}，全文 566 个段落
% 没有一个在段级覆盖它，也就是孤行寡行控制全程关闭：段首一行单独留在页底、
% 段末一行单独跑到下页页顶，Word 都不拦。
%
% \hologo{TeX} 这边管这件事的是四个罚分，含义各不相同：\cs{clubpenalty} 是在段落
% 第一行之后断页，\cs{widowpenalty} 是在段落最后一行之前断页，
% \cs{displaywidowpenalty} 是公式前那几行的寡行，\cs{brokenpenalty} 是在一个
% 断字的行之后断页。后两个 Word 根本没有对应的概念。四个一律归零。
%
% \cs{clubpenalty} 另有一处。\hologo{LaTeX} 内核的 \cs{@afterheading} 写的是
% \cs{clubpenalty}\cs{@M}，标题后头一行之后禁止断页；文档类可以把它改成 1，
% 那样第一个标题之后就一直是 1。那句是 \cs{everypar} 里的局部赋值，而那一层组
% 就是 document，所以设完不还原。这里再设一次零只管第一个标题之前那一段，
% 1 与 0 实际没有分别。本包不动 \cs{@afterheading}，那归文档类管。
%
% 样本文档上实测这一步不改变任何一页：\cs{tracingpages} 里带 \texttt{p=150} 的候选
% 断点有 110 个，罚分是真加上去了，但一次都没改变最终选中的那个断点——页面通常
% 有几个点可伸，badness 落在几十的量级，150 顶不过后面那个不带罚分的候选。
% 归零是为了别人的稿子：样本文档没炸不等于用户的稿子不炸。
%    \begin{macrocode}
\cs_new_protected:Npn \docxlike_format_penalties:
  {
    \int_set:Nn \clubpenalty { 0 }
    \int_set:Nn \@clubpenalty { 0 }
    \int_set:Nn \widowpenalty { 0 }
    \int_set:Nn \displaywidowpenalty { 0 }
    \int_set:Nn \brokenpenalty { 0 }
  }
%    \end{macrocode}
% \end{macro}
%
% \emph{全角标点少一格网格胶。} Word 的字符网格里，每个全角字符占整整一格，标点
% 不例外。\pkg{xeCJK} 只在标点\emph{空白的那一侧}发 \cs{CJKglue}：\texttt{，} 的
% 空白在右，于是 \texttt{点→汉} 有胶而 \texttt{汉→点} 没有；\texttt{（} 的空白在
% 左，于是 \texttt{汉→开} 有胶而 \texttt{开→汉} 没有。结果每个标点恰好少一格胶，
% \texttt{汉，汉} 排出来是 $12.0 + 12.672 = 24.672$，Word 是 $2 \times 12.672
% = 25.344$。
%
% 补的办法是往 \pkg{xeCJK} 的字符类记号表里前置一道 \cs{CJKglue}。用的是
% \cs{xeCJK\_pre\_inter\_class\_toks:nnn}，单下划线前缀，是 \pkg{xeCJK} 给别的
% 宏包用的模块级接口，不是 \cs{\_\_xeCJK\_} 那种私有件。发的是 \cs{CJKglue} 本身
% 而不是某个定值，所以每一级标题各自的字距都跟得上。
%
% 追加版 \cs{xeCJK\_app\_inter\_class\_toks:nnn} 不行，实测三处纹丝不动：
% \pkg{xeCJK} 自己那串记号会把后面的东西吃掉，必须前置。
%
% \emph{前面那句 \cs{nobreak} 不能省。} 胶在 \hologo{TeX} 眼里是合法断点，补进去
% 就等于在汉字与后置标点之间、以及前置标点与汉字之间各开了一个可断处，于是
% “、”“。”这类点号会被断到下一行行首，“（”会被留在行尾。实测第一版补完
% 全文出现 13 处行首点号。\cs{nobreak} 就是 \cs{penalty10000}，摆在胶前面把这个
% 断点的代价抬到禁止。
%
% \emph{版本断言。} 用到的几个接口——\cs{xeCJK\_pre\_inter\_class\_toks:nnn}、
% \cs{xeCJK\_class\_num:n}、\cs{xeCJK\_new\_class:n}、
% \cs{\_\_xeCJK\_save\_CJK\_class:n}、\cs{xeCJK\_copy\_inter\_class\_toks:nnnn}——
% 在 \pkg{xeCJK} 3.9.1（\TeX~Live 2025）与 3.10.5（2026）上都验过。够不着时报错
% 而不是静默少一格——静默走样比报错难查得多。能挡住改名，挡不住保名改行为，
% 后者只能靠每年跟着上游复验。
%    \begin{macrocode}
%    \end{macrocode}
% \emph{Word 的标点压缩是按需的排版资源，不是定值。} 修正过字体提示的测试文档
% （\file{标点压缩测试2}）把同一标点对铺在各个行位上，结果是双态分布：
% \texttt{汉→“} 松行 12.49、紧行 8.24；\texttt{汉→《} 12.49 与 9.00；
% \texttt{；→”} 12.49 与 11.24；\texttt{、→汉} 12.50 与 6.74。1--5 页版同一行里
% \texttt{；→”} 压满 6.25 而 \texttt{用→“} 只压 1.0。合起来的模型很短：
% \emph{每个全角字符自然占一格（字身加网格增量），行紧时排版器从标点的空白侧
% 按需扣，扣的上限约是空白的大小}。样本文档第 18 页“、→“”的 1.50 是两侧上限同时
% 扣满的极端（$12.45-5.77-4.73$），同页“用→《”的 12.50 是松行一点不扣；
% 按定值 kern 实现的话松行也在扣，引号就贴得过近。
%
% 照搬“自然宽=整格、可压量=空白”在 \hologo{TeX} 里走不通，断行会漂：Word 的
% 断行器把标点压缩当免费资源，挤一挤也要把“探讨””塞进第 18 页那行；
% \hologo{TeX} 的 shrink 带 badness，同样的行宁可提前断也不肯压，实测那行
% “探讨””被推去下一行。所以标点对写成\emph{反方向}：自然宽取压紧值
% （kern，断行器看到的宽度与 Word 压完一致，断点因此对齐），后面补一道
% 只伸不缩的胶，行松时把间隙撑回去。紧行伸长率为零，数值自动落在压紧值上。
% 上限沿用实测：闭号右空白半格（$0.5$\cs{ccwd}），开号左侧削 0.38 格，
% 后接开双方相加 0.88 格。松行的撑开是按伸量比例分的，撑不满 Word 的全格，
% 这是两家分配器的固有差距，记档。
%
% 胶是合法断点，所以每道胶前的断行语义都要点名：禁则处（汉→后置、开→汉、
% 开→开、开→后置）压胶前置 \cs{nobreak}；可断处（汉→开）让胶裸着；
% “后置→汉”的断点用 \pkg{xeCJK} 原生胶自带的那一个，行末减半由字符突出
% 负责（见下）；“后置→开”“后置→引”的断点是各自折处里的，断行分支把
% pair~kern 退回去，标点回整格再由突出挂半。
%
% 折处与标点对的胶只在能断行的地方发。盒子里（\cs{hbox}、表格格子、
% \pkg{ulem} 的收词盒）不断行，胶的伸缩不参与，一律换等宽或压紧 kern；
% \pkg{ulem} 遇到分组内的胶和罚分会把分组拆坏（实测报 Too~many~\}'s），
% kern 无害。\texttt{p} 型表列里的段落是普通水平模式，不受影响。
%    \begin{macrocode}
\cs_new_protected:Npn \docxlike_format_punct_oshrink:
  { \skip_horizontal:n { 0pt ~ minus ~ 0.5 \g_docxlike_page_char_dim } }
\cs_new_protected:Npn \docxlike_format_punct_cshrink:
  { \skip_horizontal:n { 0pt ~ minus ~ 0.5 \g_docxlike_page_char_dim } }
%    \end{macrocode}
%
% \emph{行中标点可压半格。} 这两个数只改可压量，自然宽一点不动，所以不影响任何
% 不参与齐行的量测；变的是断行时 \hologo{TeX} 手里有多少压缩资源。
%
% 取值按行末基线扫出来：闭侧 0.44/0.46/0.48/0.50/0.52 对应行差
% 93/88/88/83/83，开侧固定闭侧 0.50 再扫 0.30/0.38/0.44/0.50 得 86/83/81/79，
% 开侧 0.50 往上到 0.70 都是 79，是一片平台。取平台起点的 0.50，它同时是中文
% 排版“行中标点占半格”的通则值。整体行差由 88 降到 79。
%
% 同时扫过汉字字距的伸量（\cs{c\_\_docxlike\_format\_ccstretch\_tl}）：0.005/0.01/0.02/
% 0.03/0.05 得 104/91/88/101/130，0.02 两侧都更差，是最优点，不动。
%
% 这条与 Word 的机制不完全同构：样本文档里 \texttt{闭标点→开引号} 的步进在 1.248
% 到 6.000 之间变动，而那些行的汉字步进都在网格上，说明 Word 把压缩几乎全吃在
% 标点上、汉字网格不动。我们这边压缩由 \hologo{TeX} 按 badness 在汉字胶与标点
% 胶之间分配，只能把资源配比调到最接近，做不到同构。
%
% \emph{标点与标点之间的胶，盒内要换成 kern。} \pkg{ulem} 收词时按胶和罚分
% 切段，而标点对的记号表在 \pkg{xeCJK} 的标点类分组\emph{内部}执行，组还没
% 闭合就切段，直接拆坏分组（测试文档里 \cs{uline} 内的“？】”实测报
% Too~many~\}'s；同位置 v3.2a 之前的定值 kern 就没事）。盒内不断行、胶的
% 伸缩也不参与，等宽 kern 与自然宽的胶输出相同，于是每对做两个分支：
% 盒内发一格增量的 kern，盒外发 \cs{nobreak} 加可压胶。汉字与标点之间的
% 三对（汉→后、开→汉、汉→开）在分组外执行，发胶时 \cs{uline} 实测无恙，不用换。
%    \begin{macrocode}
\cs_new_protected:Npn \__docxlike_format_punct_pair_kern:
  { \tex_kern:D \dim_eval:n { 0.5 \g_docxlike_page_char_dim - 1em } \scan_stop: }
\cs_new_protected:Npn \docxlike_format_punct_ngap:
  {
    \mode_if_inner:TF
      { \tex_kern:D \dim_eval:n { \g_docxlike_page_char_dim - 1em } \scan_stop: }
      { \nobreak \CJKglue }
  }
\cs_new_protected:Npn \docxlike_format_punct_qtrim:
  {
    \discretionary { } { } { \tex_kern:D \dim_eval:n { -0.38 \g_docxlike_page_char_dim } \scan_stop: }
    \mode_if_inner:F
      { \nobreak \skip_horizontal:n { 0pt ~ plus ~ 0.38 \g_docxlike_page_char_dim } }
  }
\cs_new_protected:Npn \docxlike_format_punct_opair:
  {
    \mode_if_inner:TF
      { \tex_kern:D \dim_eval:n { \g_docxlike_page_char_dim - 1em } \scan_stop: }
      {
        \nobreak
        \skip_horizontal:n
          { \dim_eval:n { \g_docxlike_page_char_dim - 1em } ~ minus ~ 0.38 \g_docxlike_page_char_dim }
      }
  }
\cs_new_protected:Npn \docxlike_format_punct_cpair:
  {
    \__docxlike_format_punct_pair_kern:
    \mode_if_inner:F
      { \nobreak \skip_horizontal:n { 0pt ~ plus ~ 0.5 \g_docxlike_page_char_dim } }
  }
\cs_new_protected:Npn \docxlike_format_punct_bpair:
  {
    \__docxlike_format_punct_pair_kern:
    \mode_if_inner:F
      { \nobreak \skip_horizontal:n { 0pt ~ plus ~ 0.5 \g_docxlike_page_char_dim } }
    \mode_if_inner:TF
      { \tex_kern:D \dim_eval:n { -0.38 \g_docxlike_page_char_dim } \scan_stop: }
      {
        \discretionary
          { \tex_kern:D \dim_eval:n { 1em - 0.5 \g_docxlike_page_char_dim } \scan_stop: }
          { }
          { \tex_kern:D \dim_eval:n { -0.38 \g_docxlike_page_char_dim } \scan_stop: }
        \nobreak \skip_horizontal:n { 0pt ~ plus ~ 0.38 \g_docxlike_page_char_dim }
      }
  }
\cs_new_protected:Npn \docxlike_format_punct_cbreak:
  {
    \__docxlike_format_punct_pair_kern:
    \mode_if_inner:F
      {
        \discretionary
          { \tex_kern:D \dim_eval:n { 1em - 0.5 \g_docxlike_page_char_dim } \scan_stop: }
          { } { }
        \nobreak \skip_horizontal:n { 0pt ~ plus ~ 0.5 \g_docxlike_page_char_dim }
      }
  }
%    \end{macrocode}
%
% \emph{波浪线进场补半格。} 数字侧的西文字距是按字符加的，\pkg{fontspec} 会删掉
% 一段西文末字的尾距，\texttt{3→～} 那半格就丢在这里。出场方向（\texttt{～→6}）
% 提示修正后的测试文档实测 12.50，一整格增量，直接用 \cs{CJKglue}。
%    \begin{macrocode}
\cs_new_protected:Npn \docxlike_format_punct_xtra:
  { \tex_kern:D \dim_eval:n { ( \g_docxlike_page_char_dim - 1em ) / 2 } \scan_stop: }
%    \end{macrocode}
%
% \emph{行末闭号 Word 恒压半格，这边不追了，记下为什么。} 试过三条路：
% 封掉原生胶断点逼断行走折处，断点集一变全文断行漂移；折处与原生胶并存让
% \hologo{TeX} 按代价挑，行末闭号时半时全；字符突出（\cs{XeTeXprotrudechars=2}
% 加 \cs{rpcode}，挂 \cs{xeCJK\_select\_font:} 逐字体设，机制本身可用）——
% 突出量在断行账本上不是免费的，挂半格的行要多撑半格，与折处减半同价，
% \hologo{TeX} 照样不选；行末闭号后面还跟着类记号的节点，突出常被挡住。
% Word 的断行是逐行贪心、压缩不计代价，\hologo{TeX} 是全段最优、伸缩都算钱，
% 这半格是两个模型的固有差距。
%
% \emph{开引号与波浪线各单开一类。} 开引号““”“‘”进 \texttt{DocxlikeQuote}：
% 行中左侧削 0.38 格（压紧值做进 \cs{discretionary} 的不断分支，断行处与行首
% 自动不削），行松时靠 \cs{docxlike\_format\_punct\_qtrim:} 里那道只伸的胶撑回来。
% 书名号括号不削（样本文档 \texttt{用→《} 12.50），所以引号不能与 \texttt{FullLeft}
% 同类。“～”单开 \texttt{DocxlikeTilde}：不参与压缩，与数字相邻也不发中西文胶。
% 建类照抄 \pkg{xeCJK} 处理破折号的 \texttt{PoZheHao} 三步：
% \cs{xeCJK\_new\_class:n} 建类、\cs{\_\_xeCJK\_save\_CJK\_class:n} 注册类号
% （漏这步字体查类号报 Missing number，早先自建类失败就败在这里）、
% \cs{xeCJK\_copy\_inter\_class\_toks:nnnn} 逐对克隆素表（引号克隆
% \texttt{FullLeft}，波浪线克隆 \texttt{FullRight}）。克隆要在发规则之前做；
% 对象类清单用 \cs{g\_\_xeCJK\_class\_seq} 的快照，不写死——2025 的
% \pkg{xeCJK} 3.9.1 还没有 \texttt{PoZheHao}。
%
% “……”“——”并入汉字类：样本文档里它们与汉字步进完全相同，整格、可断、
% 不压缩，连“、→…”“。→…”也是整格（Word 的压缩不作用于非标点邻位）。
% \cs{xeCJKDeclareCharClass} 在 \cs{begin\{document\}} 时刻改已有类照样生效，
%
% \emph{源码在后置标点后换行会丢配对。} “、”后换行，行尾换行符变成空格记号，
% \pkg{xeCJK} 吞掉它之后““”是从 \texttt{Boundary} 类入场的，上面按类配的
% 规则一条都轮不上——测试文档里同一行的两处“”、“”曾一处 1.49 一处 12.00，就是
% 这么来的。开 \texttt{CheckFullRight} 让 \pkg{xeCJK} 在记号层直接删掉后置标点
% 后的空格，两个字符恢复相邻，配对保住。代价是“。~A”这种全角标点后手打的
% 西文空格也会被删；中文标点后本不该打空格，真要空可以写 \verb|~|。
%
% “后置→西文”的中西文胶直接前置发胶原语：字母数字要那道空隙，
% 西文标点如今在半角类里，不从这对走，不会错发。
%
% \emph{闭号接西文不减半、不加中西文胶}，Word 是全宽加约 0.25\,bp 微隙
% （关键词行“；t”实测 12.30）。
%
% \emph{半角标点类的三处修正。} 西文句读、括号、乘除号声明进
% \texttt{HalfLeft}/\texttt{HalfRight} 后，它们与字母数字分了家，三处按
% 1--5 页版实测配平：汉与开引号到半角标点，补一格网格增量（\texttt{汉→,}
% 与 \texttt{“→,} 都是一格出头）；\pkg{xeCJK} 原生在“半角标点→汉”方向
% 发一道 \cs{CJKecglue}，Word 没有（\texttt{,→汉} 实测裸宽），用
% \cs{xeCJK\_replace\_inter\_class\_toks:nnnn} 从素表里摘掉；
% “汉→半角开号”原生也发 \cs{CJKecglue}，按实测降格成 \cs{CJKglue}。
% “(→x”这类西文互邻的边界两类之间没有记号表，天然裸排，
% 英文正文里的逗号句号不受影响。
%
% \emph{乘除号不另立字符类。} 试过三版把 \texttt{×} 的字形掰成中文字体：挪进
% 汉字类，间距塌（正文“×、÷”断点翻三行，棘轮 78 冲 81）；单列一类照抄
% \texttt{HalfRight} 记号表，字体切换藏在表里被一起抄没；只抄非 \texttt{Default}
% 的类对，\cs{xeCJK\_class\_group\_end:} 是配对的分组、拆单边必失衡。字形的事
% 交给字符自己：要宋体形的占位叉由调用方按汉字排，
% 类表一个不动。
% \emph{源码里打的空格也走中西文胶。} \pkg{xeCJK} 默认把“汉字、空格、字母”里的
% 空格留成字体自带的词间空格（12\,bp 下 3.695，伸 6.47 缩 0.81），与不打空格的
% 3.531 不是一个值，同一段文字打不打空格排出来两样。Word 的中西文间距不看有没有
% 空格，所以把 \cs{xeCJK\_space\_or\_xecglue:} 直接设成发 \cs{CJKecglue}，两种写法
% 同宽。包里有个 \texttt{xCJKecglue} 开关能达到同样效果，但它连带换掉
% \cs{xeCJK\_make\_space\_node:} 一套，与我们改过的类对撞车，\TeX~Live 2025 上整篇
% 会在 \cs{\_\_xeCJK\_boundary\_group\_end:n} 报 Extra~\}，只取需要的那一条。
% 反向（字母数字→汉字）的中西胶要不要去掉网格增量另说，真要做得走
% \cs{xeCJK\_replace\_inter\_class\_toks:nnnn} 配 \texttt{Default} 到 \texttt{CJK}
% 的类对。
%
% \textbf{verb 边界的中西文胶补给老版。}\pkg{xeCJK} 2026 重写中西文边界后，
% 行内 \cs{verb} 与前面的汉字之间也空一道中西文胶，2025 及更早不空，
% 差 3.5\,bp，还会被两端对齐摊满整行。老版不空的原因是组边界会丢「上一个
% 字是汉字」的状态（2026 靠新增的 boundary 捕获管道记着），所以补法是把
% 这条管道搭出来：每个“某类→Boundary”的类对把全局布尔清掉、唯
% “CJK→Boundary”置位，“Boundary→Default”处若布尔在且身在 verb 里就
% 手动补一道中西文胶——直接发私有的 \cs{\_\_docxlike\_format\_ecglue\_emit:}，
% 不发 \cs{CJKecglue}：老版的 verb 字体钩子（\cs{\_\_xeCJK\_verb\_font\_hook:}
% 的 \texttt{nobreak\_skip\_zero} 一档）会把后者在 verb 组里垫成零，这正是
% 老版 verb 无胶的另一半机理；2026 插的是边界外捕获的胶值，同样不吃组内
% 清零，两边语义一致。位置与 2026 插的同一处。“身在 verb”由
% \cs{verbatim@font} 里的局部布尔说了算（\pkg{xeCJK} 自己也挂那儿），
% verb 组一关自动失效，别处的组边界不受影响。
%
% 除 verb 之外，2026 还把一份“排出西文的命令”名单注册成边界流
% （\cs{cs\_if\_exist:NT} \cs{cs} 那一批条件注册），命令前有汉字就插
% 中西文胶——示例包的 \cs{cs} 撞名在册，老版没这机制就差 3.5\,bp。
% 给老版按名单补：\texttt{cmd/<名>/before} 钩子在宏展开时刻执行，组还
% 没开、\pkg{xeCJK} 的“上一个是汉字”状态还在，判到就补同一道胶。
% 名单只抄实测出过差的（\cs{cs}、\cs{eqref}），其余边界流暂不移植。
% 上标的后边界另算：上标命令的实现是一整段只装着上标盒的数学模式，
% 2026 的边界体系把它当作正文的附属，尾后紧跟汉字时不发中西文胶；
% 老版按“数学模式结束接汉字”照发 3.5\,bp。测试文档里“\cs{cite} 带页码
% 再接省略号”正落在这条缝上，跨版本会翻断点。补法：老版给
% \cs{textsuperscript} 尾追加一只零宽空盒，盒把字符类状态重置回
% 边界类，这道胶就发不出来了。一般的行内公式尾（如 $x^2$ 接汉字）
% 两边都发胶，不在此列，不动。
%
% verb 的\emph{后边界}同理：2026 在 verb 内容与后面的汉字之间也空一道胶，
% 老版不空。出组那一侧“身在 verb”的局部布尔已经失效，改用一个全局
% 布尔接力——任何“某类→Boundary”时把“刚离开 verb”记下（是否在 verb
% 里，那一刻还读得到），“Boundary→汉字”处见着就补胶并清掉。
%
% 文本组的两侧边界是同一族更大的差：\cs{num}、\cs{ref}、\cs{mbox} 这类
% 命令排出来的内容裹着组或整段落在 whatsit 后头，2026 的边界体系拿组当
% 透明的，两侧照发中西文胶；老版一到组界就把字符类状态断掉，两侧全不发。
% 示例文档里 \cs{num}（“500～1000 字”一段）与 \cs{ref}（“后文
% 2.10.6 中”）都是整段翻断点的实源。全量对齐做不到（任何文本组都算），
% 只修实际出现的：进入侧把 \cs{num}、\cs{ref} 挂进上面 \texttt{cmd} 钩子
% 名单；离开侧把这两条命令包一层，尾后置“刚出组”布尔，复用 verb 的
% “Boundary→汉字”消费点补胶。布尔是全局的，怕它越过组尾漂到后文，
% 凡“Boundary→西文/标点”与每段开头都顺手清一次。
%
% 链接的\emph{前}边界不进上面的名单：2026 按链接内容的头一个节点定夺，
% 是字符就发、是盒（\cs{uline} 划线的内容整个是盒串）就不发，钩子在
% 内容排出来之前看不见这个，宁缺毋滥。示例文档里的链接内容不是划着线
% 就是顶着标点，都落在“不发”的一侧，素着排两版本来就一致。
%
% 链接（\cs{href}、\cs{cite} 这些）的后边界老版也全不发，统一从
% \cs{Hy@EndAnnot} 出手：每个标注收尾时把“刚出组”记上，补不补由
% 消费点定夺。上标里带着链接（\cs{cite} 就是）会把这布尔漏到上标
% 外头，所以上标包装在垫完盒后顺手清掉；目录、页眉这些每行单开一只
% 盒，行首上一节点是空（类型值 $-1$），消费点一并挡住不发。名字里
% 带 at 号的 \cs{Hy@EndAnnot} 在 expl3 语法下会拆词，只能走 c 型
% 按名存取。
%
% 消费时还得看一眼
% 最后的节点是不是西文字符：出组后隔了一记 \texttt{\string~} 这样的胶再遇
% 汉字，2026 不会再叠一道中西文胶（“表~\cs{ref}~所示”正是这种），
% 老版 xeCJK 的边界检查（\cs{xeCJK\_check\_for\_glue:} 一族）在“汉字、
% 空格或 \texttt{\string~}、西文”这样已经隔着胶的边界照样裸调
% \cs{CJKecglue}，间距叠成双份；2026 会把胶与中西文胶折叠成一道。
% \cs{CJKecglue} 正好是我们自己的发胶宏，中间垫一层：老版启动时把它
% 换成带闸的版本，上一节点已是胶（\cs{lastnodetype} 为 11）就不再发。
% 2026 的束定路径读的是寄存器，不经这里，直通版原样保留。
%
% 补发认“组尾还贴着”的情形：上一节点是胶（\cs{lastnodetype} 为 11)
% 说明中间已经有 \texttt{\string~} 这样的间隔，不再补；组尾的字符、
% 链接锚点这类节点则照补。
% 挂接得在类加载期做、赶在命令定义之前——\texttt{cmd} 钩子对定义在先的
% robust 命令补挂不进宏体（静默失败），2026 自己也是先挂后等定义注入。
%
% 已知未拉平：verb \emph{内部}的空格两版是不同的复合胶（2026 会被两端
% 对齐拉宽到 28.9\,bp、老版定宽 12.9，实测钉 \cs{spaceskip} 只压不平——
% 里头还有 \pkg{xeCJK} 各版的标点链），含空格的长 verb 串会让整段换
% 断点。写示例与文档时 verb 串里能省的空格就省掉。“身在 verb”的布尔在
% 全部版本生效。
% 版本指纹用
% \cs{\_\_xeCJK\_use\_ecglue\_skip:}——它是 2026 边界重写才有的名字，在就
% 说明胶已由 \pkg{xeCJK} 自己插、不挂，免得双份。DocxlikeQuote 与 DocxlikeTilde
% 两个自定义类建在排版期，它们到 Boundary 的清理挂接跟着建类走，见
% \cs{docxlike\_format\_punct\_grid:}。
%    \begin{macrocode}
\bool_new:N \g__docxlike_format_boundary_cjk_bool
\bool_new:N \l__docxlike_format_in_verb_bool
\cs_new_protected:Npn \__docxlike_format_verb_ecglue:
  {
    \bool_lazy_and:nnT
      { \l__docxlike_format_in_verb_bool } { \g__docxlike_format_boundary_cjk_bool }
      { \__docxlike_format_ecglue_emit: }
    \bool_gset_false:N \g__docxlike_format_boundary_cjk_bool
  }
\cs_new_protected:Npn \__docxlike_format_boundary_cmd_glue:
  { \xeCJK_if_last_node:nT { CJK } { \__docxlike_format_ecglue_emit: } }
%    \end{macrocode}
%
% \cs{g\_docxlike\_format\_verb\_exit\_bool} 记着“刚从行内代码里出来”，出来之后头一个汉字
% 前要补汉西隙。文档类里有类似的组（比如自己定义的缩略语命令）也想要这一段，就在组尾调
% \cs{docxlike\_format\_group\_exit:}；需要清掉时把布尔量放下。
%    \begin{macrocode}
\bool_new:N \g_docxlike_format_verb_exit_bool
\cs_new_protected:Npn \__docxlike_format_sup_orig:n #1 { }
\cs_new_protected:Npn \__docxlike_format_ecglue_guard:
  { \__docxlike_format_ecglue_emit: }
\cs_new_protected:Npn \docxlike_format_group_exit:
  { \bool_gset_true:N \g_docxlike_format_verb_exit_bool }
\cs_new_protected:Npn \__docxlike_format_num_orig: { }
\cs_new_protected:Npn \__docxlike_format_endannot_orig: { }
\cs_new_protected:Npn \__docxlike_format_ref_orig: { }
\cs_new_protected:Npn \__docxlike_format_verb_exit_mark:
  {
    \bool_if:NTF \l__docxlike_format_in_verb_bool
      { \bool_gset_true:N \g_docxlike_format_verb_exit_bool }
      { \bool_gset_false:N \g_docxlike_format_verb_exit_bool }
  }
\cs_new_protected:Npn \__docxlike_format_exit_clear_hook:n #1
  { \xeCJK_pre_inter_class_toks:nnn { Boundary } {#1}
      { \bool_gset_false:N \g_docxlike_format_verb_exit_bool } }
\cs_new_protected:Npn \__docxlike_format_verb_hook_old:n #1
  { \xeCJK_pre_inter_class_toks:nnn {#1} { Boundary }
      {
        \bool_gset_false:N \g__docxlike_format_boundary_cjk_bool
        \__docxlike_format_verb_exit_mark:
      } }
\cs_new_protected:Npn \__docxlike_format_hook_cmd_glue:n #1
  {
    \AddToHook { cmd / #1 / before }
      { \__docxlike_format_boundary_cmd_glue: }
  }
\AddToHook { package / xeCJK / after }
  {
    \cs_if_exist:NF \__xeCJK_use_ecglue_skip:
      {
        \clist_map_function:nN { cs , eqref , num , ref }
          \__docxlike_format_hook_cmd_glue:n
        \AtBeginDocument
          {
            \cs_gset_eq:NN \__docxlike_format_sup_orig:n \textsuperscript
            \cs_gset_protected:Npn \textsuperscript #1
              {
                \__docxlike_format_sup_orig:n {#1} \hbox:n { }
                \bool_gset_false:N \g_docxlike_format_verb_exit_bool
              }
            \cs_if_exist:cT { Hy@EndAnnot }
              {
                \cs_gset_eq:Nc \__docxlike_format_endannot_orig: { Hy@EndAnnot }
                \cs_gset_protected:cpn { Hy@EndAnnot }
                  {
                    \__docxlike_format_endannot_orig:
                    \bool_if:NF \l_docxlike_format_in_header_bool
                      { \docxlike_format_group_exit: }
                  }
              }
            \AddToHook { begindocument / end }
              {
                \cs_if_exist:NT \num
                  {
                    \DeclareCommandCopy \__docxlike_format_num_orig: \num
                    \RenewDocumentCommand \num { o m }
                      {
                        \IfValueTF {#1}
                          { \__docxlike_format_num_orig: [#1] {#2} }
                          { \__docxlike_format_num_orig: {#2} }
                        \docxlike_format_group_exit:
                      }
                  }
                \DeclareCommandCopy \__docxlike_format_ref_orig: \ref
                \RenewDocumentCommand \ref { s m }
                  {
                    \IfBooleanTF {#1}
                      { \__docxlike_format_ref_orig: * {#2} }
                      { \__docxlike_format_ref_orig: {#2} }
                    \docxlike_format_group_exit:
                  }
              }
          }
      }
  }
\AtBeginDocument
  {
    \bool_lazy_and:nnT
      { \cs_if_exist_p:N \xeCJK_pre_inter_class_toks:nnn }
      { ! \cs_if_exist_p:N \__xeCJK_use_ecglue_skip: }
      {
        \clist_map_function:nN
          { Default , HalfLeft , HalfRight , FullLeft , FullRight }
          \__docxlike_format_verb_hook_old:n
        \xeCJK_pre_inter_class_toks:nnn { CJK } { Boundary }
          {
            \bool_gset_true:N \g__docxlike_format_boundary_cjk_bool
            \__docxlike_format_verb_exit_mark:
          }
        \xeCJK_pre_inter_class_toks:nnn { Boundary } { Default }
          {
            \__docxlike_format_verb_ecglue:
            \bool_gset_false:N \g_docxlike_format_verb_exit_bool
          }
        \cs_gset_protected:Npn \__docxlike_format_ecglue_guard:
          {
            \int_compare:nNnF { \tex_lastnodetype:D } = { 11 }
              { \__docxlike_format_ecglue_emit: }
          }
        \clist_map_function:nN
          { HalfLeft , HalfRight , FullLeft , FullRight }
          \__docxlike_format_exit_clear_hook:n
        \xeCJK_pre_inter_class_toks:nnn { Boundary } { CJK }
          {
            \bool_if:NT \g_docxlike_format_verb_exit_bool
              {
                \bool_lazy_or:nnF
                  { \int_compare_p:nNn { \tex_lastnodetype:D } = { 11 } }
                  { \int_compare_p:nNn { \tex_lastnodetype:D } = { -1 } }
                  { \__docxlike_format_ecglue_emit: }
              }
            \bool_gset_false:N \g_docxlike_format_verb_exit_bool
          }
      }
    \cs_set_protected:Npx \verbatim@font
      {
        \exp_not:o \verbatim@font
        \exp_not:N \docxlike_format_verb_enter:
      }
    \cs_set_protected:Npn \@setupverbvisiblespace
      {
        \cs_set_eq:NN \@xobeysp \docxlike_format_verb_space_visible:
        \@setupverbvisibletab
      }
  }
%    \end{macrocode}
%
% \subsection{行内代码}
%
% \cs{verb} 里的一串要能一字一字映射到一份 Word 文档里：Consolas 五号的 run（没装
% Consolas 就是 Courier New），空格是 U+02FD 加 U+200B，长串的分隔符前是 U+200B。
% 这三样都是 Word 原生的字符与属性，Word 排出来什么样，这里就排成什么样；断点也照
% Word 的正文规则走（连字符后只有模式 15 可断、不认 \texttt{/}、超长串起行按字形硬断）。
% 用 Mac Word 16.89 试过三轮才定下来，探针在 \file{lb/code/}：
%
% \emph{空格。} 内核把 \cs{verb} 的空格做成 \cs{nobreakspace}，不可断却带字体的伸缩量，
% 两端对齐时一行里被拉到 28.9\,bp，\verb|a  b| 数不出是几个。Word 里做得到「一格、
% 刚性、可断」的写法是不断行空格 U+00A0 后跟零宽可断 U+200B：Word 不拉伸 NBSP，
% 在 U+200B 处断，行尾剩一个 NBSP 无害。要看得见就把 U+00A0 换成 U+02FD（modifier
% letter shelf）：Consolas 与 Courier New 都自带这个字，一格宽，形状是贴着基线的
% 小开口方框，正是 \cs{verb*} 那种记号；真正的 U+2423 这两副字都没有，Word 会换字体。
% 这里 U+200B 不发字符，发一个零罚分：Word 导出的 PDF 文本层里也没有它。
%
% \emph{长串。} 超过 16 个字符的串在 \texttt{/}、\texttt{.}、\texttt{-} 前面可断
% （Chicago 14.18 排 URL 的规矩：分隔符落到下一行行首，读者看得出它是串的一部分），
% 头尾四个字符内不断，扩展名（点后不超过五个字符）前不断，连续分隔符与 \texttt{://}
% 之间不断，\texttt{\_} 不断。Word 里就是在这些位置敲 U+200B。这一步在
% \hologo{LuaLaTeX} 的断行层做（串已经切成节点了，规则要看整串），
% \hologo{XeLaTeX} 下没有。
%
% \emph{字号与网格。} run 五号（等宽字体按 7/8 缩放，小四正文里正好五号）、
% 不贴字符网格：Word 里非正文字号的西文本来就不加
% 半格增量，rPr 上再写 \texttt{snapToGrid=0} 是保险；这边 Lua 见了代码属性也不加。
%
% 连字符后压不住：U+2060 与 U+FEFF 放在连字符后 Word 都不理，非断行连字符 U+2011
% 两副字都没字形。所以模式 15 的文档里 \verb|make -C| 照样会断成“make -”加“C”，
% Word 也是；默认的模式 11 不会。
%
% \env{verbatim} 环境（竖排模式进来）不动。
%    \begin{macrocode}
\sys_if_engine_luatex:T { \newattribute \docxlike@code@attr }
\cs_new_protected:Npn \docxlike_format_verb_space_visible:
  { \leavevmode \tex_char:D "02FD \scan_stop: \tex_penalty:D \c_zero_int }
\cs_new_protected:Npn \docxlike_format_verb_space_cell:
  { \leavevmode \skip_horizontal:n { \tex_fontdimen:D 2 \tex_font:D } }
\cs_new_protected:Npn \docxlike_format_verb_space:
  {
    \bool_if:NTF \g_docxlike_verb_space_visible_bool
      { \docxlike_format_verb_space_visible: } { \docxlike_format_verb_space_cell: }
  }
\cs_new_protected:Npn \docxlike_format_verb_enter:
  {
    \bool_set_true:N \l__docxlike_format_in_verb_bool
    \sys_if_engine_luatex:T { \docxlike@code@attr = 1 \scan_stop: }
    \mode_if_horizontal:T { \cs_set_eq:NN \@xobeysp \docxlike_format_verb_space: }
  }
%    \end{macrocode}
%
%    \begin{macrocode}
\cs_new_protected:Npn \docxlike_format_punct_grid:
  {
    \bool_lazy_all:nTF
      {
        { \cs_if_exist_p:N \xeCJK_pre_inter_class_toks:nnn }
        { \cs_if_exist_p:N \xeCJK_class_num:n }
        { \cs_if_exist_p:N \xeCJK_new_class:n }
        { \cs_if_exist_p:N \__xeCJK_save_CJK_class:n }
        { \seq_if_exist_p:N \g__xeCJK_class_seq }
        { \cs_if_exist_p:N \xeCJK_copy_inter_class_toks:nnnn }
        { \cs_if_exist_p:N \xeCJK_replace_inter_class_toks:nnnn }
      }
      {
        \xeCJKsetup { PunctStyle = plain , CheckFullRight = true }
%    \end{macrocode}
% 这枚覆盖原本无条件掏 \cs{CJKecglue}。老版 xeCJK 在“汉字、胶、西文”
% 这样隔着东西的边界照样调它，源码里“在~\cs{verb}”或换行空格接命令的
% 地方就叠出双份间距，2026 只在字符还贴着时才发。补一道闸：上一节点
% 已经是胶（\cs{lastnodetype} 为 11）就不再发。
%    \begin{macrocode}
        \cs_if_exist:NT \xeCJK_space_or_xecglue:
          {
            \cs_set_protected:Npn \xeCJK_space_or_xecglue:
              {
                \int_compare:nNnF { \tex_lastnodetype:D } = { 11 }
                  { \CJKecglue }
              }
          }
        \exhyphenpenalty = -5 \scan_stop:
        \docxlike_format_hyphen:
        \xeCJKDeclareCharClass { CJK } { "2026 }
        \seq_set_eq:NN \l_tmpa_seq \g__xeCJK_class_seq
        \xeCJK_new_class:n { DocxlikeQuote }
        \__xeCJK_save_CJK_class:n { DocxlikeQuote }
        \xeCJK_new_class:n { DocxlikeTilde }
        \__xeCJK_save_CJK_class:n { DocxlikeTilde }
        \seq_map_inline:Nn \l_tmpa_seq
          {
            \xeCJK_copy_inter_class_toks:nnnn { ##1 } { DocxlikeQuote }
              { ##1 } { FullLeft }
            \xeCJK_copy_inter_class_toks:nnnn { DocxlikeQuote } { ##1 }
              { FullLeft } { ##1 }
            \xeCJK_copy_inter_class_toks:nnnn { ##1 } { DocxlikeTilde }
              { ##1 } { FullRight }
            \xeCJK_copy_inter_class_toks:nnnn { DocxlikeTilde } { ##1 }
              { FullRight } { ##1 }
          }
        \xeCJK_copy_inter_class_toks:nnnn { DocxlikeQuote } { DocxlikeQuote }
          { FullLeft } { FullLeft }
        \xeCJK_copy_inter_class_toks:nnnn { DocxlikeTilde } { DocxlikeTilde }
          { FullRight } { FullRight }
        \xeCJK_copy_inter_class_toks:nnnn { DocxlikeQuote } { DocxlikeTilde }
          { FullLeft } { FullRight }
        \cs_if_exist:NF \__xeCJK_use_ecglue_skip:
          {
            \__docxlike_format_verb_hook_old:n { DocxlikeQuote }
            \__docxlike_format_verb_hook_old:n { DocxlikeTilde }
          }
        \xeCJK_copy_inter_class_toks:nnnn { DocxlikeTilde } { DocxlikeQuote }
          { FullRight } { FullLeft }
        \xeCJKDeclareCharClass { DocxlikeQuote } { "201C , "2018 }
        \xeCJKDeclareCharClass { DocxlikeTilde } { "FF5E , "2014 , "2015 }
%    \end{macrocode}
%
% \pkg{xeCJK} 在 2026 版把 \texttt{～}（U+FF5E）连同 \texttt{〜}、\texttt{゠} 从
% \option{MiddlePunct} 改成 \option{LongPunct}（\href{https://github.com/CTeX-org/ctex-kit/pull/793}{ctex-kit\#793}），
% 理由是与破折号 \texttt{—} 对齐。改动带来两件事：两侧不再生成压缩节点
% （胶加负宽 rule），这一条是对的，间距干净了；断行从禁止改成允许，这一条不对
% ——《标点符号用法》（GB/T~15834）里连接号与破折号都不得处于行首，正确的
% “与破折号对齐”应当是两个都禁止，不是两个都允许。
%
% 所以显式钉成 \option{NoBreakLongPunct}：长标点的间距，加上禁止断行。这个键
% 是 2026 新加的，旧版的 \cs{g\_\_xeCJK\_special\_punct\_clist} 里没有
% \texttt{nobreak\_long} 这一档，那边 \texttt{～} 留在 \option{MiddlePunct}，
% 断行同样禁止，只是多一层压缩节点。两头都不会出现行首波浪号。
%
% 传的是字符本身不是码点：\pkg{xeCJK} 那边 \cs{\_\_xeCJK\_add\_special\_punct:nn}
% 把参数按记号逐个取，上游自己也是先用 \cs{char\_generate:nn} 把码点转成字符
% 再传。
%    \begin{macrocode}
        \bool_lazy_and:nnT
          { \cs_if_exist_p:N \__xeCJK_add_special_punct:nn }
          { \clist_if_exist_p:N \g__xeCJK_special_punct_clist }
          {
            \clist_if_in:NnT \g__xeCJK_special_punct_clist { nobreak_long }
              {
                \use:x
                  {
                    \exp_not:N \__xeCJK_add_special_punct:nn { nobreak_long }
                      {
                        \char_generate:nn { "FF5E } { 12 }
                        \char_generate:nn { "301C } { 12 }
                        \char_generate:nn { "30A0 } { 12 }
                      }
                  }
              }
          }
        \xeCJK_pre_inter_class_toks:nnn { CJK } { FullRight }
          { \nobreak \CJKglue }
        \xeCJK_pre_inter_class_toks:nnn { FullLeft } { CJK }
          { \nobreak \CJKglue }
        \xeCJK_pre_inter_class_toks:nnn { FullLeft } { Default }
          { \docxlike_format_punct_ngap: }
        \xeCJK_pre_inter_class_toks:nnn { FullLeft } { FullRight }
          { \docxlike_format_punct_ngap: }
        \xeCJK_pre_inter_class_toks:nnn { FullLeft } { FullLeft }
          { \docxlike_format_punct_cpair: }
        \xeCJK_pre_inter_class_toks:nnn { CJK } { FullLeft }
          { \docxlike_format_punct_oshrink: }
        \xeCJK_pre_inter_class_toks:nnn { FullRight } { CJK }
          {
            \mode_if_inner:F
              {
                \tex_kern:D \dim_eval:n { -0.5 \g_docxlike_page_char_dim } \scan_stop:
                \discretionary { } { }
                  { \kern \dim_eval:n { 0.5 \g_docxlike_page_char_dim } }
                \nobreak \docxlike_format_punct_cshrink:
              }
          }
        \xeCJK_pre_inter_class_toks:nnn { FullRight } { FullRight }
          { \docxlike_format_punct_cpair: }
        \xeCJK_pre_inter_class_toks:nnn { FullRight } { FullLeft }
          { \docxlike_format_punct_cbreak: }
        \xeCJK_pre_inter_class_toks:nnn { FullRight } { Default }
          { \tex_kern:D 0.25 bp \scan_stop: }
        \xeCJK_pre_inter_class_toks:nnn { CJK } { DocxlikeQuote }
          { \docxlike_format_punct_oshrink: }
        \xeCJK_pre_inter_class_toks:nnn { FullRight } { DocxlikeQuote }
          { \docxlike_format_punct_bpair: }
        \xeCJK_pre_inter_class_toks:nnn { FullLeft } { DocxlikeQuote }
          { \docxlike_format_punct_cpair: }
        \xeCJK_pre_inter_class_toks:nnn { DocxlikeQuote } { DocxlikeQuote }
          { \docxlike_format_punct_cpair: }
        \xeCJK_pre_inter_class_toks:nnn { Default } { DocxlikeQuote }
          { }
        \xeCJK_pre_inter_class_toks:nnn { DocxlikeQuote } { CJK }
          { \nobreak \CJKglue }
        \xeCJK_pre_inter_class_toks:nnn { DocxlikeQuote } { FullLeft }
          { \docxlike_format_punct_cpair: }
        \xeCJK_pre_inter_class_toks:nnn { DocxlikeQuote } { FullRight }
          { \docxlike_format_punct_ngap: }
        \xeCJK_pre_inter_class_toks:nnn { DocxlikeQuote } { Default }
          { \docxlike_format_punct_ngap: }
        \xeCJK_pre_inter_class_toks:nnn { CJK } { DocxlikeTilde }
          { \nobreak \CJKglue }
%    \end{macrocode}
% 汉字接波浪线这一对，从 FullRight 拷来的默认体在 2026 里不出胶，
% 2025 及更早却多出 0.456\,bp（“3 万～5 万字”跨版本翻断点就是它）。
% 老版把这对整体重设，只留我们自己的禁则加行间胶，与 2026 逐位同宽。
%    \begin{macrocode}
        \cs_if_exist:NF \__xeCJK_use_ecglue_skip:
          {
            \xeCJK_inter_class_toks:nnn { CJK } { DocxlikeTilde }
              { \nobreak \CJKglue }
          }
        \xeCJK_pre_inter_class_toks:nnn { DocxlikeTilde } { Default }
          { \docxlike_format_punct_ngap: }
%    \end{macrocode}
%
% \emph{西文接全角标点那一格。} \hologo{XeTeX} 的 \option{LetterSpace} 只落在词内，
% 一段西文的末字没有尾距（\cs{hbox}\marg{12} 实测 12.227 而不是 12.454），后面接
% 汉字时那半份由中西文胶顶上，接全角标点时就整个丢了：改之前 \texttt{1）} 里数字
% 的步进只有 6.000。
%
% 补多少按样本文档反解。取样本文档与测试文档两侧的段末行（不参与齐行拉伸）逐字比对，
% 数字步进在数字前是 6.257，在全角标点前是 6.445，差 0.188；6.445 与
% “字号 12 加一整格增量 0.4543”只差 0.01，而 6.257 与“加半格”只差 0.03。
% 所以 Word 在这个边界上给的是整格，不是半格。声明进 \texttt{Default} 到
% \texttt{FullLeft}／\texttt{FullRight} 两对（全角的开号、闭号、冒号、方括号都
% 走这两类），17 个样本的均差由 $-0.442$ 收到 $+0.012$，与汉字对本身的 $+0.008$
% 同量级。
%    \begin{macrocode}
        \xeCJK_pre_inter_class_toks:nnn { Default } { DocxlikeTilde }
          { \docxlike_format_punct_xtra: }
        \xeCJK_pre_inter_class_toks:nnn { Default } { FullLeft }
          { \docxlike_format_punct_ngap: }
        \xeCJK_pre_inter_class_toks:nnn { Default } { FullRight }
          { \docxlike_format_punct_ngap: }
        \xeCJK_pre_inter_class_toks:nnn { FullRight } { DocxlikeTilde }
          { \docxlike_format_punct_ngap: }
        \xeCJK_pre_inter_class_toks:nnn { DocxlikeTilde } { FullRight }
          { \docxlike_format_punct_ngap: }
        \xeCJK_pre_inter_class_toks:nnn { DocxlikeTilde } { DocxlikeTilde }
          {
            \mode_if_inner:TF
              { \tex_kern:D \dim_eval:n { \g_docxlike_page_char_dim - 1em } \scan_stop: }
              { \CJKglue }
          }
        \xeCJK_pre_inter_class_toks:nnn { DocxlikeQuote } { DocxlikeTilde }
          { \docxlike_format_punct_ngap: }
        \xeCJK_pre_inter_class_toks:nnn { DocxlikeTilde } { DocxlikeQuote }
          { \docxlike_format_punct_ngap: }
        \xeCJKDeclareCharClass { HalfRight }
          { "21 , "29 , "2C , "2E , "3A , "3B , "3F , "5D , "D7 , "F7 }
        \xeCJKDeclareCharClass { HalfLeft } { "28 , "5B }
        \xeCJK_pre_inter_class_toks:nnn { CJK } { HalfRight }
          { \docxlike_format_punct_ngap: }
        \xeCJK_pre_inter_class_toks:nnn { FullLeft } { HalfRight }
          { \docxlike_format_punct_ngap: }
        \xeCJK_pre_inter_class_toks:nnn { DocxlikeQuote } { HalfRight }
          { \docxlike_format_punct_ngap: }
%    \end{macrocode}
%
% 摘胶别把断点摘没了。\hologo{TeX} 只在胶处断行，把“半角闭号→汉字”的
% \cs{CJKecglue} 替换成空之后，“(2022)发表”这种地方就断不开——recruit
% 章那段第一行只好退到前一个词间空格断，剩 268\,bp 硬拉满 425\,bp，每个
% 汉字之间被扯出一个字宽。补一道零宽的胶：宽度照旧是零（Word 实测裸宽），
% 断点回来（Word 正是断在这儿，闭号留行尾禁则也允许）。
% “半角开号→汉字”那对保持替换成空：开号不做行尾是禁则的真要求，
% 那儿本来就该断不开。
%    \begin{macrocode}
        \xeCJK_replace_inter_class_toks:nnnn { HalfRight } { CJK }
          { \CJKecglue } { \skip_horizontal:N \c_zero_skip }
        \xeCJK_replace_inter_class_toks:nnnn { HalfLeft } { CJK }
          { \CJKecglue } { }
        \xeCJK_replace_inter_class_toks:nnnn { CJK } { HalfLeft }
          { \CJKecglue } { \CJKglue }
%    \end{macrocode}
%
% 全角空格 U+3000 在 Word 里是空格，旁边不加中西文间距，与 \hologo{LuaLaTeX} 那边的
% \texttt{jaxspmode} 同一件事（见 \cs{docxlike\_format\_hglue:nn}）。\pkg{xeCJK} 把它归在
% \texttt{CJK} 类里，这里另开一类，各个边界照 \texttt{CJK} 抄（放在最后，上面对
% \texttt{CJK} 的改动一并带上），再把它与 \texttt{Default}（字母数字）之间的
% \cs{CJKecglue} 摘掉。
%    \begin{macrocode}
        \seq_set_eq:NN \l_tmpa_seq \g__xeCJK_class_seq
        \xeCJK_new_class:n { DocxlikeIdeoSpace }
        \__xeCJK_save_CJK_class:n { DocxlikeIdeoSpace }
        \seq_map_inline:Nn \l_tmpa_seq
          {
            \xeCJK_copy_inter_class_toks:nnnn { ##1 } { DocxlikeIdeoSpace }
              { ##1 } { CJK }
            \xeCJK_copy_inter_class_toks:nnnn { DocxlikeIdeoSpace } { ##1 }
              { CJK } { ##1 }
          }
        \xeCJK_copy_inter_class_toks:nnnn { DocxlikeIdeoSpace } { DocxlikeIdeoSpace }
          { CJK } { CJK }
        \xeCJK_replace_inter_class_toks:nnnn { Default } { DocxlikeIdeoSpace }
          { \CJKecglue } { }
        \xeCJK_replace_inter_class_toks:nnnn { DocxlikeIdeoSpace } { Default }
          { \CJKecglue } { }
        \xeCJKDeclareCharClass { DocxlikeIdeoSpace } { "3000 }
      }
      {
        \sys_if_engine_luatex:TF
          { \msg_warning:nn { docxlike } { luatex-punct-grid } }
          { \msg_error:nn { docxlike } { xecjk-punct-grid } }
      }
  }
%    \end{macrocode}
%
% \emph{正文的对齐与段间距只在写了的时候发。} \option{alignment} 查的是标题用的
% 那张表，\texttt{left} 发 \cs{raggedright}，两端对齐与分散对齐不发（\hologo{TeX}
% 本来就两端对齐）。段前段后折成 \cs{parskip}：正文段接正文段，两段之间是前一段的
% 段后与后一段的段前取大，每一对都是同一个数，页首页尾那一截 \hologo{TeX} 在分页处
% 丢掉，与 Word 丢段前段后的地方一致。两个都没写就不碰 \cs{parskip}，文档类
% 原先给的值（可能带伸缩量）照旧。
%
% \emph{标题两侧不吃 \cs{parskip}。} \cs{parskip} 加在每一段的\emph{前面}，
% Word 的段后却属于前一段：正文接标题，Word 是正文段后与标题段前取大；标题接
% 正文，Word 只有标题的段后（正文段前是零），\hologo{TeX} 却还要再加一个
% \cs{parskip}。所以两处各退一步：标题那一段与标题后的头一段，在
% \texttt{para/before} 钩子里发一段 $-$\cs{parskip} 抵掉内核刚加的那一段
% （标题后的头一段认 \cs{if@nobreak}，那是 \cs{@afterheading} 立的）；标题那一段本身由
% \cs{docxlike\_format\_par:nn} 排，它在组里把 \cs{parskip} 清零，段前与前一段的段后取大。
%
% \emph{段与段之间按行盒接。} Word 里下一段首行的基线落在“上一段末行的行深、
% 两段之间的间距、本段首行的上升部”三者之和处。撑杆已经把每一行做成 Word 的
% 行盒，只差 \hologo{TeX} 在两段之间补的那段行间胶：它按\emph{本段}的
% \cs{baselineskip} 算，本段行距比上一段大就多出一截（空白文档里正文接 24\,bp
% 标题多 10.6\,bp），小就退成 \cs{lineskip}（标题接正文多 1\,bp）。
% 所以段落起头时（\texttt{para/before}，竖直模式，内核刚加完 \cs{parskip}）把
% \cs{prevdepth} 改成“本段行距减本段首行的上升部”，首行那段行间胶就恰好是零；
% 行里有比撑杆高的东西时胶是负的，退成 \cs{lineskip}，行盒跟着内容长高，与 Word
% 一样。\cs{prevdepth} 被别处压成 $-1000$\,pt（不要行间胶）时不动。这一条改的是
% 所有段落，挂哪些文档类要各自核对，所以是开关：
% \cs{docxlike\_format\_word\_interline:} 打开，要在 \cs{docxlike\_format\_apply\_body:}
% 之前调。空白文档预设打开它。
%
%    \begin{macrocode}
\cs_generate_variant:Nn \prop_get:NnNT { NVNT }
\bool_new:N \l__docxlike_format_spaced_bool
\dim_new:N \g__docxlike_format_parskip_dim
\bool_new:N \g__docxlike_format_interline_bool
\bool_new:N \g__docxlike_format_para_hook_bool
\cs_new_protected:Npn \docxlike_format_word_interline:
  { \bool_gset_true:N \g__docxlike_format_interline_bool }
\cs_new_protected:Npn \__docxlike_format_body_parskip:
  {
    \bool_set_false:N \l__docxlike_format_spaced_bool
    \clist_map_inline:nn { before , after }
      {
        \__docxlike_format_slot_if_set:nnT { Normal } {##1}
          { \bool_set_true:N \l__docxlike_format_spaced_bool }
      }
    \bool_if:NT \l__docxlike_format_spaced_bool
      {
        \dim_gset:Nn \g__docxlike_format_parskip_dim
          { \dim_max:nn { \l_docxlike_format_before_dim } { \l_docxlike_format_after_dim } }
        \skip_set:Nn \parskip { \g__docxlike_format_parskip_dim }
      }
  }
\cs_new_protected:Npn \__docxlike_format_para_hook:
  {
    \bool_lazy_and:nnT
      { ! \g__docxlike_format_para_hook_bool }
      {
        \g__docxlike_format_interline_bool
        || \dim_compare_p:nNn { \g__docxlike_format_parskip_dim } > { 0pt }
      }
      {
        \bool_gset_true:N \g__docxlike_format_para_hook_bool
        \AddToHook { para / before } { \__docxlike_format_para_before: }
      }
  }
\cs_new_protected:Npn \__docxlike_format_para_before:
  {
    \legacy_if:nT { @nobreak } { \skip_vertical:n { - \parskip } }
    \bool_if:NT \g__docxlike_format_interline_bool
      {
        \bool_if:NTF \l_docxlike_format_strut_latin_bool
          {
            \__docxlike_format_strut_dims:nn
              { \l_docxlike_format_natural_latin_tl }
              { \l_docxlike_format_ascent_latin_tl }
          }
          {
            \__docxlike_format_strut_dims:nn
              { \l__docxlike_format_strut_natural_tl } { \l_docxlike_format_ascent_tl }
          }
        \__docxlike_format_prevdepth:n
          { \baselineskip - \l__docxlike_format_asc_dim }
      }
  }
\cs_generate_variant:Nn \docxlike_format_compute:n { V }
\cs_new_protected:Npn \__docxlike_format_prevdepth:n #1
  {
    \dim_compare:nNnT { \tex_prevdepth:D } > { -1000pt }
      { \tex_prevdepth:D = \dim_eval:n {#1} \scan_stop: }
  }
%    \end{macrocode}
%
%    \begin{macrocode}
\cs_new_protected:Npn \docxlike_format_apply_body:
  {
    \__docxlike_format_if_blank:nF { Normal }
      {
        \docxlike_format_compute:n { Normal }
        \bool_set_eq:NN \l_docxlike_format_strut_latin_bool \l__docxlike_format_latin_bool
        \dim_gset_eq:NN \g_docxlike_body_size_dim \l_docxlike_format_size_dim
        \dim_gset_eq:NN \g_docxlike_body_lead_dim \l_docxlike_format_lead_dim
        \dim_gset_eq:NN \g_docxlike_body_stretch_dim \l_docxlike_format_stretch_dim
        \dim_gset_eq:NN \g_docxlike_body_first_dim \l_docxlike_format_first_dim
        \dim_gset_eq:NN \g_docxlike_body_indent_dim \l_docxlike_format_indent_dim
        \dim_gset_eq:NN \g__docxlike_format_body_after_dim \l_docxlike_format_after_dim
        \bool_set_eq:NN \l_docxlike_format_line_snap_bool \l_docxlike_format_snap_vertical_bool
        \dim_set_eq:NN \topskip \g_docxlike_body_first_dim
        \docxlike_format_reserve_depth:
        \docxlike_format_penalties:
        \normalsize
        \__docxlike_format_get:nnNT { Normal } { pPr/jc } \l__docxlike_format_value_tl
          {
            \prop_get:NVNT \g__docxlike_format_align_prop \l__docxlike_format_value_tl
              \l__docxlike_format_value_tl { \l__docxlike_format_value_tl }
          }
        \__docxlike_format_body_parskip:
        \__docxlike_format_para_hook:
        \dim_set_eq:NN \parindent \g_docxlike_body_indent_dim
        \docxlike_format_latin_pitch:
%    \end{macrocode}
%
% \emph{字距最后设。}
%    \begin{macrocode}
        \dim_compare:nNnF { \g_docxlike_page_char_dim } = { 0pt }
          {
            \docxlike_format_hglue:nn
              { \__docxlike_format_hglue_cc: } { \__docxlike_format_hglue_ec: }
            \bool_if:NT \l_docxlike_format_snap_horizontal_bool
              {
                \bool_if:NTF \l_docxlike_format_snap_vertical_bool
                  {
                    \docxlike_format_slack:
                    \skip_set_eq:NN \rightskip \g_docxlike_page_slack_dim
                  }
                  { \skip_set:Nn \rightskip { \__docxlike_format_grid_glue: } }
              }
            \docxlike_format_punct_grid:
          }
        \AddToHook { para / begin } { \__docxlike_format_para_open: }
        \AddToHook { para / end }   { \__docxlike_format_para_close: }
        \hook_use:n { docxlike / display }
        \docxlike_format_display_hooks:
        \cs_if_exist:NT \docxlike_header_strut:
          { \cs_set_eq:NN \docxlike_header_strut: \docxlike_format_strut: }
      }
  }
%    \end{macrocode}
%
% \subsection{断字与空行}
%
% \emph{断字全文一个开关} \option{hyphenation}，正文、题注、
% 纯西文段一律跟随，一份文档不设两套规矩。
%
% 默认\emph{关}。依据三条：样本文档的 Word 导出勾了不断字；样本文档逐字量下来 Word 没有
% 断 \texttt{Nature Communications}；中文论文里的西文通常不断字。行末基线对这个
% 开关几乎无感——\cs{hyphenpenalty} 取 50/85/200/1000 都是行差 79，取 10000 是
% 80，所以默认值按排版道理定，不迁就基线。
%
% \cs{exhyphenpenalty}（在\emph{已有}的连字符处断行，\texttt{Tree-structured}
% 这种）不归这个开关管：GB/T~15834 允许，样本文档里 Word 实测也断在 \texttt{Tree-}，
% 而且这个值是 $-5$，是行末闭号那套梯子调出来的。拉丁段落也不另设（实测行差不变）。
%    \begin{macrocode}
\cs_new_protected:Npn \docxlike_format_hyphen:
  {
    \int_set:Nn \hyphenpenalty
      { \bool_if:NTF \g_docxlike_format_hyphen_bool { 85 } { 10000 } }
  }
%    \end{macrocode}
%
% \begin{macro}{\docxlike_format_apply_par:n}
% \cs{docxlike\_format\_apply\_par:n}\marg{目标} 按目标给\emph{接下来那几段}铺好版面：
% 字体、字号、行距、首行基线落点、行末余量、断词。与 \cs{docxlike\_format\_par:nn}
% 的分工是，那个排\emph{一段}，这个铺\emph{一片}。
%
% 撑杆在 \cs{l\_\_docxlike\_format\_out\_tl} 末尾，它会 \cs{leavevmode} 起段，铺版面
% 时不能让它跑。执行前摘掉、执行完挂回去——每段的撑杆由 \texttt{para/begin} 钩子
% 发，那一根仍在。
%    \begin{macrocode}
\cs_new_eq:NN \__docxlike_format_strut_saved: \prg_do_nothing:
\cs_new_protected:Npn \docxlike_format_apply_par:n #1
  {
    \__docxlike_format_if_blank:nF {#1}
      {
        \docxlike_format_apply_font:nN {#1} \l__docxlike_format_par_tl
        \cs_set_eq:NN \__docxlike_format_strut_saved: \docxlike_format_strut:
        \cs_set_eq:NN \docxlike_format_strut: \prg_do_nothing:
        \l__docxlike_format_par_tl
        \cs_set_eq:NN \docxlike_format_strut: \__docxlike_format_strut_saved:
        \dim_set_eq:NN \topskip \l_docxlike_format_first_dim
        \bool_set_eq:NN \l_docxlike_format_line_snap_bool
          \l_docxlike_format_snap_vertical_bool
        \bool_if:NF \l_docxlike_format_snap_horizontal_bool
          { \skip_zero:N \rightskip }
        \docxlike_format_hyphen:
      }
  }
%    \end{macrocode}
%    \end{macro}
%
% \begin{macro}{\parmark}
% 模拟 Word 里单独按一次回车留下的空段。空段行高 = 段落标记字体的自然行高
% $\times$ 行距，Word 的段落标记是西文字符，默认按 Times 算；倍数乘在自然行高上，
% 固定值就是那么高，最小值取两者较大者。
%
% \emph{行距与贴格跟正文走。} 回车敲出的空段继承上一段的段落属性，绝大多数空段
% 敲在正文里，所以默认按 \texttt{Normal} 目标现算一遍（\cs{docxlike\_format\_compute:n}，
% 算完 \cs{l\_docxlike\_format\_line\_unit\_tl}、\cs{l\_docxlike\_format\_line\_value\_tl} 与
% \cs{l\_docxlike\_format\_snap\_vertical\_bool} 就是正文那一档）。不能拿上一次算过的目标凑数：页眉页脚、表格都会算，留下的是谁
% 说不准（回归测试里就曾拿到单倍）。贴格的空段与正文的行走同一条换算
% （\cs{\_\_docxlike\_format\_endpoint:nnN}）：单倍凑成一整格，1.25 倍是 1.25 格。
% 这是 2026-09-21 用 Word 探针量的：无网格或段上写了
% \texttt{snapToGrid=0} 时 Times 单倍 13.8、1.25 倍 17.25，段落标记带
% \texttt{hint=eastAsia} 也不改变；贴格时单倍 19.55、1.25 倍 24.4。样本文档正文
% 是 1.25 倍不贴格，所以正文里的空段仍是 17.25。\texttt{Normal} 没写行距
% 时退回 1.25 倍不贴格。正文的行距单位是文档类登记的单位时，空段照那个单位的代码折
% （字号取段落标记的）；写了区间只取第一个端点，空段是刚性的。
%
% 恰逢换页时默认随页首丢弃（页首的空行排出来没有意义）。星号版与 Word 的
% 空段同构：每行一只高为空段行高的空盒，盒不可弃、盒间是合法断点，页底装得下
% 几行装几行，剩下的落到新页页首照常占位；\cs{vspace*} 只能整体保留，劈不开。
% 盒前 \cs{nointerlineskip}、盒后把 \cs{prevdepth} 还原，前后段的行间胶不受影响。
% 断点得明写：\cs{nointerlineskip} 把盒间的行间胶也去掉了，两只盒子之间没有胶就
% 没有合法断点，十三个空段会整块搬到下一页（逐页对拍时踩过：Word 在页底装了十个，
% 我们一个都没装）。每只盒前发一个零罚分，上一段与第一只盒之间也可断。
% 星号跟 \cs{vspace*} 的惯例，星号才保留。抄 Word 稿子时写星号版。
%
% \emph{紧跟在浮动体后面时不动浮动体的后距。} 文档类给图表前后留的距离（比如各
% 一个网格行）是它自己的规则，空段是作者另加的，两者相加。Word 里与之对应的是
% 表后头一段（哪怕是空段）段前一行。曾经试过让空段把后距减掉去凑样本文档（样本里
% 表后的空段没有段前），没有采用：浮动体的后距不能砍。
%
% 可选参数是回车几次；含等号时按键值读，写法与样式同：\option{lines} 几行；
% \option{font} 里 \option{size}、\option{asian-text-font}、\option{latin-text-font} 指定段落标记的字号
% 与字体，自然行高只认实测值，没量过的字体报错不瞎给；\option{paragraph} 里 \option{spacing} 的
% \option{line-spacing} 与 \option{at}，以及 \option{snap-to-grid-when-document-grid-is-defined} 贴不贴网格（网格没启用时写
% \texttt{true} 也不贴）。
%    \begin{macrocode}
\int_new:N \l__docxlike_parmark_lines_int
\bool_new:N \l__docxlike_parmark_snap_bool
\dim_new:N \l__docxlike_parmark_size_dim
\dim_new:N \l__docxlike_parmark_base_dim
\dim_new:N \l__docxlike_parmark_natural_dim
\dim_new:N \l__docxlike_parmark_ht_dim
\fp_new:N \l__docxlike_parmark_ratio_fp
\tl_new:N \l__docxlike_parmark_mode_tl
\tl_new:N \l__docxlike_parmark_val_tl
\tl_new:N \l__docxlike_parmark_last_tl
%    \end{macrocode}
%
% \emph{自然行高按 Word 的输出量，不按字体的 hhea 行高算}（Times 的 hhea 给
% 1.1472，12\,bp 上差 0.034）。量法：造一份 docx，八个字号（9、10.5、12、14、16、18、22、24）各排六段
% Times 单倍行距的文字，用 Word 转 PDF 再量相邻基线的步进。
%
% \begin{longtblr}[caption={Times 自然行高的标定}, label={tab:impl-natural-latin}]
%   {colspec={l r r r r r r r r}}
% \toprule
% 字号   & 9 & 10.5 & 12 & 14 & 16 & 18 & 22 & 24 \\ \midrule
% 行距   & 10.360 & 12.120 & 13.800 & 16.100 & 18.400 & 20.720 & 25.250 & 27.620 \\
% \bottomrule
% \end{longtblr}
%
% 过原点最小二乘是 1.15014，而且 12\,bp 正好 13.800、16\,bp 正好 18.400。
% 汉字那一档同法量得 1.29726，与 \cs{c\_docxlike\_format\_natural\_tl} 的 1.296875
% 一致。两个数最先见于 Typst 包 \pkg{banshi} 的标定，这里独立复量过。
%    \begin{macrocode}
\str_const:Nn \c__docxlike_parmark_alpha_str { abcdefghijklmnopqrstuvwxyz }
\prg_generate_conditional_variant:Nnn \str_if_in:Nn { NV } { TF }
\cs_new_protected:Npn \__docxlike_parmark_size:n #1
  {
    \prop_get:NnNTF \g__docxlike_size_prop {#1} \l__docxlike_size_tl
      { \dim_set:Nn \l__docxlike_parmark_size_dim { \l__docxlike_size_tl } }
      { \dim_set:Nn \l__docxlike_parmark_size_dim { \__docxlike_bp:n {#1} } }
  }
\dim_new:N \l__docxlike_parmark_asc_dim
\dim_new:N \l__docxlike_parmark_prev_dim
\tl_new:N \l__docxlike_parmark_ascent_tl
\tl_set_eq:NN \l__docxlike_parmark_ascent_tl \c_docxlike_format_ascent_latin_tl
\cs_new_protected:Npn \__docxlike_parmark_font:n #1
  {
    \docxlike_font_metric:nNNT {#1} \l__docxlike_format_coef_tl \l__docxlike_parmark_ascent_tl
      { \fp_set:Nn \l__docxlike_parmark_ratio_fp { \l__docxlike_format_coef_tl } }
  }
\cs_new_protected:Npn \__docxlike_parmark_spacing:n #1
  {
    \str_case:nnF {#1}
      {
        { single } { \tl_set:Nn \l__docxlike_parmark_mode_tl { auto }
                     \tl_set:Nn \l__docxlike_parmark_val_tl { 1 } }
        { double } { \tl_set:Nn \l__docxlike_parmark_mode_tl { auto }
                     \tl_set:Nn \l__docxlike_parmark_val_tl { 2 } }
      }
      {
        \tl_set:Nx \l__docxlike_parmark_last_tl { \str_item:nn {#1} { -1 } }
        \str_case:VnF \l__docxlike_parmark_last_tl
          {
            { ! }
              {
                \tl_set:Nn \l__docxlike_parmark_mode_tl { exact }
                \tl_set:Nx \l__docxlike_parmark_val_tl { \str_range:nnn {#1} { 1 } { -2 } }
              }
            { + }
              {
                \tl_set:Nn \l__docxlike_parmark_mode_tl { atLeast }
                \tl_set:Nx \l__docxlike_parmark_val_tl { \str_range:nnn {#1} { 1 } { -2 } }
              }
          }
          {
            \str_if_in:NVTF \c__docxlike_parmark_alpha_str \l__docxlike_parmark_last_tl
              { \tl_set:Nn \l__docxlike_parmark_mode_tl { exact } }
              { \tl_set:Nn \l__docxlike_parmark_mode_tl { auto } }
            \tl_set:Nn \l__docxlike_parmark_val_tl {#1}
          }
      }
  }
\keys_define:nn { docxlike / parmark }
  {
    lines .int_set:N = \l__docxlike_parmark_lines_int ,
    font .code:n = { \__docxlike_keys_set:nn { docxlike / parmark / font } {#1} } ,
    paragraph .code:n =
      { \__docxlike_keys_set:nn { docxlike / parmark / paragraph } {#1} } ,
    unknown .code:n =
      { \__docxlike_unknown_key: } ,
  }
\keys_define:nn { docxlike / parmark / font }
  {
    size .code:n = { \__docxlike_parmark_size:n {#1} } ,
    latin-text-font .code:n = { \__docxlike_parmark_font:n {#1} } ,
    asian-text-font .code:n = { \__docxlike_parmark_font:n {#1} } ,
    unknown .code:n =
      { \__docxlike_unknown_key: } ,
  }
\keys_define:nn { docxlike / parmark / paragraph }
  {
    spacing .code:n = { \__docxlike_parmark_rule:n {#1} } ,
    snap-to-grid-when-document-grid-is-defined .bool_set:N = \l__docxlike_parmark_snap_bool ,
    snap-to-grid-when-document-grid-is-defined .default:n = { true } ,
    unknown .code:n =
      { \__docxlike_unknown_key: } ,
  }
\keys_define:nn { docxlike / parmark / spacing }
  {
    line-spacing .tl_set:N = \l__docxlike_style_rule_tl ,
    at .tl_set:N = \l__docxlike_style_at_tl ,
    unknown .code:n =
      { \__docxlike_unknown_key: } ,
  }
\cs_new_protected:Npn \__docxlike_parmark_rule:n #1
  {
    \tl_clear:N \l__docxlike_style_rule_tl
    \tl_clear:N \l__docxlike_style_at_tl
    \keys_set:nn { docxlike / parmark / spacing } {#1}
    \tl_if_empty:NF \l__docxlike_style_rule_tl
      {
        \__docxlike_parmark_spacing:e
          {
            \str_case:Vn \l__docxlike_style_rule_tl
              {
                { single } { single }
                { 1.5-lines } { 1.5 }
                { double } { double }
                { multiple } { \l__docxlike_style_at_tl }
                { exactly } { \__docxlike_bp:V \l__docxlike_style_at_tl ! }
                { at-least } { \__docxlike_bp:V \l__docxlike_style_at_tl + }
              }
          }
      }
  }
\cs_generate_variant:Nn \__docxlike_parmark_spacing:n { e }
\cs_new_protected:Npn \__docxlike_parmark_height:
  {
    \dim_compare:nNnTF \l__docxlike_parmark_size_dim > { 0pt }
      { \dim_set_eq:NN \l__docxlike_parmark_base_dim \l__docxlike_parmark_size_dim }
      { \dim_set:Nn \l__docxlike_parmark_base_dim { 1 em } }
    \dim_set:Nn \l__docxlike_parmark_natural_dim
      { \fp_to_decimal:N \l__docxlike_parmark_ratio_fp \l__docxlike_parmark_base_dim }
    \dim_set_eq:NN \l_docxlike_format_natural_dim \l__docxlike_parmark_natural_dim
    \bool_set:Nn \l_docxlike_format_snap_vertical_bool
      {
        \bool_if_p:N \l__docxlike_parmark_snap_bool
        && \bool_if_p:N \g_docxlike_grid_bool
        && ! \dim_compare_p:nNn { \g_docxlike_page_lead_dim } = { 0pt }
      }
    \str_case:VnF \l__docxlike_parmark_mode_tl
      {
        { auto }
          {
            \__docxlike_format_endpoint:nVN { auto } \l__docxlike_parmark_val_tl
              \l__docxlike_parmark_ht_dim
          }
        { exact }
          { \dim_set:Nn \l__docxlike_parmark_ht_dim { \l__docxlike_parmark_val_tl } }
        { atLeast }
          {
            \__docxlike_format_endpoint:nVN { atLeast } \l__docxlike_parmark_val_tl
              \l__docxlike_parmark_ht_dim
          }
      }
      {
        \cs_if_exist:cTF { __docxlike_format_line_unit_ \l__docxlike_parmark_mode_tl :nN }
          {
            \dim_set_eq:NN \l_docxlike_format_size_dim \l__docxlike_parmark_base_dim
            \__docxlike_format_endpoint:VVN \l__docxlike_parmark_mode_tl
              \l__docxlike_parmark_val_tl \l__docxlike_parmark_ht_dim
          }
          { \dim_set:Nn \l__docxlike_parmark_ht_dim { \l__docxlike_parmark_val_tl } }
      }
    \dim_set:Nn \l__docxlike_parmark_asc_dim
      {
        \dim_min:nn
          {
            \l__docxlike_parmark_ascent_tl \l__docxlike_parmark_base_dim
          }
          { \l__docxlike_parmark_ht_dim }
      }
  }
\NewDocumentCommand \parmark { s O{1} }
  {
    \par
    \dim_gzero:N \g__docxlike_format_pending_after_dim
    \group_begin:
    \int_set:Nn \l__docxlike_parmark_lines_int { 1 }
    \dim_zero:N \l__docxlike_parmark_size_dim
    \fp_set:Nn \l__docxlike_parmark_ratio_fp { \c_docxlike_format_natural_latin_tl }
    \tl_set_eq:NN \l__docxlike_parmark_ascent_tl \c_docxlike_format_ascent_latin_tl
    \docxlike_format_compute:n { Normal }
    \__docxlike_format_slot_if_set:nnTF { Normal } { line }
      {
        \tl_set_eq:NN \l__docxlike_parmark_val_tl \l_docxlike_format_line_value_tl
        \tl_set_eq:NN \l__docxlike_parmark_mode_tl \l_docxlike_format_line_unit_tl
      }
      {
        \tl_set:Nn \l__docxlike_parmark_mode_tl { auto }
        \tl_set:Nn \l__docxlike_parmark_val_tl { 1.25 }
      }
    \bool_set_eq:NN \l__docxlike_parmark_snap_bool \l_docxlike_format_snap_vertical_bool
    \tl_if_in:nnTF {#2} { = }
      { \keys_set:nn { docxlike / parmark } {#2} }
      { \int_set:Nn \l__docxlike_parmark_lines_int {#2} }
    \__docxlike_parmark_height:
    \IfBooleanTF {#1}
      {
        \dim_set:Nn \l__docxlike_parmark_prev_dim { \tex_prevdepth:D }
        \prg_replicate:nn { \l__docxlike_parmark_lines_int }
          {
            \penalty \c_zero_int
            \nointerlineskip
            \hbox:n
              {
                \vrule ~ width ~ \c_zero_dim
                  ~ height ~ \l__docxlike_parmark_ht_dim
                  ~ depth ~ \c_zero_dim
              }
          }
        \dim_set:Nn \tex_prevdepth:D { \l__docxlike_parmark_prev_dim }
      }
      {
        \vspace
          { \dim_eval:n { \l__docxlike_parmark_ht_dim * \l__docxlike_parmark_lines_int } }
      }
    \group_end:
  }
%    \end{macrocode}
% \end{macro}
%
%    \begin{macrocode}
%</docxlike-format>
%    \end{macrocode}
%
% \Finale
